\part*{Part IV — Application}
\label{navigation:IV}
\addcontentsline{toc}{part}{Part IV — Application}
Begin with the distinction between managed active information and authoritative persistence. Mneme manages Context and Pragma and provides Context Load, Context Save, Adhoc Resource Access and Adhoc Search. Mnemosyne establishes authoritative durable state and persistence outcomes. Cache retention, completed execution and publication of outputs do not independently establish durable authority. The component definitions retain the candidate standing and detailed-contract reservations of the accepted baseline.

{\small Reading sources: \hyperref[element:1af8a56358f7429fabfb57be66542a99]{Mneme}; \hyperref[element:7985e647bf9e4ca3ba70faf9e3e4ae02]{Mnemosyne}.}\par
Execution, operational management and messaging are separate responsibilities. Ponos provides governed execution; Energeia owns Praxis and Journey and observes operational progression. Energeia's R1.x/R2.x boundary is Observe to Report, without autonomous changes to execution or workload. Petasos manages durable queues and PetasosParcel handoff; referenced information retains its own authority. Read these responsibilities together without turning their collaboration into an inferred component hierarchy.

{\small Reading sources: \hyperref[element:b521762c1ea444b48cd5f66be3752d30]{Ponos}; \hyperref[element:fb3f37a6ac614415ac486a6de9d35994]{Energeia}; \hyperref[element:195e5f10db8e419391354197b5d7a8cc]{Petasos}.}\par
At the interaction boundaries, Pylai supplies external interoperability while Iris supplies contextual human interaction. Themis owns authorisation semantics and consumes trusted authenticated Identity rather than authenticating it. Agora controls collaboration semantics and managed Room/Space intent; Omilia performs collaboration operations. These responsibilities do not confer ownership of information merely conveyed through them. Detailed representations and deferred mechanisms remain as qualified in their definitions.

{\small Reading sources: \hyperref[element:e3e4147fea3640198dc251ed41ee854d]{Pylai}; \hyperref[element:e927c1af96c64fcd979be659bee2cbe3]{Iris}; \hyperref[element:899d34ec491948b68aad09e507dea6d1]{Themis}; \hyperref[element:d1c4053e68384080b8eb79d2008632ba]{Agora}.}\par
Calliope-Models and Calliope-DataSets distinguish build-time representation from runtime content management. Models supplies structural definitions and library contracts. DataSets owns mutable governed content and uses the established information-management and persistence responsibilities. This distinction does not select a new deployment or storage allocation. The information map that follows remains a separate view of information concepts, not an inferred physical schema.

{\small Reading sources: \hyperref[element:307b3440a507403f8bc214d290bffb6a]{Calliope-Models}; \hyperref[element:53bf6815c7a54ed0abde9b1aad311843]{Calliope-DataSets}.}\par
\begin{center}\begin{tikzpicture}[>=Stealth,every node/.style={draw=Ink!40,rounded corners,align=left,font=\small,text width=5.2cm,inner sep=6pt}]
\node (source) at (0,0) {\hyperref[element:4af969b8c23a4dcab7569dc51b263383]{[Topic] Individual}};
\node (target0) at (8,4.25) {\hyperref[element:99dc695355384ea29ac3515b9178dfd0]{[Topic] Authority}};
\draw[->] (source.east) -- (3,0) |- (target0.west);
\node[draw=none,text width=1.8cm,font=\tiny,align=center,inner sep=0pt] at (4.15,4.63) {Composition\newline{}source-qualified / unclassified proposal};
\node (target1) at (8,2.55) {\hyperref[element:5c33b01e23964ef5935be79364e7c0f9]{[Topic] Contact}};
\draw[->] (source.east) -- (3,0) |- (target1.west);
\node[draw=none,text width=1.8cm,font=\tiny,align=center,inner sep=0pt] at (4.15,2.93) {Composition\newline{}source-qualified / unclassified proposal};
\node (target2) at (8,0.85) {\hyperref[element:95a13e957f5c471797869b901085266c]{[Topic] Demographics}};
\draw[->] (source.east) -- (3,0) |- (target2.west);
\node[draw=none,text width=1.8cm,font=\tiny,align=center,inner sep=0pt] at (4.15,1.23) {Composition\newline{}source-qualified / unclassified proposal};
\node (target3) at (8,-0.85) {\hyperref[element:3b72456b5686449f95ad26bfbbae07f3]{[Topic] Identity}};
\draw[->] (source.east) -- (3,0) |- (target3.west);
\node[draw=none,text width=1.8cm,font=\tiny,align=center,inner sep=0pt] at (4.15,-0.47) {Composition\newline{}source-qualified / unclassified proposal};
\node (target4) at (8,-2.55) {\hyperref[element:e92ea15616ea43c9a1d698e34ed1376b]{[Topic] Patient}};
\draw[->] (source.east) -- (3,0) |- (target4.west);
\node[draw=none,text width=1.8cm,font=\tiny,align=center,inner sep=0pt] at (4.15,-2.17) {Composition\newline{}source-qualified / unclassified proposal};
\node (target5) at (8,-4.25) {\hyperref[element:1163ffac87814f319578ef5cb95dcd4d]{[Topic] Person}};
\draw[->] (source.east) -- (3,0) |- (target5.west);
\node[draw=none,text width=1.8cm,font=\tiny,align=center,inner sep=0pt] at (4.15,-3.87) {Composition\newline{}source-qualified / unclassified proposal};
\end{tikzpicture}\end{center}
{\small Selected outgoing Composition assertions; arrows retain source-to-target direction. Labels retain assertion standing and record state; element standing accompanies each description. \href{structure.json}{Complete relationship qualifications and provenance} remain source-defined. No omitted edge or semantic effect is inferred.}\par
\chapter*{IV.1 Information Architecture}
\phantomsection
\label{navigation:IV.1}
\addcontentsline{toc}{chapter}{IV.1 Information Architecture}

\section*{IV.1.1 Topics}
\phantomsection
\label{navigation:IV.1.1}
\addcontentsline{toc}{section}{IV.1.1 Topics}

\subsection*{IV.1.1.1 [Topic] Care}
\phantomsection
\label{navigation:IV.1.1.1}
\addcontentsline{toc}{subsection}{IV.1.1.1 [Topic] Care}

\label{element:f702d357183847dd8f0fd09545679ffd}
\index{[Topic] Care (IV.1.1.1)}
{\small DataObject / Topic · Accepted bounded reconciliation · current}\par

\textbf{Description}: Information establishing the contexts within which healthcare occurs, including encounters, episodes, care-place associations, progression events and longitudinal relationships that provide continuity and context across care interactions.


{\footnotesize \href{../../../../continuum/application/data-object/topic/f702d357183847dd8f0fd09545679ffd.md}{Canonical record: f702d357183847dd8f0fd09545679ffd}; full qualifications and provenance.}\par
\subsubsection*{IV.1.1.1.1 [Topic] CareLocation}
\phantomsection
\label{navigation:IV.1.1.1.1}
\addcontentsline{toc}{subsubsection}{IV.1.1.1.1 [Topic] CareLocation}

\label{element:cfc241a7b6904bd19ccc1b9adbac0393}
\index{[Topic] CareLocation (IV.1.1.1.1)}
{\small DataObject / Topic · Accepted bounded reconciliation · current}\par

\textbf{Description}
Information describing the association between an Encounter and the place or places in which care occurs, including changes in location during the Encounter.


{\footnotesize \href{../../../../continuum/application/data-object/topic/cfc241a7b6904bd19ccc1b9adbac0393.md}{Canonical record: cfc241a7b6904bd19ccc1b9adbac0393}; full qualifications and provenance.}\par
\subsubsection*{IV.1.1.1.2 [Topic] Encounter}
\phantomsection
\label{navigation:IV.1.1.1.2}
\addcontentsline{toc}{subsubsection}{IV.1.1.1.2 [Topic] Encounter}

\label{element:fb464349f13b41c4ae9920a743fc6c12}
\index{[Topic] Encounter (IV.1.1.1.2)}
{\small DataObject / Topic · Explicitly APPROVED G2 decision only; G2 remains candidate/open and all residual qualifications remain unresolved. · current}\par

\textbf{Status: RESOLVED — APPROVED, with explicit unresolved further authority/evidence details.}



\textbf{Description}: Information describing a bounded interaction in which healthcare is delivered or healthcare-related activity occurs for an individual within a particular context.



\textbf{Definition}
An Encounter may provide the care-interaction context within which zero or more Healthcare Service Deliveries occur. A Healthcare Service Delivery may occur without an associated Encounter.



Encounter and HealthcareServiceDelivery are distinct Information Concepts. A single Encounter may contextualise multiple HealthcareServiceDelivery instances [0..*], and delivery may exist independently of Encounter context.



\begin{itemize}


\item Encounter ≠ HealthcareServiceDelivery. Encounter existence does not establish that any service was delivered.

\item The relationship is optional contextual association, not mandatory lifecycle progression, identity, containment or a component/state of Encounter.

\item The approved [0..*] multiplicity concerns deliveries contextualised by an Encounter; no further reverse cardinality is inferred.

\item Package1's DeliverableHealthcareService → optional Order → HealthcareServiceDelivery → ServiceOutcome semantics remain unchanged. Encounter is not made mandatory anywhere in that chain.

\item Further source/correlation qualification and universal delivery/outcome ownership are not established by this decision.

\end{itemize}


\textbf{Review consequence:} G2-Q05's Encounter/delivery existence, basic meaning and multiplicity questions are resolved. Further qualifications and the Referral portion remain unreviewed; no Referral Disposition review is commenced.



Source standing: Explicitly APPROVED G2 decision only; G2 remains candidate/open and all residual qualifications remain unresolved.


{\footnotesize \href{../../../../continuum/application/data-object/topic/fb464349f13b41c4ae9920a743fc6c12.md}{Canonical record: fb464349f13b41c4ae9920a743fc6c12}; full qualifications and provenance.}\par
\subsubsection*{IV.1.1.1.3 [Topic] EncounterContext}
\phantomsection
\label{navigation:IV.1.1.1.3}
\addcontentsline{toc}{subsubsection}{IV.1.1.1.3 [Topic] EncounterContext}

\label{element:a8be2b6474cd43cd9845b93258bd2fe3}
\index{[Topic] EncounterContext (IV.1.1.1.3)}
{\small DataObject / Topic · Accepted bounded reconciliation · current}\par

\textbf{Description}
Information bringing together the relevant subject, participants, services, places and other contextual information required to understand an Encounter.


{\footnotesize \href{../../../../continuum/application/data-object/topic/a8be2b6474cd43cd9845b93258bd2fe3.md}{Canonical record: a8be2b6474cd43cd9845b93258bd2fe3}; full qualifications and provenance.}\par
\subsubsection*{IV.1.1.1.4 [Topic] EncounterProgression}
\phantomsection
\label{navigation:IV.1.1.1.4}
\addcontentsline{toc}{subsubsection}{IV.1.1.1.4 [Topic] EncounterProgression}

\label{element:a3badcf73eaa4331b75d7b3eb6e67fc8}
\index{[Topic] EncounterProgression (IV.1.1.1.4)}
{\small DataObject / Topic · Accepted bounded reconciliation · current}\par

\textbf{Description}
Information describing changes in the administrative or operational progression of an Encounter over time, preserving the significant states and transitions associated with that interaction.


{\footnotesize \href{../../../../continuum/application/data-object/topic/a3badcf73eaa4331b75d7b3eb6e67fc8.md}{Canonical record: a3badcf73eaa4331b75d7b3eb6e67fc8}; full qualifications and provenance.}\par
\subsubsection*{IV.1.1.1.5 [Topic] Episode}
\phantomsection
\label{navigation:IV.1.1.1.5}
\addcontentsline{toc}{subsubsection}{IV.1.1.1.5 [Topic] Episode}

\label{element:3927d161556d48f4b2cc57bd46ab785d}
\index{[Topic] Episode (IV.1.1.1.5)}
{\small DataObject / Topic · Explicitly APPROVED G2 decision only; G2 remains candidate/open and all residual qualifications remain unresolved. · current}\par

\textbf{Status: RESOLVED — APPROVED, with explicit unresolved association/responsibility details.}



\textbf{Description}
Information describing a longitudinal context that relates healthcare activity over time for one or more health problems or conditions, providing continuity across relevant Encounters and other care interactions.



\textbf{Definition}
Episode is a longitudinal information context representing the treatment of one or more Problems or Conditions across a related series of Encounters, providing a coherent perspective for assurance, workflow and financial purposes.



Episode is now an independently meaningful Information Concept. The former non-authoritative Episode Context suggestion is promoted to Episode with this definition; its original rationale and standing remain historical in §7.



\begin{itemize}


\item Episode is not merely a Collection of Encounters. Encounter association contributes to longitudinal treatment context but does not define its complete meaning.

\item Episode relates one or more Problems/Conditions [1..\emph{] to a longitudinal course of treatment and may relate to multiple Encounters [0..}]. Not every Encounter for a healthcare subject necessarily belongs to a particular Episode.

\item Encounter-association mechanisms/rules, detailed identity, originating authority, Harmonia information-responsibility allocation and temporal/lifecycle semantics remain unresolved. The decision allocates no reverse cardinality, upstream owner, Feature or Service.

\item No Episode ownership of related Encounters is inferred; the association is not implementation containment. Assurance/workflow/financial perspective is not an allocation of authority for those activities.

\item This architecturally approved meaning stands independently of representation. No FHIR EpisodeOfCare derivation, mapping or equivalence is established.

\end{itemize}


\textbf{Review consequence:} G2-Q04 no longer asks whether Episode exists or what its basic meaning is; only genuinely unresolved details remain. No Problems/Conditions family is derived.



Source standing: Explicitly APPROVED G2 decision only; G2 remains candidate/open and all residual qualifications remain unresolved.


{\footnotesize \href{../../../../continuum/application/data-object/topic/3927d161556d48f4b2cc57bd46ab785d.md}{Canonical record: 3927d161556d48f4b2cc57bd46ab785d}; full qualifications and provenance.}\par
\subsection*{IV.1.1.2 [Topic] Coordination}
\phantomsection
\label{navigation:IV.1.1.2}
\addcontentsline{toc}{subsection}{IV.1.1.2 [Topic] Coordination}

\label{element:882234c291ca4c6492ad6dfbe1e23fea}
\index{[Topic] Coordination (IV.1.1.2)}
{\small DataObject / Topic · Accepted bounded reconciliation · current}\par

Information supporting the coordination of care between participants and providers, including referrals, their purpose and context, progression, supporting information and outcomes, without presuming that coordination necessarily transfers care responsibility.


{\footnotesize \href{../../../../continuum/application/data-object/topic/882234c291ca4c6492ad6dfbe1e23fea.md}{Canonical record: 882234c291ca4c6492ad6dfbe1e23fea}; full qualifications and provenance.}\par
\subsubsection*{IV.1.1.2.1 [Topic] Referral}
\phantomsection
\label{navigation:IV.1.1.2.1}
\addcontentsline{toc}{subsubsection}{IV.1.1.2.1 [Topic] Referral}

\label{element:d79d14c20bf14ecdab0434d4aa7db9b8}
\index{[Topic] Referral (IV.1.1.2.1)}
{\small DataObject / Topic · Accepted bounded reconciliation · current}\par

\textbf{Description}
Information describing a request for another healthcare participant or provider to assess, manage, provide care for, or assume responsibility concerning an individual.


{\footnotesize \href{../../../../continuum/application/data-object/topic/d79d14c20bf14ecdab0434d4aa7db9b8.md}{Canonical record: d79d14c20bf14ecdab0434d4aa7db9b8}; full qualifications and provenance.}\par
\subsubsection*{IV.1.1.2.2 [Topic] ReferralContext}
\phantomsection
\label{navigation:IV.1.1.2.2}
\addcontentsline{toc}{subsubsection}{IV.1.1.2.2 [Topic] ReferralContext}

\label{element:9eb5906e93ca4625bdb675a3397c4c8d}
\index{[Topic] ReferralContext (IV.1.1.2.2)}
{\small DataObject / Topic · Accepted bounded reconciliation · current}\par

\textbf{Description}
Information bringing together the subject, requester, recipient, purpose and other contextual information necessary to understand and act upon a Referral.


{\footnotesize \href{../../../../continuum/application/data-object/topic/9eb5906e93ca4625bdb675a3397c4c8d.md}{Canonical record: 9eb5906e93ca4625bdb675a3397c4c8d}; full qualifications and provenance.}\par
\subsubsection*{IV.1.1.2.3 [Topic] ReferralOutcome}
\phantomsection
\label{navigation:IV.1.1.2.3}
\addcontentsline{toc}{subsubsection}{IV.1.1.2.3 [Topic] ReferralOutcome}

\label{element:18ad2e58913b4a938e172df7b03f366a}
\index{[Topic] ReferralOutcome (IV.1.1.2.3)}
{\small DataObject / Topic · Explicitly APPROVED G2-D04; no promotion of Referral family/package or independent Referral Disposition. · current}\par

\textbf{Description}
Information describing the outcome or resolution of a Referral, including the resulting disposition and relevant information returned to participants involved in the referral.



\textbf{Status: RESOLVED — APPROVED. G2-Q01 is resolved; G2 remains open.}



Referral Disposition is \textbf{not retained as an independent Information Concept}. The assurance and governance requirement previously represented by that B candidate is instead represented by \textbf{Referral Outcome}.



\begin{quote}


\textbf{Referral Outcome is the governed information concerning what ultimately became of a Referral, sufficient to support assurance, governance and traceability over the Referral and its progression.}



\end{quote}


\textbf{Semantic classification:} Outcome; provenance-bearing information with Assertion characteristics where a source reports what became of the Referral. This is not ServiceOutcome equivalence or a downstream representation category.



Referral Outcome exists principally to support the ability to account for Referrals within the applicable governance/assurance scope. Where established and appropriate, it may include the resulting outcome, basis/reason, source/originator, authority, temporal context and relationships to relevant resulting activity or information. No canonical Referral Outcome value set is established.



\begin{quote}\small\ttfamily
Referral ≠ Referral Progression Event ≠ Referral Outcome\\
Referral Outcome ≠ Responsibility Assumption ≠ Transfer of Care ≠ HealthcareServiceDelivery ≠ ServiceOutcome\\
Referral information semantics ≠ Referral Administration workflow progression\\

\end{quote}


Referral Outcome may reference or relate to those Concepts where appropriate; its existence establishes none of them. Recording that a Referral was accepted, progressed or otherwise reached an outcome does not itself prove that care responsibility was actually assumed or transferred. G1 K1's intent/disposition/responsibility/Transfer of Care distinctions remain intact. Outcome is not equated with acceptance.



\textbf{Referral Progression Event:} Retained as historical information concerning significant occurrences in the administrative/business progression of a Referral. Not every Event represents a Referral Outcome, and the historical progression record is not collapsed into the final/current Outcome.



\textbf{Assurance / governance coverage — RF-IR5:} Harmonia must support assurance that Referrals within the applicable governance scope can be accounted for through sufficient progression and outcome information to establish what became of them. This requirement is applied to established Referral Administration / Manage Referral Progression / Referral Status \& Outcome Service responsibility without allocating a new owner, Feature, Service or Collection.



\begin{quote}\small\ttfamily
Referral population within governance scope → progression/outcome evidence → assurance coverage\\

\end{quote}


An inability to establish what became of a Referral is itself an assurance gap. No Referral Outcome is manufactured to make coverage appear complete. Where the outcome is not known, its absence is preserved explicitly.



\textbf{Workflow closure / finalisation:} This belongs to the progression of the Business Process / Referral Administration activity, not intrinsic Referral disposition information. The workflow may reach a state in which no further Referral Administration activity is required. This does not require a Referral Finalisation or Referral Closure Information Concept, a terminal state intrinsic to Referral, or Referral Outcome to double as workflow state. No workflow terminal states are defined and the Business Process state machine is not reproduced as the Information Concept lifecycle.



\textbf{Authority, temporal meaning and provenance:} Harmonia governs the outcome information for assurance and traceability; it does not infer substantive authority from recording or possessing it. Outcome may require provenance sufficient to establish what outcome was understood, source/originator, authority where relevant, when established/effective, and rationale/qualification where relevant. These are not all universally mandatory attributes. Detailed governance-scope criteria, evidence sufficiency, authority eligibility and contextual provenance/temporal qualification remain unestablished where existing architecture supplies no detail.



\textbf{Review history and consequence:} Q01 retains its original B question. The former Referral Disposition candidate proposed an independent Assertion for an attributable decision/reported disposition, with acceptance, decline, redirection and context-specific waitlisting examples. Architectural review rejects that separate Concept and reframes its assurance/governance purpose as Referral Outcome. Recipient-decision meanings and the canonical upstream Referral Disposition Log label are not erased. Q01 is resolved, without answering acceptance's clinical responsibility effects or the precise Referral–delivery/outcome relationships under Q05. Responsibility Assumption Assertion remains a non-authoritative D suggestion.



\textbf{Persistence boundary:} Only the Referral family, this review record and the necessary Referral navigation text are reconciled. Q02–Q07 and the complete G2-D01–D03 record remain unchanged. Frozen upstream architecture, Foundation, Package1 and initial historical report/diff are not modified. No other G2 question, G3, later Package2 family, FHIR mapping or downstream representation is commenced. G1 remains CLOSED; G2 is not CLOSED.



Source standing: Explicitly APPROVED G2-D04; no promotion of Referral family/package or independent Referral Disposition.


{\footnotesize \href{../../../../continuum/application/data-object/topic/18ad2e58913b4a938e172df7b03f366a.md}{Canonical record: 18ad2e58913b4a938e172df7b03f366a}; full qualifications and provenance.}\par
\subsubsection*{IV.1.1.2.4 [Topic] ReferralProgression}
\phantomsection
\label{navigation:IV.1.1.2.4}
\addcontentsline{toc}{subsubsection}{IV.1.1.2.4 [Topic] ReferralProgression}

\label{element:bc70a1055ca34eed8143856816587f67}
\index{[Topic] ReferralProgression (IV.1.1.2.4)}
{\small DataObject / Topic · G2-D04 explicitly retains significant historical progression information; historical describes represented occurrences, not historical architectural standing. · current}\par

\textbf{Description} 
Information describing the progression of a Referral over time, including significant changes in its state, handling or disposition.



\textbf{Status: RESOLVED — APPROVED. G2-Q01 is resolved; G2 remains open.}



Referral Disposition is \textbf{not retained as an independent Information Concept}. The assurance and governance requirement previously represented by that B candidate is instead represented by \textbf{Referral Outcome}.



\begin{quote}


\textbf{Referral Outcome is the governed information concerning what ultimately became of a Referral, sufficient to support assurance, governance and traceability over the Referral and its progression.}



\end{quote}


\textbf{Semantic classification:} Outcome; provenance-bearing information with Assertion characteristics where a source reports what became of the Referral. This is not ServiceOutcome equivalence or a downstream representation category.



Referral Outcome exists principally to support the ability to account for Referrals within the applicable governance/assurance scope. Where established and appropriate, it may include the resulting outcome, basis/reason, source/originator, authority, temporal context and relationships to relevant resulting activity or information. No canonical Referral Outcome value set is established.



\begin{quote}\small\ttfamily
Referral ≠ Referral Progression Event ≠ Referral Outcome\\
Referral Outcome ≠ Responsibility Assumption ≠ Transfer of Care ≠ HealthcareServiceDelivery ≠ ServiceOutcome\\
Referral information semantics ≠ Referral Administration workflow progression\\

\end{quote}


Referral Outcome may reference or relate to those Concepts where appropriate; its existence establishes none of them. Recording that a Referral was accepted, progressed or otherwise reached an outcome does not itself prove that care responsibility was actually assumed or transferred. G1 K1's intent/disposition/responsibility/Transfer of Care distinctions remain intact. Outcome is not equated with acceptance.



\textbf{Referral Progression Event:} Retained as historical information concerning significant occurrences in the administrative/business progression of a Referral. Not every Event represents a Referral Outcome, and the historical progression record is not collapsed into the final/current Outcome.



\textbf{Assurance / governance coverage — RF-IR5:} Harmonia must support assurance that Referrals within the applicable governance scope can be accounted for through sufficient progression and outcome information to establish what became of them. This requirement is applied to established Referral Administration / Manage Referral Progression / Referral Status \& Outcome Service responsibility without allocating a new owner, Feature, Service or Collection.



\begin{quote}\small\ttfamily
Referral population within governance scope → progression/outcome evidence → assurance coverage\\

\end{quote}


An inability to establish what became of a Referral is itself an assurance gap. No Referral Outcome is manufactured to make coverage appear complete. Where the outcome is not known, its absence is preserved explicitly.



\textbf{Workflow closure / finalisation:} This belongs to the progression of the Business Process / Referral Administration activity, not intrinsic Referral disposition information. The workflow may reach a state in which no further Referral Administration activity is required. This does not require a Referral Finalisation or Referral Closure Information Concept, a terminal state intrinsic to Referral, or Referral Outcome to double as workflow state. No workflow terminal states are defined and the Business Process state machine is not reproduced as the Information Concept lifecycle.



\textbf{Authority, temporal meaning and provenance:} Harmonia governs the outcome information for assurance and traceability; it does not infer substantive authority from recording or possessing it. Outcome may require provenance sufficient to establish what outcome was understood, source/originator, authority where relevant, when established/effective, and rationale/qualification where relevant. These are not all universally mandatory attributes. Detailed governance-scope criteria, evidence sufficiency, authority eligibility and contextual provenance/temporal qualification remain unestablished where existing architecture supplies no detail.



\textbf{Review history and consequence:} Q01 retains its original B question. The former Referral Disposition candidate proposed an independent Assertion for an attributable decision/reported disposition, with acceptance, decline, redirection and context-specific waitlisting examples. Architectural review rejects that separate Concept and reframes its assurance/governance purpose as Referral Outcome. Recipient-decision meanings and the canonical upstream Referral Disposition Log label are not erased. Q01 is resolved, without answering acceptance's clinical responsibility effects or the precise Referral–delivery/outcome relationships under Q05. Responsibility Assumption Assertion remains a non-authoritative D suggestion.



\textbf{Persistence boundary:} Only the Referral family, this review record and the necessary Referral navigation text are reconciled. Q02–Q07 and the complete G2-D01–D03 record remain unchanged. Frozen upstream architecture, Foundation, Package1 and initial historical report/diff are not modified. No other G2 question, G3, later Package2 family, FHIR mapping or downstream representation is commenced. G1 remains CLOSED; G2 is not CLOSED.



Source standing: G2-D04 explicitly retains significant historical progression information; historical describes represented occurrences, not historical architectural standing.


{\footnotesize \href{../../../../continuum/application/data-object/topic/bc70a1055ca34eed8143856816587f67.md}{Canonical record: bc70a1055ca34eed8143856816587f67}; full qualifications and provenance.}\par
\subsection*{IV.1.1.3 [Topic] Enterprise}
\phantomsection
\label{navigation:IV.1.1.3}
\addcontentsline{toc}{subsection}{IV.1.1.3 [Topic] Enterprise}

\label{element:5c5da6a06dac413aa2809a1397f8a3b9}
\index{[Topic] Enterprise (IV.1.1.3)}
{\small DataObject / Topic · Accepted bounded reconciliation · current}\par

Information concerning healthcare organisations and organisational units, including their identity, classification, structure, profile, relationships, and contact information.


{\footnotesize \href{../../../../continuum/application/data-object/topic/5c5da6a06dac413aa2809a1397f8a3b9.md}{Canonical record: 5c5da6a06dac413aa2809a1397f8a3b9}; full qualifications and provenance.}\par
\subsubsection*{IV.1.1.3.1 [Topic] Organisation}
\phantomsection
\label{navigation:IV.1.1.3.1}
\addcontentsline{toc}{subsubsection}{IV.1.1.3.1 [Topic] Organisation}

\label{element:dd3ba70128874caeb2d5da419d006c49}
\index{[Topic] Organisation (IV.1.1.3.1)}
{\small DataObject / Topic · Established Package1 conceptual information meaning; residual qualifications in selected section remain unchanged. · current}\par

\textbf{Description}
Information describing an organisation participating in healthcare, including its identity, characteristics and role within the healthcare environment.



\begin{itemize}


\item \textbf{Semantic Classification}: \texttt{\seqsplit{Entity}}

\item \textbf{Definition}: An identifiable legal, corporate, or administrative organisation operating within the healthcare domain (e.g. Health Service District, Hospital Network, Private Healthcare Provider, Diagnostic Laboratory Enterprise).

\item \textbf{Architectural Scope}: Represents the overarching enterprise entity possessing legal identity, operational mandate, and administrative responsibility.

\end{itemize}


Source standing: Established Package1 conceptual information meaning; residual qualifications in selected section remain unchanged.


{\footnotesize \href{../../../../continuum/application/data-object/topic/dd3ba70128874caeb2d5da419d006c49.md}{Canonical record: dd3ba70128874caeb2d5da419d006c49}; full qualifications and provenance.}\par
\subsubsection*{IV.1.1.3.2 [Topic] OrganisationalUnit}
\phantomsection
\label{navigation:IV.1.1.3.2}
\addcontentsline{toc}{subsubsection}{IV.1.1.3.2 [Topic] OrganisationalUnit}

\label{element:2c884cdd6a4b41a184818e3a7e540897}
\index{[Topic] OrganisationalUnit (IV.1.1.3.2)}
{\small DataObject / Topic · Established Package1 conceptual information meaning; residual qualifications in selected section remain unchanged. · current}\par

\textbf{Description}
Information describing a defined subdivision of an organisation and its position within the organisation's structural context.



\begin{itemize}


\item \textbf{Semantic Classification}: \texttt{\seqsplit{Entity / Administrative Subdivision}}

\item \textbf{Definition}: An identifiable internal division, clinical directorate, operational department, or specialty unit within a healthcare organisation (e.g. \emph{Division of Surgery}, \emph{Department of Cardiology}, \emph{Pharmacy Department}).

\item \textbf{Architectural Scope}: Represents structured administrative subdivisions that manage resources, staff, and services under the umbrella of the parent organisation.

\end{itemize}


Source standing: Established Package1 conceptual information meaning; residual qualifications in selected section remain unchanged.


{\footnotesize \href{../../../../continuum/application/data-object/topic/2c884cdd6a4b41a184818e3a7e540897.md}{Canonical record: 2c884cdd6a4b41a184818e3a7e540897}; full qualifications and provenance.}\par
\subsubsection*{IV.1.1.3.3 [Topic] OrganisationClassification}
\phantomsection
\label{navigation:IV.1.1.3.3}
\addcontentsline{toc}{subsubsection}{IV.1.1.3.3 [Topic] OrganisationClassification}

\label{element:227974cf26054bdea17038ecd6c0aa3a}
\index{[Topic] OrganisationClassification (IV.1.1.3.3)}
{\small DataObject / Topic · Established Package1 conceptual information meaning; residual qualifications in selected section remain unchanged. · current}\par

\textbf{Description}
Information describing the classifications applicable to an organisation, including its type, function or other governed categorisation.



\begin{itemize}


\item \textbf{Semantic Classification}: \texttt{\seqsplit{Assertion / Classification}}

\item \textbf{Definition}: The categorisation of the organisation according to its structural, operational, and clinical mandate (e.g. \emph{Tertiary Hospital}, \emph{Regional Health District}, \emph{Ambulatory Primary Care Clinic}, \emph{Specialist Pathology Service}).

\end{itemize}


Source standing: Established Package1 conceptual information meaning; residual qualifications in selected section remain unchanged.


{\footnotesize \href{../../../../continuum/application/data-object/topic/227974cf26054bdea17038ecd6c0aa3a.md}{Canonical record: 227974cf26054bdea17038ecd6c0aa3a}; full qualifications and provenance.}\par
\subsubsection*{IV.1.1.3.4 [Topic] OrganisationContact}
\phantomsection
\label{navigation:IV.1.1.3.4}
\addcontentsline{toc}{subsubsection}{IV.1.1.3.4 [Topic] OrganisationContact}

\label{element:cad820b906fa40f28824278fb18b7497}
\index{[Topic] OrganisationContact (IV.1.1.3.4)}
{\small DataObject / Topic · Established Package1 conceptual information meaning; residual qualifications in selected section remain unchanged. · current}\par

\textbf{Description}
Information describing how an organisation may be contacted or communicated with, including applicable addresses and communication channels.



\begin{itemize}


\item \textbf{Semantic Classification}: \texttt{\seqsplit{Assertion / Contact Bundle}}

\item \textbf{Definition}: The official administrative communication details, registered physical and postal addresses, corporate telephone/electronic points of contact, and nominated official representatives.

\end{itemize}


Source standing: Established Package1 conceptual information meaning; residual qualifications in selected section remain unchanged.


{\footnotesize \href{../../../../continuum/application/data-object/topic/cad820b906fa40f28824278fb18b7497.md}{Canonical record: cad820b906fa40f28824278fb18b7497}; full qualifications and provenance.}\par
\subsubsection*{IV.1.1.3.5 [Topic] OrganisationIdentity}
\phantomsection
\label{navigation:IV.1.1.3.5}
\addcontentsline{toc}{subsubsection}{IV.1.1.3.5 [Topic] OrganisationIdentity}

\label{element:c240be16de76415682237a238b3d3fc8}
\index{[Topic] OrganisationIdentity (IV.1.1.3.5)}
{\small DataObject / Topic · Established Package1 conceptual information meaning; residual qualifications in selected section remain unchanged. · current}\par

\textbf{Description}
Information used to identify, distinguish and correlate an organisation across authorities, systems and healthcare contexts.



\begin{itemize}


\item \textbf{Semantic Classification}: \texttt{\seqsplit{Assertion / Identity Bundle}}

\item \textbf{Definition}: The governed set of official enterprise identifiers, registry bindings, and external directory identifiers establishing the identity of the organisation.

\item \textbf{Key Semantic Decomposition}:

\item \texttt{\seqsplit{Identifier Value}}: The token or registration number identifying the organisation.

\item \texttt{\seqsplit{Identifier Authority / Namespace}}: The statutory authority, national healthcare directory, or enterprise namespace issuing the identifier.

\end{itemize}


Source standing: Established Package1 conceptual information meaning; residual qualifications in selected section remain unchanged.


{\footnotesize \href{../../../../continuum/application/data-object/topic/c240be16de76415682237a238b3d3fc8.md}{Canonical record: c240be16de76415682237a238b3d3fc8}; full qualifications and provenance.}\par
\subsubsection*{IV.1.1.3.6 [Topic] OrganisationProfile}
\phantomsection
\label{navigation:IV.1.1.3.6}
\addcontentsline{toc}{subsubsection}{IV.1.1.3.6 [Topic] OrganisationProfile}

\label{element:04adbfc449bb4b9e9ae6e7fdf8a66bf9}
\index{[Topic] OrganisationProfile (IV.1.1.3.6)}
{\small DataObject / Topic · Established Package1 conceptual information meaning; residual qualifications in selected section remain unchanged. · current}\par

\textbf{Description}
Information describing the characteristics of an organisation relevant to its participation in healthcare, including descriptive and operational attributes.



\begin{itemize}


\item \textbf{Semantic Classification}: \texttt{\seqsplit{Assertion / Profile Bundle}}

\item \textbf{Definition}: The verified set of administrative and corporate attributes describing the organisation, including registered legal name, public trading names, corporate registration numbers, and formal accreditation status.

\end{itemize}


Source standing: Established Package1 conceptual information meaning; residual qualifications in selected section remain unchanged.


{\footnotesize \href{../../../../continuum/application/data-object/topic/04adbfc449bb4b9e9ae6e7fdf8a66bf9.md}{Canonical record: 04adbfc449bb4b9e9ae6e7fdf8a66bf9}; full qualifications and provenance.}\par
\subsection*{IV.1.1.4 [Topic] Equipment}
\phantomsection
\label{navigation:IV.1.1.4}
\addcontentsline{toc}{subsection}{IV.1.1.4 [Topic] Equipment}

\label{element:8e658834311d41e68a8ecf71d8ccd5db}
\index{[Topic] Equipment (IV.1.1.4)}
{\small DataObject / Topic · Accepted bounded reconciliation · current}\par

\textbf{Description}
Information concerning clinical devices and equipment, including device definitions, physical instances, identity, operational associations, location, subject and service relationships, and communication endpoints.


{\footnotesize \href{../../../../continuum/application/data-object/topic/8e658834311d41e68a8ecf71d8ccd5db.md}{Canonical record: 8e658834311d41e68a8ecf71d8ccd5db}; full qualifications and provenance.}\par
\subsubsection*{IV.1.1.4.1 [Topic] Device}
\phantomsection
\label{navigation:IV.1.1.4.1}
\addcontentsline{toc}{subsubsection}{IV.1.1.4.1 [Topic] Device}

\label{element:25c1b751956446cea17065b6254419d5}
\index{[Topic] Device (IV.1.1.4.1)}
{\small DataObject / Topic · Established Package1 conceptual information meaning; residual qualifications in selected section remain unchanged. · current}\par

\textbf{Description}
Information describing an identifiable physical device or item of equipment, including its identity, characteristics and relationship to its defining device type.



\begin{itemize}


\item \textbf{Semantic Classification}: \texttt{\seqsplit{Entity}}

\item \textbf{Definition}: An identifiable, physical hardware asset manufactured as an instance of a \texttt{\seqsplit{Device Definition}}.

\item \textbf{Key Characteristics}: Device Identifier (such as physical serial number or asset identification), manufacturing date, device version / configuration assertion, and asset operational status (with schemes such as Unique Device Identifier [UDI] and enterprise asset tags serving as illustrative examples).

\end{itemize}


Source standing: Established Package1 conceptual information meaning; residual qualifications in selected section remain unchanged.


{\footnotesize \href{../../../../continuum/application/data-object/topic/25c1b751956446cea17065b6254419d5.md}{Canonical record: 25c1b751956446cea17065b6254419d5}; full qualifications and provenance.}\par
\subsubsection*{IV.1.1.4.2 [Topic] DeviceAssociation}
\phantomsection
\label{navigation:IV.1.1.4.2}
\addcontentsline{toc}{subsubsection}{IV.1.1.4.2 [Topic] DeviceAssociation}

\label{element:303f8586953d450d8d82d3d5f76e488a}
\index{[Topic] DeviceAssociation (IV.1.1.4.2)}
{\small DataObject / Topic · Established Package1 conceptual information meaning; residual qualifications in selected section remain unchanged. · current}\par

\textbf{Description}
Information describing the contextual association of a device with a subject, location, service or other entity, including the period for which that association applies.



\begin{itemize}


\item \textbf{Semantic Classification}: \texttt{\seqsplit{Contextual / Temporal Relationship}}

\item \textbf{Definition}: A dynamic, time-bounded association binding a \texttt{\seqsplit{Device Instance}} to a patient, physical care-place, or clinical service.

\item \textbf{Key Characteristics}: Association type, effective start and end timestamps, commissioning authority, and clinical context.

\end{itemize}


Source standing: Established Package1 conceptual information meaning; residual qualifications in selected section remain unchanged.


{\footnotesize \href{../../../../continuum/application/data-object/topic/303f8586953d450d8d82d3d5f76e488a.md}{Canonical record: 303f8586953d450d8d82d3d5f76e488a}; full qualifications and provenance.}\par
\subsubsection*{IV.1.1.4.3 [Topic] DeviceDefinition}
\phantomsection
\label{navigation:IV.1.1.4.3}
\addcontentsline{toc}{subsubsection}{IV.1.1.4.3 [Topic] DeviceDefinition}

\label{element:2a879d06b17143be8bbfc35cdaab8bdf}
\index{[Topic] DeviceDefinition (IV.1.1.4.3)}
{\small DataObject / Topic · Established Package1 conceptual information meaning; residual qualifications in selected section remain unchanged. · current}\par

\textbf{Description}
Information describing a type or model of device, including the characteristics and specifications shared by device instances of that definition.



\begin{itemize}


\item \textbf{Semantic Classification}: \texttt{\seqsplit{Definition}}

\item \textbf{Definition}: The governed type, manufacturer catalogue model, regulatory classification, and technical specification of a medical device (e.g. \emph{Infusion Pump Model XYZ}, \emph{12-Lead ECG Monitor Model ABC}).

\item \textbf{Key Characteristics}: Manufacturer name, model identifier, commercial name, regulatory classification (e.g. medical device risk class), calibration requirements, and operational capabilities.

\end{itemize}


Source standing: Established Package1 conceptual information meaning; residual qualifications in selected section remain unchanged.


{\footnotesize \href{../../../../continuum/application/data-object/topic/2a879d06b17143be8bbfc35cdaab8bdf.md}{Canonical record: 2a879d06b17143be8bbfc35cdaab8bdf}; full qualifications and provenance.}\par
\subsection*{IV.1.1.5 [Topic] Individual}
\phantomsection
\label{navigation:IV.1.1.5}
\addcontentsline{toc}{subsection}{IV.1.1.5 [Topic] Individual}

\label{element:4af969b8c23a4dcab7569dc51b263383}
\index{[Topic] Individual (IV.1.1.5)}
{\small DataObject / Topic · Accepted bounded reconciliation · current}\par

\textbf{Description}: Information concerning an individual, their identity, identifiers, demographic characteristics, relationships, and their contextual participation as a subject of healthcare.


{\footnotesize \href{../../../../continuum/application/data-object/topic/4af969b8c23a4dcab7569dc51b263383.md}{Canonical record: 4af969b8c23a4dcab7569dc51b263383}; full qualifications and provenance.}\par
\subsubsection*{IV.1.1.5.1 [Topic] Authority}
\phantomsection
\label{navigation:IV.1.1.5.1}
\addcontentsline{toc}{subsubsection}{IV.1.1.5.1 [Topic] Authority}

\label{element:99dc695355384ea29ac3515b9178dfd0}
\index{[Topic] Authority (IV.1.1.5.1)}
{\small DataObject / Topic · Accepted bounded reconciliation · current}\par

\textbf{Description}: Information describing another party's authority to represent, consent for, make decisions for, or otherwise act on behalf of an individual.


{\footnotesize \href{../../../../continuum/application/data-object/topic/99dc695355384ea29ac3515b9178dfd0.md}{Canonical record: 99dc695355384ea29ac3515b9178dfd0}; full qualifications and provenance.}\par
\subsubsection*{IV.1.1.5.2 [Topic] Contact}
\phantomsection
\label{navigation:IV.1.1.5.2}
\addcontentsline{toc}{subsubsection}{IV.1.1.5.2 [Topic] Contact}

\label{element:5c33b01e23964ef5935be79364e7c0f9}
\index{[Topic] Contact (IV.1.1.5.2)}
{\small DataObject / Topic · Accepted bounded reconciliation · current}\par

\textbf{Description}: Information describing how an individual may be contacted or communicated with, including applicable addresses and communication channels.


{\footnotesize \href{../../../../continuum/application/data-object/topic/5c33b01e23964ef5935be79364e7c0f9.md}{Canonical record: 5c33b01e23964ef5935be79364e7c0f9}; full qualifications and provenance.}\par
\subsubsection*{IV.1.1.5.3 [Topic] Demographics}
\phantomsection
\label{navigation:IV.1.1.5.3}
\addcontentsline{toc}{subsubsection}{IV.1.1.5.3 [Topic] Demographics}

\label{element:95a13e957f5c471797869b901085266c}
\index{[Topic] Demographics (IV.1.1.5.3)}
{\small DataObject / Topic · Accepted bounded reconciliation · current}\par

\textbf{Description}: Information describing the demographic characteristics of an individual relevant to their identification, administration, healthcare or communication.


{\footnotesize \href{../../../../continuum/application/data-object/topic/95a13e957f5c471797869b901085266c.md}{Canonical record: 95a13e957f5c471797869b901085266c}; full qualifications and provenance.}\par
\subsubsection*{IV.1.1.5.4 [Topic] Identity}
\phantomsection
\label{navigation:IV.1.1.5.4}
\addcontentsline{toc}{subsubsection}{IV.1.1.5.4 [Topic] Identity}

\label{element:3b72456b5686449f95ad26bfbbae07f3}
\index{[Topic] Identity (IV.1.1.5.4)}
{\small DataObject / Topic · Established Package1 conceptual information meaning; residual qualifications in selected section remain unchanged. · current}\par

\textbf{Definition}: The governed set of identity traits, identifiers, names, and demographic assertions that establish the recognised identity of a \texttt{\seqsplit{Person}}.



\textbf{Description}: Information used to establish, distinguish and correlate an individual across identities, identifiers, aliases, sources and healthcare contexts.



\textbf{Semantic Classification}: \texttt{\seqsplit{Assertion / Identity Bundle}}



\textbf{Architectural Scope}: Governed by the \texttt{\seqsplit{Person Identity}} capability. Acts as the governed identity anchor linking identifiers, aliases, and correlation graphs.



Source standing: Established Package1 conceptual information meaning; residual qualifications in selected section remain unchanged.


{\footnotesize \href{../../../../continuum/application/data-object/topic/3b72456b5686449f95ad26bfbbae07f3.md}{Canonical record: 3b72456b5686449f95ad26bfbbae07f3}; full qualifications and provenance.}\par
\subsubsection*{IV.1.1.5.5 [Topic] Patient}
\phantomsection
\label{navigation:IV.1.1.5.5}
\addcontentsline{toc}{subsubsection}{IV.1.1.5.5 [Topic] Patient}

\label{element:e92ea15616ea43c9a1d698e34ed1376b}
\index{[Topic] Patient (IV.1.1.5.5)}
{\small DataObject / Topic · Established Package1 conceptual information meaning; residual qualifications in selected section remain unchanged. · current}\par

\textbf{Definition} 
The subject-oriented context in which health information, clinical activity, or care delivery concerns a \texttt{\seqsplit{Person}}.



\textbf{Description} 
Information describing an individual in their context as a recipient or subject of healthcare, including the healthcare-specific identity and administrative context required to support care.



\textbf{Semantic Classification}: \texttt{\seqsplit{Contextual Binding}}



\textbf{Key Characteristics}: Demographic history, language/interpreter requirements, cultural and spiritual preferences, communication preferences, and vital status assertions (e.g., verified birth date, death notification).



\textbf{Architectural Scope}: Captures the individual's profile as a consumer of health services without assuming active patient status in a specific facility.



Source standing: Established Package1 conceptual information meaning; residual qualifications in selected section remain unchanged.


{\footnotesize \href{../../../../continuum/application/data-object/topic/e92ea15616ea43c9a1d698e34ed1376b.md}{Canonical record: e92ea15616ea43c9a1d698e34ed1376b}; full qualifications and provenance.}\par
\subsubsection*{IV.1.1.5.6 [Topic] Person}
\phantomsection
\label{navigation:IV.1.1.5.6}
\addcontentsline{toc}{subsubsection}{IV.1.1.5.6 [Topic] Person}

\label{element:1163ffac87814f319578ef5cb95dcd4d}
\index{[Topic] Person (IV.1.1.5.6)}
{\small DataObject / Topic · Established Package1 conceptual information meaning; residual qualifications in selected section remain unchanged. · current}\par

\textbf{Definition}: The identifiable biological and legal human entity.



\textbf{Description}: Information describing an individual as an enduring human entity, independent of any particular healthcare interaction, role or episode of care.



\textbf{Semantic Classification}: \texttt{\seqsplit{Entity}}



\textbf{Architectural Scope}: Represents the enduring human being independently of any specific clinical encounter, healthcare facility, or operational role. A \texttt{\seqsplit{Person}} exists prior to and outside of any interaction with the healthcare system.



Source standing: Established Package1 conceptual information meaning; residual qualifications in selected section remain unchanged.


{\footnotesize \href{../../../../continuum/application/data-object/topic/1163ffac87814f319578ef5cb95dcd4d.md}{Canonical record: 1163ffac87814f319578ef5cb95dcd4d}; full qualifications and provenance.}\par
\subsubsection*{IV.1.1.5.7 [Topic] Relationship}
\phantomsection
\label{navigation:IV.1.1.5.7}
\addcontentsline{toc}{subsubsection}{IV.1.1.5.7 [Topic] Relationship}

\label{element:c0ec3c75d9c54b3097e12c240ee816fe}
\index{[Topic] Relationship (IV.1.1.5.7)}
{\small DataObject / Topic · Accepted bounded reconciliation · current}\par

\textbf{Description}: Information describing relevant personal, family, carer, representative or other relationships between an individual and other people.


{\footnotesize \href{../../../../continuum/application/data-object/topic/c0ec3c75d9c54b3097e12c240ee816fe.md}{Canonical record: c0ec3c75d9c54b3097e12c240ee816fe}; full qualifications and provenance.}\par
\subsection*{IV.1.1.6 [Topic] Ordering}
\phantomsection
\label{navigation:IV.1.1.6}
\addcontentsline{toc}{subsection}{IV.1.1.6 [Topic] Ordering}

\label{element:1c4e2d5ce4574856a7d3605838e8afcb}
\index{[Topic] Ordering (IV.1.1.6)}
{\small DataObject / Topic · Accepted bounded reconciliation · current}\par

\textbf{Description}
Information concerning formal requests for services or activities, including their intent, routing, progression, modification, contextual information and correlation with resulting information or fulfilment.


{\footnotesize \href{../../../../continuum/application/data-object/topic/1c4e2d5ce4574856a7d3605838e8afcb.md}{Canonical record: 1c4e2d5ce4574856a7d3605838e8afcb}; full qualifications and provenance.}\par
\subsubsection*{IV.1.1.6.1 [Topic] Order}
\phantomsection
\label{navigation:IV.1.1.6.1}
\addcontentsline{toc}{subsubsection}{IV.1.1.6.1 [Topic] Order}

\label{element:2a791e0e1f6b45c0a57fb8233f73faa5}
\index{[Topic] Order (IV.1.1.6.1)}
{\small DataObject / Topic · Established Package1 conceptual information meaning; residual qualifications in selected section remain unchanged. · current}\par

\textbf{Description}
Information describing a formal request for a service, activity or other fulfilment to be performed, including its subject, requester, intent and requested outcome.



Harmonia enforces an essential architectural rule regarding clinical requisitions and service delivery:



\$\$\textbackslash{}text\{Order\} \textbackslash{}neq \textbackslash{}text\{Healthcare Service\}\$\$



\begin{quote}\small\ttfamily
graph TD\\
    ORDER["Order<br/>(Authoritative Requisition / Directive)"]\\
    DELIV["DeliverableHealthcareService<br/>(Bound Service Context)"]\\
    DELIVERY["HealthcareServiceDelivery<br/>(Information Referring to Service Execution)"]\\
\\
    DELIV -.->|"may be requested / directed via"| ORDER\\
    ORDER -.->|"requests / directs"| DELIVERY\\
    DELIV -->|"fulfilled by"| DELIVERY\\

\end{quote}


\par\medskip\textbf{Key Semantics \& Invariants}\par


\begin{enumerate}


\item \textbf{Associated Request Mechanism}: An \texttt{\seqsplit{Order}} is an authoritative clinical requisition, directive, or prescription requesting the delivery of a service. It is \textbf{NOT} a mandatory lifecycle stage of \texttt{\seqsplit{Healthcare Service}}.

\item \textbf{Delivery Without an Order}: Service delivery may legitimately occur without an \texttt{\seqsplit{Order}} where business or clinical semantics permit (e.g. \emph{Emergency Resuscitation}, \emph{Triage Assessment}, \emph{Direct Routine Nursing Care}, \emph{Scheduled Facility Maintenance}).

\item \textbf{Independent Lifecycles}: An \texttt{\seqsplit{Order}} possesses its own lifecycle states (\emph{Requested}, \emph{Accepted}, \emph{In-Progress}, \emph{Completed}, \emph{Discontinued}), owned by \texttt{\seqsplit{Order Administration}}, distinct from the lifecycle of the service delivery itself.

\end{enumerate}


Source standing: Established Package1 conceptual information meaning; residual qualifications in selected section remain unchanged.


{\footnotesize \href{../../../../continuum/application/data-object/topic/2a791e0e1f6b45c0a57fb8233f73faa5.md}{Canonical record: 2a791e0e1f6b45c0a57fb8233f73faa5}; full qualifications and provenance.}\par
\subsubsection*{IV.1.1.6.2 [Topic] OrderChange}
\phantomsection
\label{navigation:IV.1.1.6.2}
\addcontentsline{toc}{subsubsection}{IV.1.1.6.2 [Topic] OrderChange}

\label{element:73abd9430cfa4395ac38b41acfb169ec}
\index{[Topic] OrderChange (IV.1.1.6.2)}
{\small DataObject / Topic · Accepted bounded reconciliation · current}\par

Information describing a requested modification to an existing Order, including the nature, context and disposition of the requested change.


{\footnotesize \href{../../../../continuum/application/data-object/topic/73abd9430cfa4395ac38b41acfb169ec.md}{Canonical record: 73abd9430cfa4395ac38b41acfb169ec}; full qualifications and provenance.}\par
\subsubsection*{IV.1.1.6.3 [Topic] OrderContext}
\phantomsection
\label{navigation:IV.1.1.6.3}
\addcontentsline{toc}{subsubsection}{IV.1.1.6.3 [Topic] OrderContext}

\label{element:614c48fd853643578788b738f8017b53}
\index{[Topic] OrderContext (IV.1.1.6.3)}
{\small DataObject / Topic · Accepted bounded reconciliation · current}\par

Information bringing together the subject, requester, requested activity, intended performer and other context necessary to understand and act upon an Order.


{\footnotesize \href{../../../../continuum/application/data-object/topic/614c48fd853643578788b738f8017b53.md}{Canonical record: 614c48fd853643578788b738f8017b53}; full qualifications and provenance.}\par
\subsubsection*{IV.1.1.6.4 [Topic] OrderProgressionEvent}
\phantomsection
\label{navigation:IV.1.1.6.4}
\addcontentsline{toc}{subsubsection}{IV.1.1.6.4 [Topic] OrderProgressionEvent}

\label{element:8fb3a767d7a04d638ef76f98d61496b1}
\index{[Topic] OrderProgressionEvent (IV.1.1.6.4)}
{\small DataObject / Topic · Accepted bounded reconciliation · current}\par

\textbf{Description}
Information describing significant changes in the state and handling of an Order as it progresses from request through fulfilment or other resolution.


{\footnotesize \href{../../../../continuum/application/data-object/topic/8fb3a767d7a04d638ef76f98d61496b1.md}{Canonical record: 8fb3a767d7a04d638ef76f98d61496b1}; full qualifications and provenance.}\par
\subsubsection*{IV.1.1.6.5 [Topic] OrderResult}
\phantomsection
\label{navigation:IV.1.1.6.5}
\addcontentsline{toc}{subsubsection}{IV.1.1.6.5 [Topic] OrderResult}

\label{element:24187967f53446479c87d7676c66ba12}
\index{[Topic] OrderResult (IV.1.1.6.5)}
{\small DataObject / Topic · Accepted bounded reconciliation · current}\par

\textbf{Description}
Information establishing the relationship between an Order and information produced as a consequence of its fulfilment.


{\footnotesize \href{../../../../continuum/application/data-object/topic/24187967f53446479c87d7676c66ba12.md}{Canonical record: 24187967f53446479c87d7676c66ba12}; full qualifications and provenance.}\par
\subsubsection*{IV.1.1.6.6 [Topic] OrderRouting}
\phantomsection
\label{navigation:IV.1.1.6.6}
\addcontentsline{toc}{subsubsection}{IV.1.1.6.6 [Topic] OrderRouting}

\label{element:06bde79727884b27b3599e373b3cee5f}
\index{[Topic] OrderRouting (IV.1.1.6.6)}
{\small DataObject / Topic · Accepted bounded reconciliation · current}\par

\textbf{Description}
Information describing where and how an Order is to be directed for acceptance, processing or fulfilment.


{\footnotesize \href{../../../../continuum/application/data-object/topic/06bde79727884b27b3599e373b3cee5f.md}{Canonical record: 06bde79727884b27b3599e373b3cee5f}; full qualifications and provenance.}\par
\subsection*{IV.1.1.7 [Topic] Place}
\phantomsection
\label{navigation:IV.1.1.7}
\addcontentsline{toc}{subsection}{IV.1.1.7 [Topic] Place}

\label{element:36539e3eee7041e585aa3c5a06424bcc}
\index{[Topic] Place (IV.1.1.7)}
{\small DataObject / Topic · Accepted bounded reconciliation · current}\par

\textbf{Description}
Information concerning physical, geospatial, functional and care locations, including their identity, classification, spatial relationships, containment, characteristics, and care-place definitions.


{\footnotesize \href{../../../../continuum/application/data-object/topic/36539e3eee7041e585aa3c5a06424bcc.md}{Canonical record: 36539e3eee7041e585aa3c5a06424bcc}; full qualifications and provenance.}\par
\subsubsection*{IV.1.1.7.1 [Topic] CarePlace}
\phantomsection
\label{navigation:IV.1.1.7.1}
\addcontentsline{toc}{subsubsection}{IV.1.1.7.1 [Topic] CarePlace}

\label{element:779f2295b0b648dda9dbd21c4ade3ba0}
\index{[Topic] CarePlace (IV.1.1.7.1)}
{\small DataObject / Topic · Established Package1 conceptual information meaning; residual qualifications in selected section remain unchanged. · current}\par

\textbf{Description}
Information describing the characteristics and intended healthcare function of a place in which care may be provided, independently of its transient operational state.



\begin{itemize}


\item \textbf{Semantic Classification}: \texttt{\seqsplit{Definition}}

\item \textbf{Definition}: The static engineering, physical infrastructure, and architectural specification of a bed bay or clinical treatment point.

\item \textbf{Key Characteristics}: Physical capacity, isolation classification (e.g. \emph{Negative Pressure}, \emph{Positive Pressure}, \emph{Standard}), medical gas outlet provisions (e.g. \emph{Oxygen}, \emph{Suction}, \emph{Medical Air}), emergency power backup, and fixed telemetry wiring.

\item \textbf{Architectural Boundary}: Captures what the physical care-place \textbf{is} built to support, independent of its transient operational occupancy.

\end{itemize}


Source standing: Established Package1 conceptual information meaning; residual qualifications in selected section remain unchanged.


{\footnotesize \href{../../../../continuum/application/data-object/topic/779f2295b0b648dda9dbd21c4ade3ba0.md}{Canonical record: 779f2295b0b648dda9dbd21c4ade3ba0}; full qualifications and provenance.}\par
\subsubsection*{IV.1.1.7.2 [Topic] Geolocation}
\phantomsection
\label{navigation:IV.1.1.7.2}
\addcontentsline{toc}{subsubsection}{IV.1.1.7.2 [Topic] Geolocation}

\label{element:771b2e814946401bbe6dd6b4fdd3fb42}
\index{[Topic] Geolocation (IV.1.1.7.2)}
{\small DataObject / Topic · Established Package1 conceptual information meaning; residual qualifications in selected section remain unchanged. · current}\par

\textbf{Description}
Information describing the geographical position or address of a location, supporting its physical identification and location.



\begin{itemize}


\item \textbf{Semantic Classification}: \texttt{\seqsplit{Assertion / Spatial Coordinate}}

\item \textbf{Definition}: The physical address, geospatial coordinates (latitude/longitude), building elevation, and GIS geofence defining the physical location on Earth.

\end{itemize}


Source standing: Established Package1 conceptual information meaning; residual qualifications in selected section remain unchanged.


{\footnotesize \href{../../../../continuum/application/data-object/topic/771b2e814946401bbe6dd6b4fdd3fb42.md}{Canonical record: 771b2e814946401bbe6dd6b4fdd3fb42}; full qualifications and provenance.}\par
\subsubsection*{IV.1.1.7.3 [Topic] Location}
\phantomsection
\label{navigation:IV.1.1.7.3}
\addcontentsline{toc}{subsubsection}{IV.1.1.7.3 [Topic] Location}

\label{element:9ee4df4d76864bb1ac90fe1077c56c6f}
\index{[Topic] Location (IV.1.1.7.3)}
{\small DataObject / Topic · Established Package1 conceptual information meaning; residual qualifications in selected section remain unchanged. · current}\par

\textbf{Description}
Information describing an identifiable physical, geographical or functional location relevant to healthcare delivery or administration.



\begin{itemize}


\item \textbf{Semantic Classification}: \texttt{\seqsplit{Entity}}

\item \textbf{Definition}: An identifiable physical, geospatial, or functional place where healthcare is delivered, managed, or supported.

\item \textbf{Architectural Scope}: Represents physical sites (e.g. \emph{Main Hospital Campus}), buildings, wards, rooms, and specific bed bays, as well as mobile care units (e.g. \emph{Mobile Outreach Clinic}).

\end{itemize}


Source standing: Established Package1 conceptual information meaning; residual qualifications in selected section remain unchanged.


{\footnotesize \href{../../../../continuum/application/data-object/topic/9ee4df4d76864bb1ac90fe1077c56c6f.md}{Canonical record: 9ee4df4d76864bb1ac90fe1077c56c6f}; full qualifications and provenance.}\par
\subsubsection*{IV.1.1.7.4 [Topic] LocationClassification}
\phantomsection
\label{navigation:IV.1.1.7.4}
\addcontentsline{toc}{subsubsection}{IV.1.1.7.4 [Topic] LocationClassification}

\label{element:3896c910c28646098f0707efe88169ee}
\index{[Topic] LocationClassification (IV.1.1.7.4)}
{\small DataObject / Topic · Established Package1 conceptual information meaning; residual qualifications in selected section remain unchanged. · current}\par

\textbf{Description}
Information describing the type, function or other governed classification applicable to a location.



\begin{itemize}


\item \textbf{Semantic Classification}: \texttt{\seqsplit{Assertion / Classification}}

\item \textbf{Definition}: The structural and functional categorisation of the location (e.g. \emph{Campus}, \emph{Building}, \emph{Floor}, \emph{Ward / Wing}, \emph{Patient Room}, \emph{Care-Place / Bed Bay}, \emph{Operating Theatre}, \emph{Consulting Suite}, \emph{Mobile Unit}).

\end{itemize}


Source standing: Established Package1 conceptual information meaning; residual qualifications in selected section remain unchanged.


{\footnotesize \href{../../../../continuum/application/data-object/topic/3896c910c28646098f0707efe88169ee.md}{Canonical record: 3896c910c28646098f0707efe88169ee}; full qualifications and provenance.}\par
\subsubsection*{IV.1.1.7.5 [Topic] LocationIdentity}
\phantomsection
\label{navigation:IV.1.1.7.5}
\addcontentsline{toc}{subsubsection}{IV.1.1.7.5 [Topic] LocationIdentity}

\label{element:5f8299e31ee5410583007473c75d578f}
\index{[Topic] LocationIdentity (IV.1.1.7.5)}
{\small DataObject / Topic · Established Package1 conceptual information meaning; residual qualifications in selected section remain unchanged. · current}\par

\textbf{Description}
Information used to identify, distinguish and correlate a location across organisations, authorities, systems and healthcare contexts.



\begin{itemize}


\item \textbf{Semantic Classification}: \texttt{\seqsplit{Assertion / Identity Bundle}}

\item \textbf{Definition}: The governed set of enterprise and local identifiers that uniquely distinguish the location within an enterprise or regional jurisdiction.

\item \textbf{Key Semantic Decomposition}:

\item \texttt{\seqsplit{Identifier Value}}: The location code or name (e.g. \emph{CAMPUS-NORTH}, \emph{WARD-3B}, \emph{BED-304-A}).

\item \texttt{\seqsplit{Identifier Authority / Namespace}}: The health enterprise, government department, or facilities management namespace.

\end{itemize}


Source standing: Established Package1 conceptual information meaning; residual qualifications in selected section remain unchanged.


{\footnotesize \href{../../../../continuum/application/data-object/topic/5f8299e31ee5410583007473c75d578f.md}{Canonical record: 5f8299e31ee5410583007473c75d578f}; full qualifications and provenance.}\par
\subsection*{IV.1.1.8 [Topic] Scheduling}
\phantomsection
\label{navigation:IV.1.1.8}
\addcontentsline{toc}{subsection}{IV.1.1.8 [Topic] Scheduling}

\label{element:1225bfab0a6c4312875091a027f24275}
\index{[Topic] Scheduling (IV.1.1.8)}
{\small DataObject / Topic · Accepted bounded reconciliation · current}\par

\textbf{Description}
Information concerning the scheduling and administrative context of healthcare interactions, including appointments, scheduling changes, synchronisation state, attendance-related information and consolidated schedule perspectives.


{\footnotesize \href{../../../../continuum/application/data-object/topic/1225bfab0a6c4312875091a027f24275.md}{Canonical record: 1225bfab0a6c4312875091a027f24275}; full qualifications and provenance.}\par
\subsubsection*{IV.1.1.8.1 [Topic] Appointment}
\phantomsection
\label{navigation:IV.1.1.8.1}
\addcontentsline{toc}{subsubsection}{IV.1.1.8.1 [Topic] Appointment}

\label{element:9badf0d360b64215bc60e74eec954018}
\index{[Topic] Appointment (IV.1.1.8.1)}
{\small DataObject / Topic · Explicitly APPROVED G2 decision only; G2 remains candidate/open and all residual qualifications remain unresolved. · current}\par

""Description**
Information describing a planned healthcare interaction, including its subject, participants, intended service, time, place and scheduling context.



\textbf{Status: RESOLVED — APPROVED, with explicit unresolved cardinality/qualification details.}



\begin{quote}


\textbf{An Appointment may provide the prospective scheduling context for an Encounter. Where an Encounter occurs without a pre-existing Appointment, such as an emergency or walk-in interaction, an Appointment may be created retrospectively to represent the scheduling/administrative context associated with that Encounter.}



\end{quote}


Appointment and Encounter are distinct Information Concepts. Their relationship is principally scheduling context, not mandatory lifecycle progression.



\begin{quote}\small\ttfamily
Appointment ≠ Encounter\\
Appointment existence does not imply Encounter occurred\\
Encounter does not prove prospective Appointment existed\\
Retrospective Appointment establishment is not evidence of prospective scheduling\\

\end{quote}


\begin{itemize}


\item No Appointment → Encounter cardinality is established. Cancellation, rescheduling, recurring appointments, group appointments and other scheduling contexts remain unreviewed for multiplicity.

\item Prospective and retrospective context do not create separate Appointment Concepts. Scheduling-origin/context qualification is identified as a property/qualification for later review, without a defined representation or mandatory property structure.

\item The decision does not assign the originator/authority for retrospective creation or transfer external booking authority to Harmonia.

\end{itemize}


\textbf{Review consequence:} G2-Q02 reflects approved prospective/retrospective meaning and retains only unreviewed qualification/classification/authority/cardinality details. G2-Q03 preserves actual care-interaction meaning while retaining administrative tracking/Assembly questions.



Source standing: Explicitly APPROVED G2 decision only; G2 remains candidate/open and all residual qualifications remain unresolved.


{\footnotesize \href{../../../../continuum/application/data-object/topic/9badf0d360b64215bc60e74eec954018.md}{Canonical record: 9badf0d360b64215bc60e74eec954018}; full qualifications and provenance.}\par
\subsubsection*{IV.1.1.8.2 [Topic] PatientSchedule}
\phantomsection
\label{navigation:IV.1.1.8.2}
\addcontentsline{toc}{subsubsection}{IV.1.1.8.2 [Topic] PatientSchedule}

\label{element:c109f726f6ae4ee1b7fde1c012cfae74}
\index{[Topic] PatientSchedule (IV.1.1.8.2)}
{\small DataObject / Topic · Accepted bounded reconciliation · current}\par

\textbf{Schedule}
Information presenting the scheduled healthcare interactions relevant to an individual, providing a consolidated patient-oriented perspective across applicable appointments and scheduling sources.


{\footnotesize \href{../../../../continuum/application/data-object/topic/c109f726f6ae4ee1b7fde1c012cfae74.md}{Canonical record: c109f726f6ae4ee1b7fde1c012cfae74}; full qualifications and provenance.}\par
\subsubsection*{IV.1.1.8.3 [Topic] ScheduleChange}
\phantomsection
\label{navigation:IV.1.1.8.3}
\addcontentsline{toc}{subsubsection}{IV.1.1.8.3 [Topic] ScheduleChange}

\label{element:0e98b52f18274a81ad3e3b9ff9cca747}
\index{[Topic] ScheduleChange (IV.1.1.8.3)}
{\small DataObject / Topic · Accepted bounded reconciliation · current}\par

\textbf{Description}
Information describing a significant change to scheduled healthcare activity, including creation, rescheduling, cancellation or other alteration affecting the planned interaction.


{\footnotesize \href{../../../../continuum/application/data-object/topic/0e98b52f18274a81ad3e3b9ff9cca747.md}{Canonical record: 0e98b52f18274a81ad3e3b9ff9cca747}; full qualifications and provenance.}\par
\subsubsection*{IV.1.1.8.4 [Topic] ScheduleState}
\phantomsection
\label{navigation:IV.1.1.8.4}
\addcontentsline{toc}{subsubsection}{IV.1.1.8.4 [Topic] ScheduleState}

\label{element:6c4e9a91baf7460387d298366c5c6d0a}
\index{[Topic] ScheduleState (IV.1.1.8.4)}
{\small DataObject / Topic · Accepted bounded reconciliation · current}\par

\textbf{Description}
Information describing the synchronisation and current known state of scheduling information across participating sources or contexts.


{\footnotesize \href{../../../../continuum/application/data-object/topic/6c4e9a91baf7460387d298366c5c6d0a.md}{Canonical record: 6c4e9a91baf7460387d298366c5c6d0a}; full qualifications and provenance.}\par
\subsection*{IV.1.1.9 [Topic] Service}
\phantomsection
\label{navigation:IV.1.1.9}
\addcontentsline{toc}{subsection}{IV.1.1.9 [Topic] Service}

\label{element:0b990e0c64684bb8a45643e48eb4470a}
\index{[Topic] Service (IV.1.1.9)}
{\small DataObject / Topic · Accepted bounded reconciliation · current}\par

\textbf{Description}
Information describing healthcare services from definition and contextualised availability through delivery, outcome and assurance, including the relationships between services, providers, locations and the contexts in which services are provided.


{\footnotesize \href{../../../../continuum/application/data-object/topic/0b990e0c64684bb8a45643e48eb4470a.md}{Canonical record: 0b990e0c64684bb8a45643e48eb4470a}; full qualifications and provenance.}\par
\subsubsection*{IV.1.1.9.1 [Topic] ServiceAssurance}
\phantomsection
\label{navigation:IV.1.1.9.1}
\addcontentsline{toc}{subsubsection}{IV.1.1.9.1 [Topic] ServiceAssurance}

\label{element:07a1f83c9b4e462eaef1595c2a50d560}
\index{[Topic] ServiceAssurance (IV.1.1.9.1)}
{\small DataObject / Topic · Established Package1 conceptual information meaning; residual qualifications in selected section remain unchanged. · current}\par

\textbf{Description}
Information describing the assurance context associated with a healthcare service, including the evidence or determination that its provision satisfies applicable expectations, requirements or obligations.



\begin{itemize}


\item \textbf{Semantic Classification}: \texttt{\seqsplit{Accountability / Audit Record}}

\item \textbf{Definition}: The formal retrospective quality evaluation, clinical audit, accreditation review, or funding/billing reconciliation assessed against delivered services and outcomes.

\item \textbf{Architectural Scope \& Ownership Assessment}: Evaluates single service deliveries or aggregated cohorts across an episode of care. Because Harmonia's core operational capabilities do not claim R1 architectural ownership for statutory health insurance billing reconciliation or external accreditation governance, \texttt{\seqsplit{AssuredHealthcareService}} is designated as a \textbf{Reference / Contextual Concept} where Harmonia consumes or supports assurance interfaces without owning the external statutory authority.

\end{itemize}


Source standing: Established Package1 conceptual information meaning; residual qualifications in selected section remain unchanged.


{\footnotesize \href{../../../../continuum/application/data-object/topic/07a1f83c9b4e462eaef1595c2a50d560.md}{Canonical record: 07a1f83c9b4e462eaef1595c2a50d560}; full qualifications and provenance.}\par
\subsubsection*{IV.1.1.9.2 [Topic] ServiceDelivery}
\phantomsection
\label{navigation:IV.1.1.9.2}
\addcontentsline{toc}{subsubsection}{IV.1.1.9.2 [Topic] ServiceDelivery}

\label{element:4627ee33845744f2ae2e10c84e313755}
\index{[Topic] ServiceDelivery (IV.1.1.9.2)}
{\small DataObject / Topic · Established Package1 conceptual information meaning; residual qualifications in selected section remain unchanged. · current}\par

\textbf{Description}
Information describing the actual provision of a healthcare service within a particular subject, provider, place and delivery context.



\begin{itemize}


\item \textbf{Semantic Classification}: \texttt{\seqsplit{Activity / Fulfilment}}

\item \textbf{Definition}: An Information Concept referring to the real-world execution or performance of a healthcare service for a specific target subject (or facility target) at a concrete point in time. It represents information about the activity rather than becoming the activity or implying that Harmonia performs it.

\item \textbf{Key Characteristics}: Actual start/end timestamps, performing practitioners, delivery location, executing devices/equipment, and execution status progression.

\item \textbf{Healthcare Subject Invariant}: While clinical service deliveries typically associate with a \texttt{\seqsplit{Healthcare Subject Context}}, broader service semantics (e.g., \emph{Facility Sanitisation}, \emph{Water Quality Testing}, \emph{Equipment Calibration}) do not universally require a Healthcare Subject.

\end{itemize}


Source standing: Established Package1 conceptual information meaning; residual qualifications in selected section remain unchanged.


{\footnotesize \href{../../../../continuum/application/data-object/topic/4627ee33845744f2ae2e10c84e313755.md}{Canonical record: 4627ee33845744f2ae2e10c84e313755}; full qualifications and provenance.}\par
\subsubsection*{IV.1.1.9.3 [Topic] ServiceOffering}
\phantomsection
\label{navigation:IV.1.1.9.3}
\addcontentsline{toc}{subsubsection}{IV.1.1.9.3 [Topic] ServiceOffering}

\label{element:441b0891822f4add834ea930c3fe58b7}
\index{[Topic] ServiceOffering (IV.1.1.9.3)}
{\small DataObject / Topic · Established Package1 conceptual information meaning; residual qualifications in selected section remain unchanged. · current}\par

\textbf{Description}
Information describing a healthcare service that an organisation or provider makes available, including its nature, characteristics and conditions of availability.



\begin{itemize}


\item \textbf{Semantic Classification}: \texttt{\seqsplit{Definition}}

\item \textbf{Definition}: The canonical service offering defined within a healthcare enterprise or jurisdictional catalogue (e.g. \emph{12-Lead Electrocardiography}, \emph{Acute Inpatient Haemodialysis}, \emph{Wound Debridement Protocol}, \emph{Ward Sanitisation Protocol}).

\item \textbf{Key Characteristics}: Service classification, specialty taxonomy, clinical indications, contraindications, required workforce competencies, supported delivery modalities, and permitted composition structures.

\item \textbf{Recursive Composition}: An \texttt{\seqsplit{OfferedHealthcareService}} may contain other \texttt{\seqsplit{OfferedHealthcareService}} definitions (e.g. \emph{Cardiology Service} contains \emph{Echocardiography} and \emph{Stress ECG}).

\end{itemize}


Source standing: Established Package1 conceptual information meaning; residual qualifications in selected section remain unchanged.


{\footnotesize \href{../../../../continuum/application/data-object/topic/441b0891822f4add834ea930c3fe58b7.md}{Canonical record: 441b0891822f4add834ea930c3fe58b7}; full qualifications and provenance.}\par
\subsubsection*{IV.1.1.9.4 [Topic] ServiceOutcome}
\phantomsection
\label{navigation:IV.1.1.9.4}
\addcontentsline{toc}{subsubsection}{IV.1.1.9.4 [Topic] ServiceOutcome}

\label{element:4fb589629e7f405db44d6a87a2b5e36b}
\index{[Topic] ServiceOutcome (IV.1.1.9.4)}
{\small DataObject / Topic · Established Package1 conceptual information meaning; residual qualifications in selected section remain unchanged. · current}\par

\textbf{Description}
Information describing the outcome resulting from the delivery of a healthcare service, including the relevant disposition or consequences of that delivery.



\begin{itemize}


\item \textbf{Semantic Classification}: \texttt{\seqsplit{Outcome}}

\item \textbf{Definition}: An Information Concept concerning a result, finding, change, consequence or administrative artefact associated with a service delivery. These information kinds remain distinct; representing them does not confer responsibility for their production, clinical interpretation or originating authority.

\item \textbf{Key Scope}: Spans across:

\item \textbf{Clinical Outcomes}: Diagnoses, physiological improvements, symptom resolutions;

\item \textbf{Diagnostic Outcomes}: Pathology findings, radiology reports, ECG interpretations;

\item \textbf{Therapeutic Outcomes}: Medication administration confirmations, surgical revisions;

\item \textbf{Operational Outcomes}: Sanitised bed bay certification, repaired asset state;

\item \textbf{Administrative Outcomes}: Completed intake assessments, eligibility confirmations.

\item \textbf{Preserved uncertainty}: Exact information-kind semantics and a universal Delivery/Outcome ownership model are not established. No single representation model, lifecycle or performing Capability is inferred.

\end{itemize}


Source standing: Established Package1 conceptual information meaning; residual qualifications in selected section remain unchanged.


{\footnotesize \href{../../../../continuum/application/data-object/topic/4fb589629e7f405db44d6a87a2b5e36b.md}{Canonical record: 4fb589629e7f405db44d6a87a2b5e36b}; full qualifications and provenance.}\par
\subsubsection*{IV.1.1.9.5 [Topic] ServiceProvision}
\phantomsection
\label{navigation:IV.1.1.9.5}
\addcontentsline{toc}{subsubsection}{IV.1.1.9.5 [Topic] ServiceProvision}

\label{element:8d41e9fb6a164deb93afb5fc4c671e10}
\index{[Topic] ServiceProvision (IV.1.1.9.5)}
{\small DataObject / Topic · Established Package1 conceptual information meaning; residual qualifications in selected section remain unchanged. · current}\par

\textbf{Description}
Information describing an offered service contextualised for actual provision, including the provider, location, applicable conditions and other bindings necessary for the service to be delivered.



\begin{itemize}


\item \textbf{Semantic Classification}: \texttt{\seqsplit{Contextual Binding}}

\item \textbf{Definition}: The contextual binding of an \texttt{\seqsplit{OfferedHealthcareService}} to a specific providing organisation, delivery location, participating practitioner roles, operational availability schedules, and eligibility rule sets.

\item \textbf{Key Characteristics}: Bound service definition, performing organisation, delivery site, participating practitioner role bindings, intake schedule, and catchment criteria.

\end{itemize}


Source standing: Established Package1 conceptual information meaning; residual qualifications in selected section remain unchanged.


{\footnotesize \href{../../../../continuum/application/data-object/topic/8d41e9fb6a164deb93afb5fc4c671e10.md}{Canonical record: 8d41e9fb6a164deb93afb5fc4c671e10}; full qualifications and provenance.}\par
\subsection*{IV.1.1.10 [Topic] TaskCompletion}
\phantomsection
\label{navigation:IV.1.1.10}
\addcontentsline{toc}{subsection}{IV.1.1.10 [Topic] TaskCompletion}

\label{element:8a9351cd170940d0b36f8d56584f15ec}
\index{[Topic] TaskCompletion (IV.1.1.10)}
{\small DataObject / Topic · Accepted bounded reconciliation · current}\par

\textbf{Description} 
Information describing actionable work from reusable definition and instantiated task through undertaking and governed resolution, including execution participation, completion information, outcomes and resulting outputs.


{\footnotesize \href{../../../../continuum/application/data-object/topic/8a9351cd170940d0b36f8d56584f15ec.md}{Canonical record: 8a9351cd170940d0b36f8d56584f15ec}; full qualifications and provenance.}\par
\subsubsection*{IV.1.1.10.1 [Topic] ExecutionConcurrency}
\phantomsection
\label{navigation:IV.1.1.10.1}
\addcontentsline{toc}{subsubsection}{IV.1.1.10.1 [Topic] ExecutionConcurrency}

\label{element:0781dd44748d42338c18bf01993af656}
\index{[Topic] ExecutionConcurrency (IV.1.1.10.1)}
{\small DataObject / Topic · Accepted bounded reconciliation · current}\par

\textbf{Description}
Information governing or describing concurrency applicable to the execution of work associated with a Task.


{\footnotesize \href{../../../../continuum/application/data-object/topic/0781dd44748d42338c18bf01993af656.md}{Canonical record: 0781dd44748d42338c18bf01993af656}; full qualifications and provenance.}\par
\subsubsection*{IV.1.1.10.2 [Topic] OutcomeConcurrency}
\phantomsection
\label{navigation:IV.1.1.10.2}
\addcontentsline{toc}{subsubsection}{IV.1.1.10.2 [Topic] OutcomeConcurrency}

\label{element:72c505782347448d8f89d74ef76359af}
\index{[Topic] OutcomeConcurrency (IV.1.1.10.2)}
{\small DataObject / Topic · Accepted bounded reconciliation · current}\par

\textbf{Description}
Information describing the concurrency rules applicable when determining or establishing the outcome of a Task.


{\footnotesize \href{../../../../continuum/application/data-object/topic/72c505782347448d8f89d74ef76359af.md}{Canonical record: 72c505782347448d8f89d74ef76359af}; full qualifications and provenance.}\par
\subsubsection*{IV.1.1.10.3 [Topic] OutputRules}
\phantomsection
\label{navigation:IV.1.1.10.3}
\addcontentsline{toc}{subsubsection}{IV.1.1.10.3 [Topic] OutputRules}

\label{element:a035fe1913db4d2589a2e4edde8db737}
\index{[Topic] OutputRules (IV.1.1.10.3)}
{\small DataObject / Topic · Accepted bounded reconciliation · current}\par

\textbf{Description}
Information describing the rules governing the outputs established as part of Task resolution.


{\footnotesize \href{../../../../continuum/application/data-object/topic/a035fe1913db4d2589a2e4edde8db737.md}{Canonical record: a035fe1913db4d2589a2e4edde8db737}; full qualifications and provenance.}\par
\subsubsection*{IV.1.1.10.4 [Topic] Task}
\phantomsection
\label{navigation:IV.1.1.10.4}
\addcontentsline{toc}{subsubsection}{IV.1.1.10.4 [Topic] Task}

\label{element:c43139289cb649b4836521a96a51604c}
\index{[Topic] Task (IV.1.1.10.4)}
{\small DataObject / Topic · Bounded adjudicated conceptual Task meaning; section 5 reservations, including runtime TaskOutcome identity/lifecycle, remain unresolved. · current}\par

\textbf{Description}
Information describing an identifiable instance of work to be completed, including its purpose, context, state and applicable completion expectations.



For example, the archetype \textbf{"Climb a ladder"} describes a kind of work;
\textbf{"ladder-climb-0023"} identifies a particular \texttt{\seqsplit{ActionableTask}} of that kind.
These are conceptual examples, not an identifier scheme.



An \texttt{\seqsplit{ActionableTask}} SHALL NOT mean the reusable definition or the execution /
undertaking of the work. Its contextual work meaning remains distinct from
information about particular undertakings.



An \texttt{\seqsplit{ActionableTask}} explicitly governs execution concurrency, which
undertakings participate in resolving its outcome, and how their outputs
establish its \texttt{\seqsplit{Task.Output(s)}}. The following three controls are independent
semantic properties, not an implementation structure or execution mechanism.



\par\medskip\textbf{2.2.1 ActionableTask.ExecutionConcurrency}\par


\begin{quote}\small\ttfamily
ActionableTask.ExecutionConcurrency = [1..x]\\

\end{quote}


\texttt{\seqsplit{ExecutionConcurrency}} establishes how many \texttt{\seqsplit{FulfillmentTask}} undertakings of
the \texttt{\seqsplit{ActionableTask}} may execute concurrently. It concerns execution
concurrency only; it does not determine which undertakings establish the
\texttt{\seqsplit{TaskOutcome}}. This permitted concurrency range does not require an
undertaking to exist or execute: an ActionableTask may still have zero
undertakings. No universal value for \texttt{\seqsplit{x}} or concurrency execution mechanism
is established here.



\par\medskip\textbf{2.2.2 ActionableTask.OutcomeConcurrency.Mode}\par


\begin{quote}\small\ttfamily
ActionableTask.OutcomeConcurrency.Mode =\\
    FirstToFinish | AssignedToFinish | Aggregate\\

\end{quote}


This property determines \textbf{which FulfillmentTask undertaking or undertakings
participate in resolving the ActionableTask outcome}.



\begin{longtable}{>{\raggedright\arraybackslash}p{\dimexpr\linewidth/2-2\tabcolsep\relax}>{\raggedright\arraybackslash}p{\dimexpr\linewidth/2-2\tabcolsep\relax}}
Mode & Participation semantics \\
\textbf{FirstToFinish} & The applicable first finishing/qualifying \texttt{\seqsplit{FulfillmentTask}} establishes participation in resolution of the outcome. This defines neither detailed terminal-state qualification nor cancellation or execution mechanics for other undertakings. \\
\textbf{AssignedToFinish} & A particular assigned \texttt{\seqsplit{FulfillmentTask}} has responsibility for establishing the outcome. Other undertakings may exist or execute concurrently; their existence or completion does not independently establish the ActionableTask outcome. The assignment mechanism is not defined here. \\
\textbf{Aggregate} & Multiple applicable \texttt{\seqsplit{FulfillmentTask}} undertakings participate in resolving the ActionableTask outcome. This selects the participating undertakings; it does not require recursive TaskOutcome composition or determine how their outputs form the resulting output. \\
\end{longtable}


\par\medskip\textbf{2.2.3 ActionableTask.Outcome.OutputRules.Mode}\par


\begin{quote}\small\ttfamily
ActionableTask.Outcome.OutputRules.Mode = Direct | Collection | Derived\\

\end{quote}


This independent property determines \textbf{how the outputs of the participating
FulfillmentTasks establish the ActionableTask's Task.Output(s)}.



\begin{longtable}{>{\raggedright\arraybackslash}p{\dimexpr\linewidth/3-2\tabcolsep\relax}>{\raggedright\arraybackslash}p{\dimexpr\linewidth/3-2\tabcolsep\relax}>{\raggedright\arraybackslash}p{\dimexpr\linewidth/3-2\tabcolsep\relax}}
Mode & Output semantics & Conceptual path \\
\textbf{Direct} & The applicable participating \texttt{\seqsplit{FulfillmentTask}} output is used directly as the ActionableTask output. & \texttt{\seqsplit{FulfillmentTask.Output}} → \texttt{\seqsplit{TaskOutcome}} → \texttt{\seqsplit{Task.Output}}. \\
\textbf{Collection} & The applicable participating \texttt{\seqsplit{FulfillmentTask}} outputs collectively form the ActionableTask output, preserving them as the relevant collection/set. Collection does not itself imply derivation of new domain meaning. & \texttt{\seqsplit{FulfillmentTask.Output [0..*]}} → Collection → \texttt{\seqsplit{TaskOutcome}} → \texttt{\seqsplit{Task.Output(s)}}. \\
\textbf{Derived} & A new ActionableTask output is derived from the applicable participating \texttt{\seqsplit{FulfillmentTask}} outputs according to explicitly defined rules. & \texttt{\seqsplit{FulfillmentTask.Output(s)}} → Explicit Rules → \texttt{\seqsplit{TaskOutcome}} → \texttt{\seqsplit{Task.Output(s)}}. \\
\end{longtable}


\texttt{\seqsplit{Derived}} is an information semantic. It prescribes no implementation
mechanism for executing the explicit rules. Harmonia need not infer the
encapsulated domain meaning of participating outputs in order to apply those
rules; the
\hyperref[element:10284034c8534b639284386b064f2c54]{observable-information/domain-meaning boundary}
continues to apply. Developer-defined Ergo logic and optional runtime AI
retain their existing boundaries; neither is prescribed as the means of
derivation.



\par\medskip\textbf{2.2.4 Orthogonality of the Three Controls}\par


\begin{longtable}{>{\raggedright\arraybackslash}p{\dimexpr\linewidth/2-2\tabcolsep\relax}>{\raggedright\arraybackslash}p{\dimexpr\linewidth/2-2\tabcolsep\relax}}
ActionableTask concern & Question answered \\
\textbf{ExecutionConcurrency} & How many FulfillmentTasks may execute concurrently? \\
\textbf{OutcomeConcurrency.Mode} & Which FulfillmentTask undertaking(s) participate in resolving the ActionableTask outcome? \\
\textbf{Outcome.OutputRules.Mode} & How are the participating FulfillmentTask outputs used to establish the ActionableTask Task.Output(s)? \\
\end{longtable}


These controls SHALL remain orthogonal and SHALL NOT be collapsed into one
lifecycle, state or concurrency mechanism. Multiple undertakings may execute
concurrently while only one participates in the TaskOutcome. Multiple may
participate while their outputs are collected or used to derive another
output. No universal mode combination, default or execution algorithm is
inferred from these examples.



Source standing: Bounded adjudicated conceptual Task meaning; section 5 reservations, including runtime TaskOutcome identity/lifecycle, remain unresolved.


{\footnotesize \href{../../../../continuum/application/data-object/topic/c43139289cb649b4836521a96a51604c.md}{Canonical record: c43139289cb649b4836521a96a51604c}; full qualifications and provenance.}\par
\subsubsection*{IV.1.1.10.5 [Topic] TaskDefinition}
\phantomsection
\label{navigation:IV.1.1.10.5}
\addcontentsline{toc}{subsubsection}{IV.1.1.10.5 [Topic] TaskDefinition}

\label{element:1d6d4c3a2df440be94262855fa6564eb}
\index{[Topic] TaskDefinition (IV.1.1.10.5)}
{\small DataObject / Topic · Bounded adjudicated conceptual Task meaning; section 5 reservations, including runtime TaskOutcome identity/lifecycle, remain unresolved. · current}\par

\textbf{Description}
Information defining a reusable form of actionable work, including the characteristics, expectations and constraints applicable when instances of that work are created.



\textbf{ActionableTaskArchetype} is the reusable definition/archetype for a kind of
actionable work. It describes the meaning and constraints of work that may
subsequently be instantiated. It SHALL NOT be treated as an instance of work.



Its semantics may include, where applicable:



\begin{itemize}


\item the expected nature of the work;

\item permissible parameters and constraints;

\item prerequisites and performer requirements;

\item completion or satisfaction expectations; and

\item expectations concerning information/content consumed, manipulated or
  produced by the work.

\end{itemize}


An archetype may establish expectations concerning encapsulated content
without defining that content's structural representation or requiring
Harmonia to comprehend its domain semantics. The
\hyperref[element:10284034c8534b639284386b064f2c54]{observable-information boundary}
applies.



Source standing: Bounded adjudicated conceptual Task meaning; section 5 reservations, including runtime TaskOutcome identity/lifecycle, remain unresolved.


{\footnotesize \href{../../../../continuum/application/data-object/topic/1d6d4c3a2df440be94262855fa6564eb.md}{Canonical record: 1d6d4c3a2df440be94262855fa6564eb}; full qualifications and provenance.}\par
\subsubsection*{IV.1.1.10.6 [Topic] TaskFulfillment}
\phantomsection
\label{navigation:IV.1.1.10.6}
\addcontentsline{toc}{subsubsection}{IV.1.1.10.6 [Topic] TaskFulfillment}

\label{element:30f4deffcd694f07bf54b2b192695d4f}
\index{[Topic] TaskFulfillment (IV.1.1.10.6)}
{\small DataObject / Topic · Bounded adjudicated conceptual Task meaning; section 5 reservations, including runtime TaskOutcome identity/lifecycle, remain unresolved. · current}\par

\textbf{Description}
Information describing an identifiable undertaking performed toward completion of a Task, allowing the work performed and its relationship to the originating Task to be represented.



\textbf{FulfillmentTask} represents an identifiable undertaking of an
\texttt{\seqsplit{ActionableTask}}: doing or attempting to do that particular work. It carries
the execution-related information necessary to represent the undertaking,
including, where applicable:



\begin{itemize}


\item performer / executing actor and execution context;

\item commencement and timing;

\item progression and execution state;

\item completion, failure, suspension or other applicable execution information;
  and

\item other execution metadata/scaffolding needed to represent the undertaking.

\end{itemize}


It has its own execution state and may produce output. It does \textbf{not
independently own TaskOutcome}. One or more undertakings may contribute to
resolution of the parent ActionableTask according to that ActionableTask's
explicitly defined outcome policy.



An \texttt{\seqsplit{ActionableTask}} may have \textbf{zero, one or multiple \texttt{\seqsplit{FulfillmentTasks}}}.
Multiple undertakings MAY execute concurrently where the ActionableTask's
\texttt{\seqsplit{ExecutionConcurrency}} permits them. Execution state SHALL NOT be collapsed into
\texttt{\seqsplit{ActionableTask}} on the assumption that each work instance has exactly one
execution.



Completion of one \texttt{\seqsplit{FulfillmentTask}} does \textbf{not inherently establish
satisfaction or completion of its parent \texttt{\seqsplit{ActionableTask}}}. Satisfaction
depends on the definition and context of the work, with outcome resolution
governed by the ActionableTask's explicit outcome policy. No universal
execution state machine or satisfaction rule is established here. Information
representing an undertaking remains distinct from the undertaking itself.



Source standing: Bounded adjudicated conceptual Task meaning; section 5 reservations, including runtime TaskOutcome identity/lifecycle, remain unresolved.


{\footnotesize \href{../../../../continuum/application/data-object/topic/30f4deffcd694f07bf54b2b192695d4f.md}{Canonical record: 30f4deffcd694f07bf54b2b192695d4f}; full qualifications and provenance.}\par
\subsubsection*{IV.1.1.10.7 [Topic] TaskOutcome}
\phantomsection
\label{navigation:IV.1.1.10.7}
\addcontentsline{toc}{subsubsection}{IV.1.1.10.7 [Topic] TaskOutcome}

\label{element:08d8ee19a30f418f9d5bc497ced0d1ef}
\index{[Topic] TaskOutcome (IV.1.1.10.7)}
{\small DataObject / Topic · Bounded adjudicated conceptual Task meaning; section 5 reservations, including runtime TaskOutcome identity/lifecycle, remain unresolved. · current}\par

\textbf{Description}
Information describing the governed resolution of a Task, including its completion disposition, resulting information and applicable outputs.



\begin{quote}


\textbf{TaskOutcome is the governed resolution of an ActionableTask according to
that ActionableTask's explicitly defined outcome policy.}



\end{quote}


TaskOutcome belongs to and resolves the ActionableTask. It does not
independently belong to a FulfillmentTask. Conceptually it:



\begin{itemize}


\item \textbf{references FulfillmentTask [0..*]}: identifies the undertaking or
  undertakings that participated in establishing the outcome, as determined
  by \texttt{\seqsplit{OutcomeConcurrency.Mode}};

\item \textbf{contains Task Completion Metadata}: Harmonia-governed information
  concerning resolution/completion of the ActionableTask; no complete
  metadata schema is defined here; and

\item \textbf{contains / represents Task.Output(s)}: the resulting ActionableTask
  output established from the participating undertakings' outputs according
  to \texttt{\seqsplit{Outcome.OutputRules.Mode}}.

\end{itemize}


The encapsulated domain content of a \texttt{\seqsplit{Task.Output}} may remain semantically
opaque to Harmonia. Neither output information nor TaskOutcome becomes the
physical, clinical or other phenomenon represented. No detailed Task.Output
payload structure is established here.



\textbf{Recursive TaskOutcome containment is excluded from this model.} Multiple
execution results are represented through references from the single
governed TaskOutcome to the relevant participating FulfillmentTask
undertakings, rather than through recursive outcome aggregation hierarchies.
Aggregate participation does not imply recursive composition. The remaining
questions in §5 remain unresolved.



Source standing: Bounded adjudicated conceptual Task meaning; section 5 reservations, including runtime TaskOutcome identity/lifecycle, remain unresolved.



Under \hyperref[element:8882573df7ce43c6bc0f66da0a676139]{AX-17}, the following
remaining Task / Work matters require subsequent architectural adjudication
or separately governed detailed definition:



\begin{itemize}


\item TaskOutcome identity rules;

\item detailed cardinalities beyond the established relationships, including
  whether an outcome is mandatory;

\item detailed TaskOutcome lifecycle/state semantics;

\item detailed Task Completion Metadata and Task.Output payload structures, and
  their representation mechanisms;

\item detailed finishing/qualification rules, assignment mechanisms and
  concurrency execution mechanics; and

\item task-specific explicit output rules and their downstream execution
  mechanisms.

\end{itemize}


Outcome ownership, participation modes, output-rule modes, their orthogonality
and exclusion of recursive TaskOutcome composition are adjudicated above;
they SHALL NOT be presented as still provisional or unresolved. This
refinement allocates no previously unestablished archetype-authoring or
outcome-authority owner.



\textbf{ReportedTask remains unresolved.} Its existing intersection with
Definition → Contextualisation → Fulfilment → Outcome → Accountability does
not reaffirm its previous definition or require a separate accountable report.
Its necessity and semantics may overlap accountability, assurance, audit,
historical representation and evidence, and require separate adjudication.
It is neither removed nor normalised into one of those concerns here. No
ReportedTask relationships, identity, immutability or lifecycle are settled.



The metamodel, reusable pattern and general lifecycle/containment guidance
SHALL NOT be used to fill these reserved questions by inference. The bounded
semantics above are established; the remaining Task / Work family and Domain04
are not declared complete.


{\footnotesize \href{../../../../continuum/application/data-object/topic/08d8ee19a30f418f9d5bc497ced0d1ef.md}{Canonical record: 08d8ee19a30f418f9d5bc497ced0d1ef}; full qualifications and provenance.}\par
\subsection*{IV.1.1.11 [Topic] Workforce}
\phantomsection
\label{navigation:IV.1.1.11}
\addcontentsline{toc}{subsection}{IV.1.1.11 [Topic] Workforce}

\label{element:9f689e261140436eb94295d63e2b3370}
\index{[Topic] Workforce (IV.1.1.11)}
{\small DataObject / Topic · Accepted bounded reconciliation · current}\par

\textbf{Description}: Information concerning healthcare practitioners and workforce participants, including professional identity, registration, qualifications, privileges, organisational affiliations, roles, and communication endpoints.


{\footnotesize \href{../../../../continuum/application/data-object/topic/9f689e261140436eb94295d63e2b3370.md}{Canonical record: 9f689e261140436eb94295d63e2b3370}; full qualifications and provenance.}\par
\subsubsection*{IV.1.1.11.1 [Topic] CommunicationAddress}
\phantomsection
\label{navigation:IV.1.1.11.1}
\addcontentsline{toc}{subsubsection}{IV.1.1.11.1 [Topic] CommunicationAddress}

\label{element:410be08eea0b4653a88f4e4235354cea}
\index{[Topic] CommunicationAddress (IV.1.1.11.1)}
{\small DataObject / Topic · Accepted bounded reconciliation · current}\par

\textbf{Description}
Information identifying an electronic communication destination associated with a practitioner or professional context, including the addressing and connection information required to communicate with that destination.


{\footnotesize \href{../../../../continuum/application/data-object/topic/410be08eea0b4653a88f4e4235354cea.md}{Canonical record: 410be08eea0b4653a88f4e4235354cea}; full qualifications and provenance.}\par
\subsubsection*{IV.1.1.11.2 [Topic] Practitioner}
\phantomsection
\label{navigation:IV.1.1.11.2}
\addcontentsline{toc}{subsubsection}{IV.1.1.11.2 [Topic] Practitioner}

\label{element:4e8f2bd581a5418889813cf7ee414750}
\index{[Topic] Practitioner (IV.1.1.11.2)}
{\small DataObject / Topic · Established Package1 conceptual information meaning; residual qualifications in selected section remain unchanged. · current}\par

\textbf{Definition}: An identifiable individual healthcare workforce or professional entity.



\textbf{Description}: Information describing an individual participating professionally in healthcare, independent of a particular organisational position, assignment or service context.



\textbf{Semantic Classification}: \texttt{\seqsplit{Entity}}



\textbf{Architectural Scope}: Represents the individual professional independently of statutory registration status. A \texttt{\seqsplit{Practitioner}} may be known to Harmonia prior to verification, when registration is pending or not applicable (e.g. unregistered healthcare workers, support staff), or when registration is externally authoritative.



\textbf{Boundary Rule}: \$\textbackslash{}text\{Practitioner\} \textbackslash{}neq \textbackslash{}text\{Professional Registration\}\$. Registration is a governed assertion concerning a Practitioner, not the definition of the Practitioner entity itself.



Source standing: Established Package1 conceptual information meaning; residual qualifications in selected section remain unchanged.


{\footnotesize \href{../../../../continuum/application/data-object/topic/4e8f2bd581a5418889813cf7ee414750.md}{Canonical record: 4e8f2bd581a5418889813cf7ee414750}; full qualifications and provenance.}\par
\subsubsection*{IV.1.1.11.3 [Topic] PractitionerRole}
\phantomsection
\label{navigation:IV.1.1.11.3}
\addcontentsline{toc}{subsubsection}{IV.1.1.11.3 [Topic] PractitionerRole}

\label{element:cb7ac98f62aa4599ba57263338cb8b3c}
\index{[Topic] PractitionerRole (IV.1.1.11.3)}
{\small DataObject / Topic · Established Package1 conceptual information meaning; residual qualifications in selected section remain unchanged. · current}\par

\textbf{Definition}: The contextual binding of a \texttt{\seqsplit{Practitioner}} to a specific clinical capacity, service, facility, or department within an organisation (e.g. \emph{Specialist Physician in Cardiology Department}, \emph{On-Call Registrar}).



\textbf{Description}: Information describing the contextual relationship between a practitioner and the organisation, service, location or role through which they participate in healthcare delivery.



\textbf{Semantic Classification}: \texttt{\seqsplit{Contextual Binding}}



\textbf{Key Characteristics}: Contextual title/role, participating organisation, participating healthcare service, physical practice locations, active schedule, and linked communication endpoints.



Source standing: Established Package1 conceptual information meaning; residual qualifications in selected section remain unchanged.


{\footnotesize \href{../../../../continuum/application/data-object/topic/cb7ac98f62aa4599ba57263338cb8b3c.md}{Canonical record: cb7ac98f62aa4599ba57263338cb8b3c}; full qualifications and provenance.}\par
\subsubsection*{IV.1.1.11.4 [Topic] Privilege}
\phantomsection
\label{navigation:IV.1.1.11.4}
\addcontentsline{toc}{subsubsection}{IV.1.1.11.4 [Topic] Privilege}

\label{element:75e194f6168945588ed30cc6c78d04cc}
\index{[Topic] Privilege (IV.1.1.11.4)}
{\small DataObject / Topic · Established Package1 conceptual information meaning; residual qualifications in selected section remain unchanged. · current}\par

\textbf{Definition}: An enterprise-approved scope of clinical practice, procedural authorization, admission entitlement, or prescribing privilege granted by a healthcare organisation.



\textbf{Description}: Information describing the activities, services or clinical practices a practitioner is authorised to undertake within a particular organisational or service context.



\textbf{Semantic Classification}: \texttt{\seqsplit{Assertion / Entitlement Directive}}



\textbf{Key Characteristics}: Granting organisation, authorized clinical scope / procedure codes, supervision requirements, effective validity period, and approving clinical governance authority.



Source standing: Established Package1 conceptual information meaning; residual qualifications in selected section remain unchanged.


{\footnotesize \href{../../../../continuum/application/data-object/topic/75e194f6168945588ed30cc6c78d04cc.md}{Canonical record: 75e194f6168945588ed30cc6c78d04cc}; full qualifications and provenance.}\par
\subsubsection*{IV.1.1.11.5 [Topic] ProfessionalIdentity}
\phantomsection
\label{navigation:IV.1.1.11.5}
\addcontentsline{toc}{subsubsection}{IV.1.1.11.5 [Topic] ProfessionalIdentity}

\label{element:2f2868f077fc47afbb541925a9801b2c}
\index{[Topic] ProfessionalIdentity (IV.1.1.11.5)}
{\small DataObject / Topic · Established Package1 conceptual information meaning; residual qualifications in selected section remain unchanged. · current}\par

\textbf{Definition}: The governed set of professional identifiers, credential assertions, and demographic traits establishing the professional identity of a \texttt{\seqsplit{Practitioner}}.



\textbf{Description}: Information establishing and correlating the professional identity of a practitioner across organisations, authorities and healthcare contexts.



\textbf{Semantic Classification}: \texttt{\seqsplit{Assertion / Identity Bundle}}
\textbf{Key Semantic Decomposition}:
  - \texttt{\seqsplit{Identifier Value}}: The token or registration number identifying the practitioner.
  - \texttt{\seqsplit{Identifier Authority / Namespace}}: The statutory body, national authority, or enterprise issuing the identifier.



Source standing: Established Package1 conceptual information meaning; residual qualifications in selected section remain unchanged.


{\footnotesize \href{../../../../continuum/application/data-object/topic/2f2868f077fc47afbb541925a9801b2c.md}{Canonical record: 2f2868f077fc47afbb541925a9801b2c}; full qualifications and provenance.}\par
\subsubsection*{IV.1.1.11.6 [Topic] ProfessionalProfile}
\phantomsection
\label{navigation:IV.1.1.11.6}
\addcontentsline{toc}{subsubsection}{IV.1.1.11.6 [Topic] ProfessionalProfile}

\label{element:703bb80c062743eeb7ed8dea3ffeaaa0}
\index{[Topic] ProfessionalProfile (IV.1.1.11.6)}
{\small DataObject / Topic · Established Package1 conceptual information meaning; residual qualifications in selected section remain unchanged. · current}\par

\textbf{Definition}: The verified set of educational qualifications, clinical competencies, certified subspecialties, and academic credentials possessed by a practitioner.



\textbf{Description}: Information describing the professional characteristics of a practitioner, including qualifications, competencies, specialties and other information relevant to professional practice.



\textbf{Semantic Classification}: \texttt{\seqsplit{Assertion / Profile Bundle}}



Source standing: Established Package1 conceptual information meaning; residual qualifications in selected section remain unchanged.


{\footnotesize \href{../../../../continuum/application/data-object/topic/703bb80c062743eeb7ed8dea3ffeaaa0.md}{Canonical record: 703bb80c062743eeb7ed8dea3ffeaaa0}; full qualifications and provenance.}\par
\subsubsection*{IV.1.1.11.7 [Topic] Registration}
\phantomsection
\label{navigation:IV.1.1.11.7}
\addcontentsline{toc}{subsubsection}{IV.1.1.11.7 [Topic] Registration}

\label{element:26c72794993044af98b5329b8b14ae89}
\index{[Topic] Registration (IV.1.1.11.7)}
{\small DataObject / Topic · Established Package1 conceptual information meaning; residual qualifications in selected section remain unchanged. · current}\par

\textbf{Definition}: A formal, externally governed assertion attesting that a practitioner is registered with a statutory registration board or regulatory authority.



\textbf{Description}: Information describing a practitioner's professional registration, including the registering authority, profession, registration status, scope and applicable period.



\textbf{Semantic Classification}: \texttt{\seqsplit{Assertion / Governed Relationship}}



\textbf{Key Characteristics}: Registration authority, registration category, endorsed specialties, condition/restriction status, and validity period.



\textbf{Authority Invariant}: Originates from an external regulatory authority. Harmonia records verification provenance without claiming originating registration authority.



Source standing: Established Package1 conceptual information meaning; residual qualifications in selected section remain unchanged.


{\footnotesize \href{../../../../continuum/application/data-object/topic/26c72794993044af98b5329b8b14ae89.md}{Canonical record: 26c72794993044af98b5329b8b14ae89}; full qualifications and provenance.}\par
\section*{IV.1.2 Reusable Information Concepts}
\phantomsection
\label{navigation:IV.1.2}
\addcontentsline{toc}{section}{IV.1.2 Reusable Information Concepts}

\subsection*{IV.1.2.1 PractitionerCapacity}
\phantomsection
\label{navigation:IV.1.2.1}
\addcontentsline{toc}{subsection}{IV.1.2.1 PractitionerCapacity}

\label{element:008f622bf9ce4d40b8a85024748bd485}
\index{PractitionerCapacity (IV.1.2.1)}
{\small DataObject · Candidate — adjudicated Domain05 Step-B baseline; final Solution Architecture unapproved · current}\par

For a Practitioner, the available PractitionerCapacity set is constrained by
the practitioner's governed profile and applicable operational context. The
Practitioner selects \textbf{one active PractitionerCapacity} for the applicable
session/operational or roster context. It supplies operational context used by
Themis when deriving effective authorisation; \textbf{it does not itself confer authority}.
Selection UI and the exact PractitionerCapacity ↔ FHIR PractitionerRole mapping
remain deferred.



\begin{quote}


\textbf{Authentication establishes who the Individual is. PractitionerCapacity
establishes the capacity in which a Practitioner is currently acting. Themis
establishes what that Identity, acting in that capacity, may do.}



\end{quote}


\begin{quote}\small\ttfamily
graph TD\\
    Auth["Authenticated Practitioner"] --> Eligible["Eligible PractitionerCapacity set"]\\
    Profile["Governed practitioner profile"] --> Eligible\\
    Context["Applicable operational / roster context"] --> Eligible\\
    Eligible -->|"Practitioner selects ONE"| Capacity["Active PractitionerCapacity: context, no authority by itself"]\\
    Capacity -->|"authorisation context"| PDP["Themis: resolve governed Criteria / assignments / policy state"]\\
    PDP --> Criteria["Identity-specific precomputed PEP-Criteria through PEP-Context"]\\
    Criteria --> PEP["Constrained local Themis PEP"]\\
    Capacity -.->|"operational responsibility context"| Responsibility["Agora: membership intent; detailed context use deferred"]\\

\end{quote}


Themis answers \textbf{“May this Identity perform this action on this information?”}
Future Agora collaboration detail concerns \textbf{“Is this Practitioner currently
responsible for receiving/responding to this operational activity?”} These are
separate responsibilities even where they consume the same PractitionerCapacity.
Receipt of an escalation or approval notification grants no information access.
Following it into Iris or another application surface requires normal Themis PEP
enforcement using current valid PEP-Criteria. U16 establishes
\hyperref[element:d1c4053e68384080b8eb79d2008632ba]{Agora control and Omilia operations};
detailed membership, escalation and approval-routing mechanics remain deferred.


{\small Recorded unresolved matters: canonicalName; exact scope remains in the canonical record.}\par
{\footnotesize \href{../../../../continuum/application/data-object/008f622bf9ce4d40b8a85024748bd485.md}{Canonical record: 008f622bf9ce4d40b8a85024748bd485}; full qualifications and provenance.}\par
\chapter*{IV.2 System-of-Systems Architecture}
\phantomsection
\label{navigation:IV.2}
\addcontentsline{toc}{chapter}{IV.2 System-of-Systems Architecture}

\section*{IV.2.1 System Map}
\phantomsection
\label{navigation:IV.2.1}
\addcontentsline{toc}{section}{IV.2.1 System Map}

The existing application responsibilities form the system map. These links identify independently defined systems and components; no new enclosing system or relationship is asserted.\par
\begin{longtable}{p{.22\linewidth}p{.73\linewidth}}
IV.3.1.1.1 & \hyperref[element:d1c4053e68384080b8eb79d2008632ba]{Agora} \\
IV.3.2.1.1 & \hyperref[element:466ad8f75a954c0581b5bba090443849]{Calliope} \\
IV.3.3.1.1 & \hyperref[element:3f64d339648246e99e8866c8089f8960]{DigitalTwin} \\
IV.3.4.1.1 & \hyperref[element:d739c75dffe2453081735710b285f03b]{Dokimasia} \\
IV.3.5.1.1 & \hyperref[element:fb3f37a6ac614415ac486a6de9d35994]{Energeia} \\
IV.3.6.1.1 & \hyperref[element:33b8421ffc444b78b2d059201a58a688]{Harmonia.Omilia.Client} \\
IV.3.7.1.1 & \hyperref[element:369907025d06492bb414a72cea834a77]{Harmonia.Omilia.Server} \\
IV.3.8.1.1 & \hyperref[element:feea742463b84021b66853d287d8916f]{Harmonia.Sylloge} \\
IV.3.9.1.1 & \hyperref[element:e927c1af96c64fcd979be659bee2cbe3]{Iris} \\
IV.3.10.1.1 & \hyperref[element:fefda110e2274884ad58906733606efd]{Kleio} \\
IV.3.11.1.1 & \hyperref[element:1af8a56358f7429fabfb57be66542a99]{Mneme} \\
IV.3.12.1.1 & \hyperref[element:7985e647bf9e4ca3ba70faf9e3e4ae02]{Mnemosyne} \\
IV.3.13.1.1 & \hyperref[element:b521762c1ea444b48cd5f66be3752d30]{Ponos} \\
IV.3.14.1.1 & \hyperref[element:e3e4147fea3640198dc251ed41ee854d]{Pylai} \\
IV.3.15.1.1 & \hyperref[element:899d34ec491948b68aad09e507dea6d1]{Themis} \\
IV.2.3.1.1.1 & \hyperref[element:195e5f10db8e419391354197b5d7a8cc]{Petasos} \\
\end{longtable}
\section*{IV.2.2 System Responsibility Map}
\phantomsection
\label{navigation:IV.2.2}
\addcontentsline{toc}{section}{IV.2.2 System Responsibility Map}

\subsection*{IV.2.2.1 Provider Registry solution collaboration}
\phantomsection
\label{navigation:IV.2.2.1}
\addcontentsline{toc}{subsection}{IV.2.2.1 Provider Registry solution collaboration}

\label{element:9305dc8d55b34b70950190e28b5464cf}
\index{Provider Registry solution collaboration (IV.2.2.1)}
{\small ApplicationCollaboration · Candidate — accepted Step-B baseline; final Solution Architecture unapproved · current}\par

\begin{itemize}


\item \textbf{Role/responsibility:} documented application solution for directory resources
  and governed changes across existing components, rather than a new enclosing
  platform component. \href{\detokenize{../../../../../docs/markdown/05-information-systems-architecture/sources.md#l15}}{L15/L18}.

\item \textbf{Features/contained constructs:} synchronous read/search, request-for-change,
  per-resource validation/identifier uniqueness/reference checks, candidate automated
  approval, durable commit and observable task progress. Seven per-resource Ergon
  activities and a registry change sequence are documented behavioural specialisations;
  Erga supplies the distinct collection/assembly/library responsibility.

\item \textbf{Information:} Practitioner, PractitionerRole, Organization, Location,
  HealthcareService, Endpoint and Group; change Pragma/input/output, errors and
  resulting versions. Endpoint and Group remain independently referenced resources.

\item \textbf{Persistence:} clinical/FHIR Mnemosyne slice; current-row versus history
  structures unclear. Iris is presentation-only; server-side governance remains
  with Mnemosyne Clinical, Ergon activities within Erga, and Themis under
  AGENTS Invariant 3.

\item \textbf{Collaboration/exclusions:} Pylai interaction, Themis, Petasos, Ponos/Praxis/Erga,
  Mneme/Mnemosyne and Iris Administration. Does not automatically own all provider,
  care-team or national-source facts, or imply an EMPI.

\item \textbf{Standing:} documented candidate domain solution, CDG-03/11/12/21/32/33.
  Direct persistence search and published async response claims are not approved
  external contracts by inclusion here.

\end{itemize}

{\small Recorded unresolved matters: canonicalName; exact scope remains in the canonical record.}\par
{\footnotesize \href{../../../../continuum/application/application-collaboration/9305dc8d55b34b70950190e28b5464cf.md}{Canonical record: 9305dc8d55b34b70950190e28b5464cf}; full qualifications and provenance.}\par
\section*{IV.2.3 Middleware Architecture}
\phantomsection
\label{navigation:IV.2.3}
\addcontentsline{toc}{section}{IV.2.3 Middleware Architecture}

\subsection*{IV.2.3.1 Petasos}
\phantomsection
\label{navigation:IV.2.3.1}
\addcontentsline{toc}{subsection}{IV.2.3.1 Petasos}

\subsubsection*{IV.2.3.1.1 Component Architecture}
\phantomsection
\label{navigation:IV.2.3.1.1}
\addcontentsline{toc}{subsubsection}{IV.2.3.1.1 Component Architecture}

\paragraph*{IV.2.3.1.1.1 Petasos}
\phantomsection
\label{navigation:IV.2.3.1.1.1}
\addcontentsline{toc}{paragraph}{IV.2.3.1.1.1 Petasos}

\label{element:195e5f10db8e419391354197b5d7a8cc}
\index{Petasos (IV.2.3.1.1.1)}
{\small ApplicationComponent · Candidate — accepted Step-B baseline; final Solution Architecture unapproved · current}\par

\begin{itemize}


\item \textbf{Role/responsibility:} \textbf{Harmonia's durable message/event queue-management
  component}, a first-class shared messaging/event capability. Petasos manages
  creation, operation and operational state/status of its durable queues.
  \href{\detokenize{../../../../../docs/markdown/05-information-systems-architecture/sources.md#u09}}{U09} extends the U05 handoff boundary, preserving AX-10.

\item \textbf{Purpose-defined queues:} purposes include inter-Ergon information flow,
  defined communication/event sharing, notification/event sharing and other
  architecturally established Harmonia event-sharing services. Queue purpose is
  architecturally defined; queues are not arbitrary generic queues. Final taxonomy,
  names, registration/configuration and topology remain unresolved.

\item \textbf{Features:} queue management and durable Parcel communication/event handling;
  persistent handoff obligations, bounded backlog, destination mechanics and
  handoff redelivery/recovery/dead-letter and designated replay responsibilities
  remain. Detailed delivery/ordering/duplicate/replay/recovery contracts are not
  newly selected; CDG-16/17/19/34/37/44 retain uncertainty.

\end{itemize}






\par\medskip\textbf{PetasosParcel and information ownership}\par


Every Petasos queue carries the single common datatype \textbf{PetasosParcel}.
It represents the queued Harmonia communication/event and is conceptually
analogous to \textbf{FHIR Communication}, without established literal equivalence.
No schema or mapping to historical PetasosMessage/ErgonEvent is inferred.



\begin{quote}\small\ttfamily
PetasosParcel          pass-by-value; the queued communication/event\\
Referenced Information predominantly pass-by-reference;\\
                       remains managed by its authoritative component\\

\end{quote}


\textbf{PetasosParcel is passed by value.} Petasos does not transport only references
and does not duplicate all associated information into the Parcel. Petasos owns
Parcel messaging/event semantics and queue operations; references confer no
ownership of the referenced information.



For the \textbf{ActionableTask handoff special case}, ActionableTask/Pragma and payload
remain managed through Mneme. A Parcel may reference that ActionableTask and is
itself queued by value. Petasos does not transport another consumer's
FulfillmentTask; each consumer undertaking the ActionableTask establishes its own
distinct FulfillmentTask. The established no-payload-copy rule for this case
does not make all Parcels reference-only. \href{\detokenize{../../../../../docs/markdown/05-information-systems-architecture/pathways-journeys-and-handoff.md#persistent-handoff}}{Handoff view}.



\begin{itemize}


\item \textbf{Persistence:} durable queues/Parcels and handoff recovery state are Petasos
  responsibility; referenced managed-information durable establishment remains
  Mnemosyne. Queue persistence mechanisms are not selected and queues do not
  become automatically permanent audit evidence.

\item \textbf{Operational observation:} Petasos owns queue operational state/status.
  Energeia may \textbf{observe/report} relevant state, without acquiring queue management
  or autonomous queue control. Detailed metrics/state structures remain unresolved.

\item \textbf{Themis collaboration:} for a governed FORWARD action, Petasos is the host/
  enforcement location because it controls the queue operation. Applicable
  enforcement SHALL use \textbf{Themis-owned PEP capability} and relevant \textbf{PEP-Criteria through PEP-Context}. Themis owns authorisation/enforcement semantics; Petasos owns the
  operation. Detailed protocol remains unresolved.

\item \textbf{Other collaborations/exclusions:} serves Ergon, ingress, execution,
  publication, Agora and architecturally established communication/notification
  services. It is not subordinate to Ponos, ActionableTask-only, a reference
  transport helper or merely a generic enterprise broker abstraction. Mneme owns
  Publishing/Subscription/Pattern semantic selection; Ponos owns execution.

\item \textbf{Standing:} U09 resolves queue responsibility and common datatype/pass semantics;
  \href{\detokenize{../../../../../docs/markdown/05-information-systems-architecture/gaps.md#petasos-detailed-contracts}}{CDG-44} retains schema, queue and interaction
  detail. Internal middleware remains Domain05; external protocols remain Domain06.

\end{itemize}

{\small Recorded unresolved matters: canonicalName; exact scope remains in the canonical record.}\par
{\footnotesize \href{../../../../continuum/application/application-component/195e5f10db8e419391354197b5d7a8cc.md}{Canonical record: 195e5f10db8e419391354197b5d7a8cc}; full qualifications and provenance.}\par
\subsubsection*{IV.2.3.1.2 Information Architecture}
\phantomsection
\label{navigation:IV.2.3.1.2}
\addcontentsline{toc}{subsubsection}{IV.2.3.1.2 Information Architecture}

\paragraph*{IV.2.3.1.2.1 PetasosParcel Topic}
\phantomsection
\label{navigation:IV.2.3.1.2.1}
\addcontentsline{toc}{paragraph}{IV.2.3.1.2.1 PetasosParcel Topic}

\label{element:f0a88f05a0cb438cb42572ec07fb8513}
\index{PetasosParcel Topic (IV.2.3.1.2.1)}
{\small DataObject · Bounded adjudicated application information meaning; Domain05 remains the accepted candidate Step-B baseline, not final Solution Architecture. Detailed representation is deferred; only this defined logical object is projected. · current}\par

\begin{quote}


\textbf{Mneme determines semantic eligibility and eligible recipient/consumer
relationships under Publishing, Subscription and Pattern semantics. Petasos
establishes and manages the concrete persistent handoff obligation to the
destination and reliably effects that handoff. Ponos manages execution.}



\end{quote}


Eligible consumer relationships express semantic selection. They do not assign
queue binding, destination mechanics or concrete handoff state to Mneme.



Petasos is Harmonia's durable message/event queue-management component.
It creates, operates and owns the operational state/status of \textbf{purpose-defined
queues} for inter-Ergon flow and other architecturally established communication,
notification and event sharing. Every queue carries \textbf{PetasosParcel by value};
associated information is predominantly by reference to governed information
managed elsewhere. Parcel is conceptually analogous to FHIR Communication;
literal equivalence and detailed schema are not established.



For the ActionableTask handoff special case, governed ActionableTask/Pragma and
Information Unit payload remain managed through Mneme. A by-value Parcel may
carry references and communication/handoff metadata required for resolution and
recovery. Petasos acquires no management authority over the referenced information
and does not copy the managed ActionableTask payload in this pattern.



\begin{quote}\small\ttfamily
Mneme manages ActionableTask / Pragma and payload\\
  --> semantic selection / eligibility\\
Petasos persistently manages handoff obligation\\
  --> queues PetasosParcel by value, referring to ActionableTask\\
Consumer resolves governed information through Mneme\\
  --> establishes its own distinct FulfillmentTask undertaking\\
Ponos manages execution / progression\\
  --> Ergon supplies bounded behaviour within Ergo's contract\\

\end{quote}


\textbf{Petasos distributes the opportunity to undertake work, not another consumer's
FulfillmentTask}, in this pattern. Each consumer undertaking the ActionableTask
establishes a distinct FulfillmentTask; no undertaking is transported on its behalf.
The Ergo framework works with \textbf{Mneme Context Load} to materialise anticipated
values before \textbf{Context Loader} exposes them to business logic. \textbf{Context Unloader}
returns changed working values and new Content as \textbf{Pragma Output} across that
boundary; \textbf{Mneme Context Save} manages subsequent consequences. The Unloader is
not persistence. Bounded Adhoc Resource Access/Search remains supported; Mnemosyne
establishes/reports durable state. \hyperref[element:49c9eb2bef0c47b8b5328621d5fc5479]{U10/U11 lifecycle}. Handoff acceptance, information commit, undertaking
completion, work resolution and external acknowledgement remain distinct.



ActionableTask handoff is one Petasos use case. Communication/Notification Events
and other established event-sharing purposes use the common PetasosParcel datatype
without becoming ActionableTasks or losing their own semantics. The no-payload-copy
rule above SHALL NOT be generalised to reference-only Petasos transport or to
copying all referenced information into every Parcel. Final purpose taxonomy,
names, Parcel schema and detailed delivery contracts remain \href{\detokenize{../../../../../docs/markdown/05-information-systems-architecture/gaps.md#petasos-detailed-contracts}}{CDG-44}.



Petasos owns queue operational state/status; Energeia may observe/report it,
without acquiring queue management or autonomous control. For governed FORWARD,
Petasos hosts enforcement because it controls the operation, using \textbf{Themis-owned
PEP capability and relevant PEP-Criteria through PEP-Context}. Themis owns semantics; Petasos
owns the queue operation. Detailed protocols/state metrics remain unresolved.



Source standing: Bounded adjudicated application information meaning; Domain05 remains the accepted candidate Step-B baseline, not final Solution Architecture. Detailed representation is deferred; only this defined logical object is projected.


{\footnotesize \href{../../../../continuum/application/data-object/f0a88f05a0cb438cb42572ec07fb8513.md}{Canonical record: f0a88f05a0cb438cb42572ec07fb8513}; full qualifications and provenance.}\par
\chapter*{IV.3 Applications}
\phantomsection
\label{navigation:IV.3}
\addcontentsline{toc}{chapter}{IV.3 Applications}

\section*{IV.3.1 Agora}
\phantomsection
\label{navigation:IV.3.1}
\addcontentsline{toc}{section}{IV.3.1 Agora}

\subsection*{IV.3.1.1 Application Overview}
\phantomsection
\label{navigation:IV.3.1.1}
\addcontentsline{toc}{subsection}{IV.3.1.1 Application Overview}

\subsubsection*{IV.3.1.1.1 Agora}
\phantomsection
\label{navigation:IV.3.1.1.1}
\addcontentsline{toc}{subsubsection}{IV.3.1.1.1 Agora}

\label{element:d1c4053e68384080b8eb79d2008632ba}
\index{Agora (IV.3.1.1.1)}
{\small ApplicationComponent · Candidate — accepted Step-B baseline; final Solution Architecture unapproved · current}\par

\begin{itemize}


\item \textbf{Role/responsibility:} \textbf{Harmonia.Agora}, collaboration control and semantics;
  native Teams/Slack-like collaboration. \textbf{Agora controls; Omilia operates.}
  \hyperref[element:d1c4053e68384080b8eb79d2008632ba]{U16 canonical view}
  replaces the earlier projection-only account.

\item \textbf{Features:} controls designated managed Room/Space lifecycle and membership
  intent and Harmonia publication; listens to dialogue/activity in \textbf{all} managed
  Rooms/Spaces through Omilia.Server's Application Services API. Defined meaningful
  interactions may be captured as Communication Information Units and/or Pragma
  for governed Harmonia processing. Not every message becomes work.

\item \textbf{Information:} owns/manages native collaboration record/state: dialogue/actions,
  Room/Space and membership state, references/activity and native collaboration
  media. Information merely conveyed through Agora retains its originating owner.
  Capture derives new governed information/activity; original dialogue remains Agora's.

\item \textbf{Membership:} Agora controls intent; Omilia.Server performs membership
  operations. Roster Management is SEP; only applicable Shift Start/End periods
  are required for R1.x/R2.x operational context. No roster/workforce authority or
  inference that PractitionerCapacity/Provider Registry is a roster.

\item \textbf{Persistence/audit:} application ownership is established, not a storage
  allocation. AX-05 remains controlling; collaboration persistence/media mechanics,
  mapping/transaction structures and retention are deferred. Dialogue retention
  and logging are distinct from audit; applicable occurrences go to Kleio, which
  retains authoritative Audit Information Management and existing capture rules.

\item \textbf{Collaborations/exclusions:} Omilia.Server performs underlying operations;
  Agora acquires no ownership of those mechanics. Petasos mediates execution
  collaboration, with no direct Ponos dependency. AGENTS Invariant 10 preserves
  Themis default-deny, Matrix DTO encapsulation and zero-PHI Room metadata;
  diagnostic logging remains zero-PHI. Membership/notification confers no access
  grant, clinical-work authority or originating-information ownership.

\item \textbf{Standing:} U16 adjudicated candidate; exact taxonomy, membership derivation,
  API, representations, persistence, security/privacy and human-experience
  integration remain \hyperref[element:d1c4053e68384080b8eb79d2008632ba]{deferred};
  no direct Iris-to-Agora/Omilia interaction is permitted (U19).
  Partial transaction effects remain CDG-19. \hyperref[element:feea742463b84021b66853d287d8916f]{Sylloge C49} owns F31/D15
  Binary Content Management under U17; Agora retains collaboration references.
  Historical class/table names are implementation evidence.

\end{itemize}

{\small Recorded unresolved matters: canonicalName; exact scope remains in the canonical record.}\par
{\footnotesize \href{../../../../continuum/application/application-component/d1c4053e68384080b8eb79d2008632ba.md}{Canonical record: d1c4053e68384080b8eb79d2008632ba}; full qualifications and provenance.}\par
\section*{IV.3.2 Calliope}
\phantomsection
\label{navigation:IV.3.2}
\addcontentsline{toc}{section}{IV.3.2 Calliope}

\subsection*{IV.3.2.1 Application Overview}
\phantomsection
\label{navigation:IV.3.2.1}
\addcontentsline{toc}{subsection}{IV.3.2.1 Application Overview}

\subsubsection*{IV.3.2.1.1 Calliope}
\phantomsection
\label{navigation:IV.3.2.1.1}
\addcontentsline{toc}{subsubsection}{IV.3.2.1.1 Calliope}

\label{element:466ad8f75a954c0581b5bba090443849}
\index{Calliope (IV.3.2.1.1)}
{\small ApplicationComponent · Candidate — accepted Step-B baseline; final Solution Architecture unapproved · current}\par

\begin{itemize}


\item \textbf{Role/responsibility:} management/governance of foundational models, definitions
  and mutable datasets, including first-class Security and Assurance Nomenclature
  and authoritative review of ontology sets for Harmonia use. \href{\detokenize{../../../../../docs/markdown/05-information-systems-architecture/sources.md#u13}}{U13},
  \href{\detokenize{../../../../../docs/markdown/05-information-systems-architecture/calliope-models-and-datasets.md}}{canonical Calliope view}.

\item \textbf{Distinct submodules:} \textbf{Calliope-Models C41} supplies design/build-time
  Models, Metadata Types, Criteria Definitions, Topics, Filter Types, Filter Elements,
  core Harmonia element libraries and security/assurance definitions.
  \textbf{Calliope-DataSets C42} manages runtime mutable ValueSets, CodeSets, Structured
  Metadata and applicable governed ontology content. No common lifecycle or
  undifferentiated metadata repository is inferred.

\item \textbf{Authority:} development representations support type-safe human/IDE/AI-assisted
  programming; current membership/content uses runtime authority. Models/build
  definitions establish structural access contracts; DataSets do not bootstrap them.
  Definition ownership confers neither runtime execution nor originating clinical authority.

\item \textbf{Persistence/collaboration:} DataSets uses established Mneme/Mnemosyne
  responsibilities; Calliope is not a Persistence Manager. Governed workflow may
  reuse Ponos/Ergo/Praxis, without another engine or detailed management UI.
  Themis decides/enforces, Dokimasia assures and Kleio audits.

\item \textbf{Standing:} U13 settles the two responsibilities, not their packaging, runtime
  ontology boundary or representation/workflow mechanisms (CDG-24/49). Existing
  pure-library dependency constraints apply to Models; U15 resolves the
  \href{\detokenize{../../../../../docs/markdown/05-information-systems-architecture/calliope-models-and-datasets.md#calliope-realisation}}{realisation boundary}.
  General converter-support
  versus activity allocation remains CDG-05; no terminology product is selected.

\end{itemize}






\par\medskip\textbf{Controlled authorisation nomenclature}\par


\href{\detokenize{../../../../../docs/markdown/05-information-systems-architecture/sources.md#u12}}{U12} establishes these \textbf{release-bounded ValueSets}:



\begin{longtable}{>{\raggedright\arraybackslash}p{\dimexpr\linewidth/2-2\tabcolsep\relax}>{\raggedright\arraybackslash}p{\dimexpr\linewidth/2-2\tabcolsep\relax}}
ValueSet & Architectural shorthand / meaning \\
\texttt{\seqsplit{Application.Function.AuthorisationPoint.Names}} & \textbf{Function.Name}: architecturally significant application AuthorisationPoints at which enforcement is meaningful. An AuthorisationPoint identifies a governed function; it is not itself a permission. \\
\texttt{\seqsplit{Application.Function.AllowedAction.Names}} & \textbf{Function.Scope}: the bounded set of actions applicable at AuthorisationPoints. Existing concepts include \textbf{ACCESS}, \textbf{CHANGE} and \textbf{FORWARD}; these are not asserted to be the complete final ValueSet. \\
\texttt{\seqsplit{Information.Compartments}} & Governed information compartments; an applicable dimension is expressed as \textbf{Information.Compartment}. \\
\texttt{\seqsplit{Information.Classification}} & Governed information classifications. \\
\end{longtable}


\texttt{\seqsplit{Application.Function.AuthorisationPoint.Names}} SHALL be comprehensive enough
to satisfy business security requirements for the release while constrained
enough to avoid unnecessary management, development and support burden.
Complete release-specific contents remain \href{\detokenize{../../../../../docs/markdown/05-information-systems-architecture/gaps.md#identity-deferred-design}}{deferred}.



\textbf{Calliope defines/governs the nomenclature. Themis owns policy association,
resolution and enforcement.} ValueSet membership alone grants no authority
and does not establish every Cartesian combination as meaningful or permitted.



The controlled vocabularies fall within \hyperref[element:466ad8f75a954c0581b5bba090443849]{Calliope-DataSets};
their structural definitions/contracts belong to Calliope-Models. This allocation
does not reopen U12 or permit unrestricted runtime changes to release-bounded
content; detailed release/change rules remain deferred.


{\small Recorded unresolved matters: canonicalName; exact scope remains in the canonical record.}\par
{\footnotesize \href{../../../../continuum/application/application-component/466ad8f75a954c0581b5bba090443849.md}{Canonical record: 466ad8f75a954c0581b5bba090443849}; full qualifications and provenance.}\par
\subsection*{IV.3.2.2 Component Architecture}
\phantomsection
\label{navigation:IV.3.2.2}
\addcontentsline{toc}{subsection}{IV.3.2.2 Component Architecture}

\subsubsection*{IV.3.2.2.1 Calliope-Models}
\phantomsection
\label{navigation:IV.3.2.2.1}
\addcontentsline{toc}{subsubsection}{IV.3.2.2.1 Calliope-Models}

\label{element:307b3440a507403f8bc214d290bffb6a}
\index{Calliope-Models (IV.3.2.2.1)}
{\small ApplicationComponent · Candidate — accepted Step-B baseline; final Solution Architecture unapproved · current}\par

Distinct Calliope design/build-time submodule: Models, Metadata Types, Criteria
Definitions, Topics, Filter Types/Elements, core element libraries and security/
assurance nomenclature definitions. Design → Code → Build → Test → Release → Deploy.
Provides structural access contracts and useful development representations of
known runtime DataSets as normal library dependencies; those representations do
not establish current membership. Models' pure-library dependency semantics remain.
Definition ownership is not runtime operational ownership. \hyperref[element:307b3440a507403f8bc214d290bffb6a]{U13 lifecycle and authority}.


{\small Recorded unresolved matters: canonicalName; exact scope remains in the canonical record.}\par
{\footnotesize \href{../../../../continuum/application/application-component/307b3440a507403f8bc214d290bffb6a.md}{Canonical record: 307b3440a507403f8bc214d290bffb6a}; full qualifications and provenance.}\par
\subsubsection*{IV.3.2.2.2 Calliope-DataSets}
\phantomsection
\label{navigation:IV.3.2.2.2}
\addcontentsline{toc}{subsubsection}{IV.3.2.2.2 Calliope-DataSets}

\label{element:53bf6815c7a54ed0abde9b1aad311843}
\index{Calliope-DataSets (IV.3.2.2.2)}
{\small ApplicationComponent · Candidate — accepted Step-B baseline; final Solution Architecture unapproved · current}\par

Distinct Calliope runtime management/repository responsibility for mutable governed
ValueSets, CodeSets, Structured Metadata and applicable ontology content. Owns
governed content, not persistence machinery; uses established Mneme/Mnemosyne and,
where workflow is required, existing Ponos/Ergo/Praxis. Models/build definitions
establish its structural contracts. Detailed ontology boundary, workflow, UI and
physical realisation remain \hyperref[element:466ad8f75a954c0581b5bba090443849]{deferred}.


{\small Recorded unresolved matters: canonicalName; exact scope remains in the canonical record.}\par
{\footnotesize \href{../../../../continuum/application/application-component/53bf6815c7a54ed0abde9b1aad311843.md}{Canonical record: 53bf6815c7a54ed0abde9b1aad311843}; full qualifications and provenance.}\par
\section*{IV.3.3 DigitalTwin}
\phantomsection
\label{navigation:IV.3.3}
\addcontentsline{toc}{section}{IV.3.3 DigitalTwin}

\subsection*{IV.3.3.1 Application Overview}
\phantomsection
\label{navigation:IV.3.3.1}
\addcontentsline{toc}{subsection}{IV.3.3.1 Application Overview}

\subsubsection*{IV.3.3.1.1 DigitalTwin}
\phantomsection
\label{navigation:IV.3.3.1.1}
\addcontentsline{toc}{subsubsection}{IV.3.3.1.1 DigitalTwin}

\label{element:3f64d339648246e99e8866c8089f8960}
\index{DigitalTwin (IV.3.3.1.1)}
{\small ApplicationComponent · Candidate — accepted Step-B baseline; final Solution Architecture unapproved · current}\par

\begin{itemize}


\item \textbf{Role/features:} specialised Ponos-based execution component for entity-specific
  operational coordination associated with Energeia. \href{\detokenize{../../../../../docs/markdown/05-information-systems-architecture/sources.md#u02}}{U02} establishes
  specialisation/reuse, not containment by a Ponos runtime. The Digital Twin
  responsibility in \href{\detokenize{../../../../../docs/markdown/05-information-systems-architecture/sources.md#a02}}{A02} remains controlling: entity-related
  execution SHOULD use the Twin unless explicitly justified otherwise.

\item \textbf{Twin-specific execution:} one input queue containing Stimulus and Effector.
  At most one queued input is processed at a time; interested Behaviours execute
  \textbf{sequentially}, subject to the failure dispositions below. Each receives its
  own BehaviourStimulus/BehaviourEffector and fresh Praxis-defined Context after
  its predecessor's completion/save on the normal path. No shared immutable
  input-level Context.
  \href{\detokenize{../../../../../docs/markdown/05-information-systems-architecture/design-notes/DesignNote-Digital-Twin-Behaviour-and-Execution-Semantics.md}}{U19 serialisation and lifecycle}.
  Bounded resources service runnable pending work; stopped Twins retain pending
  inputs. Concurrency is across independent Twins, without a thread per Twin.

\item \textbf{Sequence/failure/recovery:} R1.x/R2.x framework loads/exposes Behaviours and
  provides/maintains the sequencing API; developer configuration supplies policy.
  Without it, order is unspecified. The applicable sequence remains stable for
  the current input; changes apply to subsequently processed inputs.
  \texttt{\seqsplit{ContinueBehaviourOnFailure}} contains failure to the current execution while
  keeping the Behaviour active and other interested Behaviours and queue processing
  unaffected. \texttt{\seqsplit{StopBehaviourOnFailure}} stops the failing Behaviour immediately,
  lets remaining interested Behaviours finish the current input, then halts the
  Twin. \texttt{\seqsplit{StopDigitalTwinOnFailure}} terminates the failing Behaviour and halts the
  Twin immediately, preventing remaining current-input and subsequent queued
  processing. Stopped Twins do not
  automatically resume: authorised administrative input through Iris/Iris-BackEnd
  or Omilia/Agora leads to Administrative Function → Administrative Ergon →
  governed Twin operation. Neither channel directly manipulates the Twin.
  Detailed recovery of a stopped Behaviour remains deferred (CDG-54).

\item \textbf{Information:} uses context drawn from governed entity information and
  relationships; entity/resource/Twin identity are not aliases. Historical patient,
  practitioner, wardsperson, organisation, care-team, ward and bed examples are
  illustrative archetypes, not an approved exhaustive component catalogue.

\item \textbf{Persistence/significant information:} definition/configuration is a C\&G
  classification candidate; operational state/activity is an O\&A candidate under
  U06. Detailed coordination-state Data Objects, activation/deactivation,
  preservation and transient/durable qualification remain CDG-07. A Twin is not
  inherently a permanent object, table or FHIR resource.

\item \textbf{Collaboration/exclusions:} executes Behaviour C25, which specialises Ergon
  and retains all common Ergo/Ergon facilities. Uses Mneme Pragma Management and
  candidate Stimulus/Effector specialisations; Ponos supplies common execution,
  Petasos supplies messaging/recoverable handoff. Follows the
  \href{\detokenize{../../../../../docs/markdown/05-information-systems-architecture/README.md#specialisation-principle}}{specialisation principle}; adds genuinely
  Twin-specific semantics instead of parallel platform mechanisms.
  Not a second engine, unmanaged persistent actor, separate
  policy/search service or authority over originating clinical facts/clinical work.

\item \textbf{Standing:} U02 specialisation and U19 Behaviour execution semantics established;
  representation/runtime details remain CDG-07/54. DigitalTwin and Dokimasia are sibling
  Ponos specialisations; neither contains the other or manages Pragma by execution.

\end{itemize}

{\small Recorded unresolved matters: canonicalName; exact scope remains in the canonical record.}\par
{\footnotesize \href{../../../../continuum/application/application-component/3f64d339648246e99e8866c8089f8960.md}{Canonical record: 3f64d339648246e99e8866c8089f8960}; full qualifications and provenance.}\par
\subsection*{IV.3.3.2 Component Architecture}
\phantomsection
\label{navigation:IV.3.3.2}
\addcontentsline{toc}{subsection}{IV.3.3.2 Component Architecture}

\subsubsection*{IV.3.3.2.1 Behaviour}
\phantomsection
\label{navigation:IV.3.3.2.1}
\addcontentsline{toc}{subsubsection}{IV.3.3.2.1 Behaviour}

\label{element:a26f06d66bec4981b266e202c46a9afb}
\index{Behaviour (IV.3.3.2.1)}
{\small ApplicationComponent · Candidate — accepted Step-B baseline; final Solution Architecture unapproved · current}\par

\begin{itemize}


\item \textbf{Realisation:} each DigitalTwin has \textbf{Behaviour [1..*]}. Each Behaviour is
  realised through \textbf{Praxis → Pathway → Ergo [exactly 1]}. A one-Ergo Praxis
  retains required Context, expected criteria/outcomes, fulfilment semantics and
  housekeeping; it is neither an orchestrator nor a choreographer.
  \href{\detokenize{../../../../../docs/markdown/05-information-systems-architecture/design-notes/DesignNote-Digital-Twin-Behaviour-and-Execution-Semantics.md}}{U19 DesignNote}.

\item \textbf{Role/features:} Ergon specialisation for DigitalTwin operation. It inherits/
  reuses Context Loading and Context Unloading, process-management mechanisms,
  Petasos/platform integration, assurance integration, logging/observability,
  reporting, checkpoint/platform integration and other common Ergo/Ergon facilities.
  \href{\detokenize{../../../../../docs/markdown/05-information-systems-architecture/sources.md#u02}}{U02}, \href{\detokenize{../../../../../docs/markdown/05-information-systems-architecture/README.md#specialisation-principle}}{common specialisation principle}.

\item \textbf{Information/persistence/collaboration:} receives its own identifiable
  \textbf{BehaviourStimulus / BehaviourEffector} derived from an input in which it has
  registered interest. Operates relative to the Twin's singular represented
  entity. Its Praxis-defined Context is loaded fresh after the preceding
  Behaviour completes and saves its consequences; Mneme manages Load/Save and
  Pragma, Mnemosyne establishes durable state and Petasos handles handoff.

\item \textbf{Execution/output/failure:} interested Behaviours execute sequentially for
  each input. Outputs may enter normal governed processing, resolving Twin-relative
  references only where required. Same-Twin Stimulus/Effector outputs join its
  normal queue without recursive interruption. The developer's
  \texttt{\seqsplit{ContinueBehaviourOnFailure}} contains failure to the current execution; the
  Behaviour remains active, other interested Behaviours and normal queue processing
  continue. \texttt{\seqsplit{StopBehaviourOnFailure}} stops the failing Behaviour immediately,
  lets remaining interested Behaviours complete the current input, then halts
  the Twin. \texttt{\seqsplit{StopDigitalTwinOnFailure}} terminates the failing Behaviour and halts
  the Twin immediately; remaining current-input and subsequent queued processing
  stop. Sequencing policy is developer-owned through the framework API; the
  applicable sequence stays stable throughout an input, and configuration changes
  apply to subsequent inputs. No priority metadata or Pathway orchestration (U19).

\item \textbf{Standing/exclusions:} discrete candidate activity specialisation, not an
  independent activity framework parallel to Ergo/Ergon. DigitalTwin adds only
  genuinely Twin-specific behaviour instead of duplicating platform mechanisms;
  Behaviour acquires no standalone Information Resource merely by being activity.

\end{itemize}

{\small Recorded unresolved matters: canonicalName; exact scope remains in the canonical record.}\par
{\footnotesize \href{../../../../continuum/application/application-component/a26f06d66bec4981b266e202c46a9afb.md}{Canonical record: a26f06d66bec4981b266e202c46a9afb}; full qualifications and provenance.}\par
\section*{IV.3.4 Dokimasia}
\phantomsection
\label{navigation:IV.3.4}
\addcontentsline{toc}{section}{IV.3.4 Dokimasia}

\subsection*{IV.3.4.1 Application Overview}
\phantomsection
\label{navigation:IV.3.4.1}
\addcontentsline{toc}{subsection}{IV.3.4.1 Application Overview}

\subsubsection*{IV.3.4.1.1 Dokimasia}
\phantomsection
\label{navigation:IV.3.4.1.1}
\addcontentsline{toc}{subsubsection}{IV.3.4.1.1 Dokimasia}

\label{element:d739c75dffe2453081735710b285f03b}
\index{Dokimasia (IV.3.4.1.1)}
{\small ApplicationComponent · Candidate — accepted Step-B baseline; final Solution Architecture unapproved · current}\par

\begin{itemize}


\item \textbf{Role/features:} specialised Ponos-based execution component associated with Energeia
  coordinating independent governed assurance activity over a subject. Journey
  C35 is the common observed-progression evidence subject/substrate for Pathway
  and Pub/Sub
  assurance under U04/U05; this does not limit all assurance subjects to Journey.
  \href{\detokenize{../../../../../docs/markdown/05-information-systems-architecture/sources.md#u02}}{U02} settles the candidate specialisation relationship; actual
  Assurance Praxis evaluates available
  evidence against applicable basis and produces governed assurance information.
  \href{\detokenize{../../../../../docs/markdown/05-information-systems-architecture/sources.md#a04}}{A04 §2}, \href{\detokenize{../../../../../docs/markdown/05-information-systems-architecture/sources.md#l31}}{L31}. \textbf{Praxis is the assurance basis;
  Journey is the observed evidence subject} under U06. Pathway/Route expectations
  contribute through applicable Praxis rather than constituting the general basis.

\item \textbf{Information:} evidence/context/basis and assurance assessment/finding/outcome
  areas are documented. Journey observations remain distinct from inferred
  interpretation and assurance findings. Dokimasia activity/findings/results are
  O\&A classification candidates under U06; detailed Data Objects, identity and
  preservation remain unresolved. Candidate vocabulary is not promoted into an
  Information Family.

\item \textbf{Activity/collaboration:} executes discrete Dokimasia Ergon C26 specialisations,
  including candidate AI Assurance Ergon C27. Reuses Ponos execution, Ergo/Ergon
  activity, Mneme Pragma Management, process-management, platform integration,
  reporting and observability. This is specialisation/reuse, not runtime containment
  or a parallel assurance execution platform.

\item \textbf{Information/persistence:} corresponding Pragma specialisation semantics are
  required; assurance request/input and assurance result/finding are candidate
  roles whose names, schemas, cardinalities and lifecycle remain unresolved.
  Evidence-access/control independence, detailed workflow composition and durable
  structures remain CDG-08/09. Common framework reuse does not let the assured
  activity determine, suppress, manufacture or alter its own assurance outcome.

\item \textbf{Exclusions:} evidence, workflow completion and operational status are not
  assurance; custody does not make Kleio the evaluator. Dokimasia does not manage
  operational subject escalation/remediation and is not an established Twin subtype.
  AI invocation confers no authority on the model/service. Dokimasia can manage
  its own assurance work while the subject retains its operational management.

\item \textbf{Standing:} conceptual independence remains established by A04; execution
  and discrete activity responsibilities are adjudicated for the candidate by U02.
  Detailed AI architecture requires CDG-35 discovery. No upstream responsibility
  is reopened; detailed Praxis/Journey and assurance contracts remain unresolved.

\end{itemize}

{\small Recorded unresolved matters: canonicalName; exact scope remains in the canonical record.}\par
{\footnotesize \href{../../../../continuum/application/application-component/d739c75dffe2453081735710b285f03b.md}{Canonical record: d739c75dffe2453081735710b285f03b}; full qualifications and provenance.}\par
\subsection*{IV.3.4.2 Component Architecture}
\phantomsection
\label{navigation:IV.3.4.2}
\addcontentsline{toc}{subsection}{IV.3.4.2 Component Architecture}

\subsubsection*{IV.3.4.2.1 Dokimasia Ergon}
\phantomsection
\label{navigation:IV.3.4.2.1}
\addcontentsline{toc}{subsubsection}{IV.3.4.2.1 Dokimasia Ergon}

\label{element:34d7c844da3344a09e96ad821ff9fcea}
\index{Dokimasia Ergon (IV.3.4.2.1)}
{\small ApplicationComponent · Candidate — accepted Step-B baseline; final Solution Architecture unapproved · current}\par

\begin{itemize}


\item \textbf{Role/features:} discrete Ergon specialisation representing a bounded assurance
  activity executed within a Dokimasia workflow. Reuses normal Ergo/Ergon
  execution, Context Loading/Unloading, process-management, integration,
  checkpoint, logging/observability and reporting contracts. \href{\detokenize{../../../../../docs/markdown/05-information-systems-architecture/sources.md#u02}}{U02}.

\item \textbf{Information/persistence/collaboration:} uses governed assurance input,
  applicable basis and evidence, and may produce assurance information through
  Mneme Pragma Management. No detailed schema, evidence-access API, storage or
  additional Information Resource is inferred. Executed through Dokimasia's
  specialised Ponos model.

\item \textbf{Standing/exclusions:} discrete candidate component responsibility established;
  subject to \href{\detokenize{../../../../../docs/markdown/05-information-systems-architecture/component-register.md#dokimasia-terminal-boundaries}}{Dokimasia terminal constraints}.
  AI Assurance Ergon C27 specialises this component. Other potential assurance
  classifications, including deterministic/human subclasses, require adjudication.

\end{itemize}

{\small Recorded unresolved matters: canonicalName; exact scope remains in the canonical record.}\par
{\footnotesize \href{../../../../continuum/application/application-component/34d7c844da3344a09e96ad821ff9fcea.md}{Canonical record: 34d7c844da3344a09e96ad821ff9fcea}; full qualifications and provenance.}\par
\subsubsection*{IV.3.4.2.2 AI Assurance Ergon}
\phantomsection
\label{navigation:IV.3.4.2.2}
\addcontentsline{toc}{subsubsection}{IV.3.4.2.2 AI Assurance Ergon}

\label{element:8fe50975372b4d399a4bcb9f3d9348f7}
\index{AI Assurance Ergon (IV.3.4.2.2)}
{\small ApplicationComponent · Candidate — accepted Step-B baseline; final Solution Architecture unapproved · current}\par

\begin{itemize}


\item \textbf{Role/features:} candidate Dokimasia Ergon specialisation where AI-based
  processing is itself an explicit governed participant in Dokimasia workflow.
  Its existence and discrete responsibility are established by \href{\detokenize{../../../../../docs/markdown/05-information-systems-architecture/sources.md#u02}}{U02};
  its detailed Information Systems architecture is not finalised.

\item \textbf{Distinction:} ordinary Ergon may call AI internally as an implementation
  choice. AI Assurance Ergon makes that AI activity an explicit governed assurance
  workflow activity. This allocates neither every use of AI to this component
  nor a mandatory AI step to every Dokimasia workflow.

\item \textbf{Information/persistence/collaboration:} inherits C26/C23/C05 contracts and
  common Pragma Management, process-management, platform integration and reporting.
  AI invocation confers no originating domain or assurance-determination authority
  on the service/model. Context, disclosure, results, provenance/evidence and
  detailed workflow interactions require \href{\detokenize{../../../../../docs/markdown/05-information-systems-architecture/gaps.md#dokimasia-ai-discovery}}{CDG-35 discovery}.

\item \textbf{Standing/exclusions:} established component boundary with candidate detailed
  design; no LLM, agent framework, SDK, model, provider or service technology
  selected. Discovery precedes final implementation modelling. Inherits both
  Dokimasia terminal rules; unresolved design does not invalidate its boundary.

\end{itemize}

{\small Recorded unresolved matters: canonicalName; exact scope remains in the canonical record.}\par
{\footnotesize \href{../../../../continuum/application/application-component/8fe50975372b4d399a4bcb9f3d9348f7.md}{Canonical record: 8fe50975372b4d399a4bcb9f3d9348f7}; full qualifications and provenance.}\par
\section*{IV.3.5 Energeia}
\phantomsection
\label{navigation:IV.3.5}
\addcontentsline{toc}{section}{IV.3.5 Energeia}

\subsection*{IV.3.5.1 Application Overview}
\phantomsection
\label{navigation:IV.3.5.1}
\addcontentsline{toc}{subsection}{IV.3.5.1 Application Overview}

\subsubsection*{IV.3.5.1.1 Energeia}
\phantomsection
\label{navigation:IV.3.5.1.1}
\addcontentsline{toc}{subsubsection}{IV.3.5.1.1 Energeia}

\label{element:fb3f37a6ac614415ac486a6de9d35994}
\index{Energeia (IV.3.5.1.1)}
{\small ApplicationComponent · Candidate — accepted Step-B baseline; final Solution Architecture unapproved · current}\par

\begin{itemize}


\item \textbf{Role/responsibility:} discrete Software Component owning Harmonia operational
  management: “Is Harmonia itself operating as expected?” It observes work-to-be-done,
  work-being-done and work-done, outstanding ActionableTasks, FulfillmentTask
  activity, Journey progression, stalled/incomplete activity, Praxis context,
  abnormal behaviour and operational information for reporting/assurance.

\item \textbf{Candidate services/features:} \textbf{Praxis Management Services} (F11) and
  \textbf{Journey Management Services} (F38), plus operational observation
  and reporting (F39). Detailed Feature/Service/internal-construct/subordinate-
  Component decomposition remains unresolved (CDG-39); service-set names do not
  manufacture independently deployable components.

\item \textbf{Information:} owns Praxis C06 and Journey C35 at the application level.
  Mneme manages their governed information and retains ActionableTask,
  FulfillmentTask and Pragma management. Observing operational significance does
  not transfer management of that information or referenced clinical facts.

\item \textbf{Persistence:} Mnemosyne C\&G provides Praxis/configuration durability;
  Mnemosyne O\&A provides Journey/operational-assurance durability, through Mneme
  governed management. Energeia application ownership is distinct from durable
  persistence responsibility; physical structures remain unresolved.

\item \textbf{Collaboration/relationships:} associated execution/activity responsibilities
  retain U02 specialisation/reuse. Praxis aggregates Pathway [0..1] and has a
  working Aggregation relationship to Journey with cardinalities unresolved.
  Dokimasia uses Praxis as assurance basis and Journey as observed evidence
  subject, preserving independence and terminal rules. Iris presents operational
  views; presentation does not own Harmonia operational management.

\item \textbf{Release boundary:} for \textbf{R1.x/R2.x, Observe → Report}. Energeia
  SHALL NOT autonomously change execution, redistribute workload, modify Routes,
  Pathways, Pub/Sub or Praxis, intervene in resilience or adjust capacity merely
  from workload, failure, Journey, capacity, assurance or anomaly observations.
  Existing explicitly governed execution/recovery behaviour retains its own
  authority; observation grants no new adaptive authority.

\item \textbf{Separate management plane:} external IT/Ops manages infrastructure, network,
  containers/Kubernetes, JVM/runtime, databases, technical health/capacity and
  platform resilience. Energeia manages/observes Harmonia as Harmonia. This
  distinction selects no deployment topology or infrastructure integration.

\item \textbf{Deferred direction:} \textbf{R3 Adaptive Load and Resilience Service} (F40):
  Observe → Analyse → Decide → Adapt, with feedback to observation. R1.x/R2.x
  observability should permit later capability without embedding it now.
  Policy/control loops, algorithms, mechanisms, infrastructure integration and
  AI decision architecture are deliberately undesigned (CDG-41).

\item \textbf{Standing:} U06 adjudicated these candidate responsibilities; it did not
  itself establish whole-candidate acceptance or Step B. Current A4/B1 standing
  is recorded in the \href{\detokenize{../../../../../docs/markdown/05-information-systems-architecture/README.md}}{orientation}. Possible upstream omissions require
  later review, without deriving Business Architecture or capability demand here.

\end{itemize}

{\small Recorded unresolved matters: canonicalName; exact scope remains in the canonical record.}\par
{\footnotesize \href{../../../../continuum/application/application-component/fb3f37a6ac614415ac486a6de9d35994.md}{Canonical record: fb3f37a6ac614415ac486a6de9d35994}; full qualifications and provenance.}\par
\subsection*{IV.3.5.2 Information Architecture}
\phantomsection
\label{navigation:IV.3.5.2}
\addcontentsline{toc}{subsection}{IV.3.5.2 Information Architecture}

\subsubsection*{IV.3.5.2.1 Praxis}
\phantomsection
\label{navigation:IV.3.5.2.1}
\addcontentsline{toc}{subsubsection}{IV.3.5.2.1 Praxis}

\label{element:0f961462340348b290fd844357d54d67}
\index{Praxis (IV.3.5.2.1)}
{\small DataObject · Candidate — accepted Step-B baseline; final Solution Architecture unapproved · current}\par

\begin{itemize}


\item \textbf{Role/ownership:} higher-order operational construct defining expected/permitted
  context and the operational/assurance envelope. \textbf{Energeia C39 owns Praxis}
  and provides \textbf{Praxis Management Services F11} under \href{\detokenize{../../../../../docs/markdown/05-information-systems-architecture/sources.md#u06}}{U06}.
  It is not another execution engine or necessarily an independent component.

\item \textbf{Fulfilment semantics:} Pathway defines the conceptual Ergo sequence; Praxis
  defines how it is fulfilled: applicable Context definitions required by the
  Ergo framework, segment criteria/outcomes, continuation derivation semantics
  and other Praxis-level housekeeping. Ponos/Ergo performs output processing and
  routing. Context requirements are concrete in R1.x/R2.x; criteria-driven dynamic
  Context and dynamically configured reusable/generic Ergo remain deferred R3.x
  direction, without an AI-driven selection/configuration mechanism.

\item \textbf{Information/relationships:} candidate interface configuration, Topic exclusions
  and potentially inclusions, operational constraints/configuration and assurance
  rules. \textbf{Aggregation: Pathway [0..1]}; working \textbf{Aggregation: Journey} with
  cardinalities unresolved. No Composition or shared lifecycle is inferred.
  Praxis may exist without a Pathway for Publish/Subscribe activity.
  The specialised \hyperref[element:a26f06d66bec4981b266e202c46a9afb]{Digital Twin Behaviour pattern} uses a Pathway with
  exactly one Ergo; it retains full Praxis semantics and does not change the
  general Aggregation cardinality.

\item \textbf{Management/persistence:} Mneme manages governed representations; Mnemosyne
  \textbf{C\&G} provides durable persistence responsibility. Ponos consumes Praxis for
  execution and Dokimasia uses it as the primary normative assurance basis;
  neither use nor persistence transfers application ownership.

\item \textbf{Retained detail:} schema, lifecycle/versioning, activation/deactivation,
  replacement/retirement and in-flight effects remain CDG-01/06. Historical
  TaskSequence/PraxisDefinition/PraxisImplementation and A04 actual-assurance-
  activity descriptions do not establish equivalence or a definition-to-instance
  model. U06 supersedes the earlier wholly unresolved broader-Praxis standing.

\item \textbf{Logical boundary:} \href{\detokenize{../../../../../docs/markdown/05-information-systems-architecture/component-register.md#c46}}{Praxis.Horos C46}, with IngressHoros/EgressHoros,
  names the Praxis boundary. EgressHoros hands interaction state/outcome back
  before Ergo.ContextUnloader completes; it is not a technical endpoint.

\item \textbf{Exclusions:} Praxis is not Pathway, Route, a sequence of Ergo/Ergon, an
  orchestrator, choreographer, runtime engine or Pub/Sub graph. Mneme semantic
  selection, Petasos concrete persistent handoff and Ponos execution remain distinct.
  No detailed Praxis schema is invented.

\end{itemize}

{\small Recorded unresolved matters: canonicalName; exact scope remains in the canonical record.}\par
{\footnotesize \href{../../../../continuum/application/data-object/0f961462340348b290fd844357d54d67.md}{Canonical record: 0f961462340348b290fd844357d54d67}; full qualifications and provenance.}\par
\subsubsection*{IV.3.5.2.2 Journey}
\phantomsection
\label{navigation:IV.3.5.2.2}
\addcontentsline{toc}{subsubsection}{IV.3.5.2.2 Journey}

\label{element:359b901df27f4b43b5a2ff14a8484e34}
\index{Journey (IV.3.5.2.2)}
{\small DataObject · Candidate — accepted Step-B baseline; final Solution Architecture unapproved · current}\par

First-class observational construct / candidate Application Data Object: governed
observed progression reconstructed from available execution evidence, whether
Pathway-controlled, Pub/Sub or mixed. It may be graph-shaped and have no Pathway.
Primary reporting, analysis and assurance subject; Dokimasia may
interpret it under applicable basis without replacing observations with inference.
\textbf{Energeia owns Journey and Journey Management Services F38} under U06. Mneme
manages governed information; Mnemosyne \textbf{O\&A} provides durable persistence
responsibility. Journey aggregates/references observations/information from Mneme,
Petasos, Ponos, Ergon, DigitalTwin and other components/evidence mechanisms where
appropriate, retaining their ownership. Comparatively small discrete Journey
information may describe substantial activity/evidence: aggregate semantically,
without indiscriminate duplication. Detailed representation/construction, lifecycle,
retention, graph/evidence/causal/correlation/reference schema and inference contracts
remain unresolved (\href{\detokenize{../../../../../docs/markdown/05-information-systems-architecture/gaps.md#journey-contracts}}{CDG-38}); Journey is not a shard key.
Journey is neither an engine, AI narrative nor automatic permanent AuditEvent.
For Pathway-based Praxis fulfilment, continuation-Pragma lineage provides an
important observable progression thread. \textbf{Journey lineage is not complete
Pragma provenance}: independent outputs may establish their own activity,
Pub/Sub interactions, Praxes and Journeys. Mneme may manage richer relationships.
Journey is not a complete provenance graph, execution log, audit trail or Petasos
queue history. \hyperref[element:359b901df27f4b43b5a2ff14a8484e34]{U18 reconciliation}.


{\small Recorded unresolved matters: canonicalName; exact scope remains in the canonical record.}\par
{\footnotesize \href{../../../../continuum/application/data-object/359b901df27f4b43b5a2ff14a8484e34.md}{Canonical record: 359b901df27f4b43b5a2ff14a8484e34}; full qualifications and provenance.}\par
\section*{IV.3.6 Harmonia.Omilia.Client}
\phantomsection
\label{navigation:IV.3.6}
\addcontentsline{toc}{section}{IV.3.6 Harmonia.Omilia.Client}

\subsection*{IV.3.6.1 Application Overview}
\phantomsection
\label{navigation:IV.3.6.1}
\addcontentsline{toc}{subsection}{IV.3.6.1 Application Overview}

\subsubsection*{IV.3.6.1.1 Harmonia.Omilia.Client}
\phantomsection
\label{navigation:IV.3.6.1.1}
\addcontentsline{toc}{subsubsection}{IV.3.6.1.1 Harmonia.Omilia.Client}

\label{element:33b8421ffc444b78b2d059201a58a688}
\index{Harmonia.Omilia.Client (IV.3.6.1.1)}
{\small ApplicationComponent · Candidate — accepted Step-B baseline; final Solution Architecture unapproved · current}\par

\begin{itemize}


\item \textbf{Role/responsibility:} Harmonia collaboration client providing the associated
  collaboration user experience; communicates with Omilia.Server using Matrix
  protocol. Presentation does not confer information or collaboration-control
  authority.

\item \textbf{Provenance/standing:} derived from Element; U16 adjudicated candidate.
  Detailed modifications and human-experience integration remain deferred;
  Iris has no direct Agora/Omilia interaction (U19).

\end{itemize}

{\small Recorded unresolved matters: canonicalName; exact scope remains in the canonical record.}\par
{\footnotesize \href{../../../../continuum/application/application-component/33b8421ffc444b78b2d059201a58a688.md}{Canonical record: 33b8421ffc444b78b2d059201a58a688}; full qualifications and provenance.}\par
\section*{IV.3.7 Harmonia.Omilia.Server}
\phantomsection
\label{navigation:IV.3.7}
\addcontentsline{toc}{section}{IV.3.7 Harmonia.Omilia.Server}

\subsection*{IV.3.7.1 Application Overview}
\phantomsection
\label{navigation:IV.3.7.1}
\addcontentsline{toc}{subsection}{IV.3.7.1 Application Overview}

\subsubsection*{IV.3.7.1.1 Harmonia.Omilia.Server}
\phantomsection
\label{navigation:IV.3.7.1.1}
\addcontentsline{toc}{subsubsection}{IV.3.7.1.1 Harmonia.Omilia.Server}

\label{element:369907025d06492bb414a72cea834a77}
\index{Harmonia.Omilia.Server (IV.3.7.1.1)}
{\small ApplicationComponent · Candidate — accepted Step-B baseline; final Solution Architecture unapproved · current}\par

\begin{itemize}


\item \textbf{Role/responsibility:} owns/performs underlying Room/Space, user, membership,
  dialogue/message, collaboration-event and applicable collaboration-media
  operations. Exposes the \textbf{Application Services API} for Agora operation requests
  and dialogue/activity delivery from Agora-managed Rooms/Spaces.

\item \textbf{Boundary:} operational worker, not owner of Harmonia collaboration semantics,
  membership intent or originating clinical/operational information. Agora owns
  native collaboration information; AX-05 durable-state responsibility remains.

\item \textbf{Provenance/standing:} derived from Matrix Synapse; U16 adjudicated candidate.
  API technology, modifications and persistence mechanics are deferred, not
  inferred from historical Synapse deployment/configuration.

\end{itemize}

{\small Recorded unresolved matters: canonicalName; exact scope remains in the canonical record.}\par
{\footnotesize \href{../../../../continuum/application/application-component/369907025d06492bb414a72cea834a77.md}{Canonical record: 369907025d06492bb414a72cea834a77}; full qualifications and provenance.}\par
\subsection*{IV.3.7.2 Interfaces}
\phantomsection
\label{navigation:IV.3.7.2}
\addcontentsline{toc}{subsection}{IV.3.7.2 Interfaces}

\subsubsection*{IV.3.7.2.1 Omilia Application Services API}
\phantomsection
\label{navigation:IV.3.7.2.1}
\addcontentsline{toc}{subsubsection}{IV.3.7.2.1 Omilia Application Services API}

\label{element:49c52c0466bc4760a11e0ca6fe8142c0}
\index{Omilia Application Services API (IV.3.7.2.1)}
{\small ApplicationInterface · Candidate — adjudicated Domain05 Step-B baseline; final Solution Architecture unapproved · current}\par

Omilia.Server exposes an \textbf{Application Services API} used by Agora both to
request collaboration operations and to receive managed-Room dialogue/activity.



\begin{quote}\small\ttfamily
flowchart LR\\
    Agora["Harmonia.Agora: collaboration control"] -->|"Application Services API: operation requests"| Server["Harmonia.Omilia.Server: collaboration operations"]\\
    Server -->|"Application Services API: managed-Room dialogue / activity"| Agora\\
    Client["Harmonia.Omilia.Client: collaboration experience"] <-->|"Matrix protocol"| Server\\

\end{quote}


Agora decides to create/change a managed Room/Space, add/remove a participant or
publish information; Omilia.Server performs the operation. Room activity returns
through the same architectural boundary. The exact API contract, technology and
mapping to upstream Matrix interfaces remain deferred. Historical Admin,
Client-Server and AS-transaction endpoints do not define this new contract.



Workflow/execution collaboration continues through Petasos with \textbf{no direct
Agora-to-Ponos dependency}. Matrix DTO encapsulation, Themis default-deny
governance and zero-PHI Room aliases/names/topics remain repository constraints.
Membership or receipt of a notification does not confer information-access rights.



\href{\detokenize{../../../../../docs/markdown/05-information-systems-architecture/sources.md#u19}}{U19} establishes the administrative interaction boundary:
Administrator → Omilia → Agora → Administrative Function → specialised
Administrative Ergon → governed subsystem operation. Stopped-Twin resumption
uses this pattern, with normal authorisation, Context, execution and Kleio audit
controls; Agora and Omilia do not directly resume or manipulate the Twin.
The arrows describe semantic progression and retain Petasos execution handoff
and Agora's no-direct-Ponos-dependency rule. Iris provides a separate path through
Iris-BackEnd and does \textbf{not} directly interact with Agora or Omilia. Broader
Administrative Function reconciliation is \href{\detokenize{../../../../../docs/markdown/05-information-systems-architecture/gaps.md#administrative-function-finding}}{deferred}.



Identity mapping/provisioning is collaboration integration, not authentication
authority. Existing isolation constraints are not relaxed; detailed security,
privacy and external/federated interactions remain deferred.


{\small Recorded unresolved matters: canonicalName; exact scope remains in the canonical record.}\par
{\footnotesize \href{../../../../continuum/application/application-interface/49c52c0466bc4760a11e0ca6fe8142c0.md}{Canonical record: 49c52c0466bc4760a11e0ca6fe8142c0}; full qualifications and provenance.}\par
\section*{IV.3.8 Harmonia.Sylloge}
\phantomsection
\label{navigation:IV.3.8}
\addcontentsline{toc}{section}{IV.3.8 Harmonia.Sylloge}

\subsection*{IV.3.8.1 Application Overview}
\phantomsection
\label{navigation:IV.3.8.1}
\addcontentsline{toc}{subsection}{IV.3.8.1 Application Overview}

\subsubsection*{IV.3.8.1.1 Harmonia.Sylloge}
\phantomsection
\label{navigation:IV.3.8.1.1}
\addcontentsline{toc}{subsubsection}{IV.3.8.1.1 Harmonia.Sylloge}

\label{element:feea742463b84021b66853d287d8916f}
\index{Harmonia.Sylloge (IV.3.8.1.1)}
{\small ApplicationComponent · Candidate — accepted Step-B baseline; final Solution Architecture unapproved · current}\par

\begin{itemize}


\item \textbf{Role/responsibility:} \textbf{Binary Content Management subsystem} responsible
  for governed association, resolution, retrieval and assembly of binary content
  associated with Harmonia-managed information. \hyperref[element:feea742463b84021b66853d287d8916f]{U17 canonical view}.

\item \textbf{Feature/information:} F31/D15; binary-content metadata, Information Unit
  associations, location/resolution knowledge, retrieval, on-demand assembly from
  one or more payloads and knowledge of Harmonia-managed versus externally
  managed content. Logical information and binary payload remain distinct.

\item \textbf{Custody/persistence:} physical custody may remain external; management does
  not require copying every payload into Harmonia. Where applicable, Mnemosyne
  provides Harmonia persistence services without acquiring Binary Content
  Management. Sylloge is not a generic persistence manager.

\item \textbf{Collaborations/exclusions:} Mneme remains Context and Pragma Manager with its
  established service boundary; applicable processing components may obtain/resolve
  content through Sylloge. Agora retains native collaboration media and the
  record/reference/activity concerning formal binary content; Sylloge retains
  that content's Binary Content Management. Omilia acquires no such responsibility.
  Source information/domain authority and external physical custody are preserved.

\item \textbf{Standing:} U17 adjudicated Candidate Architecture; former A2 responsibility
  question CDG-25 closed. \hyperref[element:feea742463b84021b66853d287d8916f]{Storage, representations, APIs, algorithms and lifecycle
  mechanisms} remain
  non-blocking later design. No module, dependency or deployment allocation.

\end{itemize}

{\small Recorded unresolved matters: canonicalName; exact scope remains in the canonical record.}\par
{\footnotesize \href{../../../../continuum/application/application-component/feea742463b84021b66853d287d8916f.md}{Canonical record: feea742463b84021b66853d287d8916f}; full qualifications and provenance.}\par
\section*{IV.3.9 Iris}
\phantomsection
\label{navigation:IV.3.9}
\addcontentsline{toc}{section}{IV.3.9 Iris}

\subsection*{IV.3.9.1 Application Overview}
\phantomsection
\label{navigation:IV.3.9.1}
\addcontentsline{toc}{subsection}{IV.3.9.1 Application Overview}

\subsubsection*{IV.3.9.1.1 Iris}
\phantomsection
\label{navigation:IV.3.9.1.1}
\addcontentsline{toc}{subsubsection}{IV.3.9.1.1 Iris}

\label{element:e927c1af96c64fcd979be659bee2cbe3}
\index{Iris (IV.3.9.1.1)}
{\small ApplicationComponent · Candidate — accepted Step-B baseline; final Solution Architecture unapproved · current}\par

\begin{itemize}


\item \textbf{Role/responsibility:} contextual human interaction, presentation, input and
  governed requests. \href{\detokenize{../../../../../docs/markdown/05-information-systems-architecture/sources.md#a02}}{A02}, \href{\detokenize{../../../../../docs/markdown/05-information-systems-architecture/sources.md#l11}}{L11/L21/L23}.

\item \textbf{Features/contained constructs:} clinical, operations and administration
  applications (C16), mediation (C17), shared shell/design system and role-tailored
  workspaces. These are documented presentation responsibilities, not product modules
  promoted solely by their directory names.

\item \textbf{Information/persistence:} displays contextual resource and operational views,
  captures user intent/assertions, transient browser selections/filters/state.
  No authoritative clinical/registry or persistence ownership; authority of human
  input comes from the actor/process/source, not Iris.

\item \textbf{Collaboration/exclusions:} uses defined governed information/activity interfaces
  and Themis. No raw-cache/database alternative, JPA/SQL, backend registry validation,
  autonomous workflow engine or clinical decision authority.
  Iris does not directly interact with Agora or Omilia. Administrative intent
  passes through Iris-BackEnd to an Administrative Function and specialised
  Administrative Ergon, not direct subsystem manipulation (U19).

\item \textbf{Standing:} established high-level responsibility; candidate detailed views,
  CDG-11/24/26/28. Diagnostic zero-PHI controls do not turn clinical presentation
  into a technical-identifiers-only application.

\end{itemize}

{\small Recorded unresolved matters: canonicalName; exact scope remains in the canonical record.}\par
{\footnotesize \href{../../../../continuum/application/application-component/e927c1af96c64fcd979be659bee2cbe3.md}{Canonical record: e927c1af96c64fcd979be659bee2cbe3}; full qualifications and provenance.}\par
\subsection*{IV.3.9.2 Component Architecture}
\phantomsection
\label{navigation:IV.3.9.2}
\addcontentsline{toc}{subsection}{IV.3.9.2 Component Architecture}

\subsubsection*{IV.3.9.2.1 Iris backend-for-frontend mediation}
\phantomsection
\label{navigation:IV.3.9.2.1}
\addcontentsline{toc}{subsubsection}{IV.3.9.2.1 Iris backend-for-frontend mediation}

\label{element:37b82d161c5640298e951acc361b338a}
\index{Iris backend-for-frontend mediation (IV.3.9.2.1)}
{\small ApplicationComponent · Candidate — accepted Step-B baseline; final Solution Architecture unapproved · current}\par

\begin{itemize}


\item \textbf{Role/features:} decoupled presentation service mediation and clinical versus
  operations request separation; telemetry aggregation and provider contracts.
  \href{\detokenize{../../../../../docs/markdown/05-information-systems-architecture/sources.md#l11}}{L11/L21/L23}.

\item \textbf{Information/contained constructs:} projected resources, task/sequence views,
  subsystem status, dependency/queue/gateway observations and alert diagnostics;
  operations aggregator and health-provider support, not a new information authority.

\item \textbf{Persistence/collaboration:} uses Mneme-facing governed access, Themis and
  approved execution requests; backend/presentation state is not durable truth.
  Documented raw Hot Rod/cache paths remain conflicting (CDG-11).
  \textbf{Iris-BackEnd} mediates administrative intent into governed Administrative
  Function/Administrative Ergon execution. It does not directly resume a stopped
  Twin or provide a privileged subsystem backdoor (U19).

\item \textbf{Exclusions/standing:} no direct JPA/PostgreSQL or server-side registry governance.
  Candidate mediation component; ports/server frameworks/provider class catalogues
  are support details, and telemetry persistence/provenance remains CDG-28.

\end{itemize}

{\small Recorded unresolved matters: canonicalName; exact scope remains in the canonical record.}\par
{\footnotesize \href{../../../../continuum/application/application-component/37b82d161c5640298e951acc361b338a.md}{Canonical record: 37b82d161c5640298e951acc361b338a}; full qualifications and provenance.}\par
\section*{IV.3.10 Kleio}
\phantomsection
\label{navigation:IV.3.10}
\addcontentsline{toc}{section}{IV.3.10 Kleio}

\subsection*{IV.3.10.1 Application Overview}
\phantomsection
\label{navigation:IV.3.10.1}
\addcontentsline{toc}{subsection}{IV.3.10.1 Application Overview}

\subsubsection*{IV.3.10.1.1 Kleio}
\phantomsection
\label{navigation:IV.3.10.1.1}
\addcontentsline{toc}{subsubsection}{IV.3.10.1.1 Kleio}

\label{element:fefda110e2274884ad58906733606efd}
\index{Kleio (IV.3.10.1.1)}
{\small ApplicationComponent · Candidate — accepted Step-B baseline; final Solution Architecture unapproved · current}\par

\begin{itemize}


\item \textbf{Role/responsibility:} \textbf{Harmonia's authoritative Audit Information Management
  component.} Kleio owns authoritative Audit Capture, creation/management of
  authoritative audit information, canonical FHIR \textbf{AuditEvent} representation,
  preservation/retrieval of audit history, semantically significant user/system
  audit activity and provision/preparation of historical audit context for
  operational and assurance use. \href{\detokenize{../../../../../docs/markdown/05-information-systems-architecture/sources.md#u07}}{U07}, \hyperref[element:21d6864657964ba5b03e3f6b5eb29cd5]{AX-08}.

\item \textbf{Audit Capture:} clients report semantically significant auditable occurrences
  and necessary context through a \textbf{Kleio-provided contract/library/API}.
  \textbf{Clients SHALL NOT be required to construct FHIR AuditEvent resources.} Kleio
  constructs and manages the authoritative representation; client API and mapping
  details remain unresolved.

\item \textbf{Information:} FHIR AuditEvent is the authoritative audit-event representation,
  without requiring every internal Kleio structure to be an AuditEvent. Indexes,
  correlation structures and prepared historical context may support management.
  \textbf{Authoritative audit evidence SHALL remain distinguishable from derived,
  correlated, contextualised or prepared historical views.} Existing meaningful
  provenance responsibilities remain; AuditEvent does not redefine all evidence.

\item \textbf{User activity:} semantically significant access/view, governed search,
  create/change/delete, submit/approve/reject, disclosure/export, operational action
  and relevant denied activity may require capture. These are examples, not a
  universal event taxonomy; “delete” grants no physical managed-information DELETE.
  Generic UI telemetry/clickstream is outside Kleio; UI security assurance and
  governance remain \textbf{Somebody Else's Problem (SEP)}.

\item \textbf{Collaboration/exclusions:} \textbf{Audit is not Governance. Audit is not Assurance.}
  Kleio may consume Journey information and Dokimasia outputs and provide
  historical audit context back to Dokimasia, without evaluating assurance or
  governing the subject. Dokimasia's governed record/result remains authoritative
  evidence of its assurance activity; its \href{\detokenize{../../../../../docs/markdown/05-information-systems-architecture/component-register.md#dokimasia-terminal-boundaries}}{terminal rules}
  prohibit a separate own-execution AuditEvent and recursive assurance. Consumption
  of its outputs does not circumvent those rules.

\end{itemize}






\par\medskip\textbf{Mandatory capture authority and the audited Identity}\par


\begin{quote}


\textbf{Subject of audit evidence ≠ authority performing audit capture.}



\end{quote}


Kleio performs mandatory authoritative audit capture under \textbf{intrinsic platform
authority} associated with its established architectural audit responsibility.
The originating/acting Identity remains evidence recorded by Kleio; it is not
the authority under which the mandatory AuditEvent is persisted. Kleio does not
borrow that Identity's PEP-Criteria merely to record its activity, including
activity denied to that Identity.



Identity-derived authority granted to human/system Identities remains resolved
and enforced through Themis. Kleio capture executes within established security
context and remains governed by Harmonia security policy, with Themis retaining
policy evaluation ownership. Mandatory capture is neither a security bypass nor
a policy exception. \href{\detokenize{../../../../governance/architectural-axioms.md#kleio-mandatory-capture-authority}}{AX-07}
incorporates this adjudication.



Where authoritative audit capture is required for a governed operation and
Kleio cannot exercise that platform authority sufficiently to capture the
required evidence, the governed operation \textbf{SHALL fail closed}. Durable target
availability remains necessary; authority alone establishes no successful capture.



This authority applies specifically to Kleio mandatory authoritative audit
capture. It grants no equivalent authority to other components or other Kleio
activities. Any proposed equivalent responsibility requires separate architecture
adjudication; representation, policy/enforcement mechanics and transaction
coordination remain \href{\detokenize{../../../../../docs/markdown/05-information-systems-architecture/gaps.md#kleio-detailed-contracts}}{CDG-43}.







\par\medskip\textbf{Immutable history and longitudinal evidence}\par


Authoritative AuditEvents are immutable. An existing AuditEvent SHALL NOT be
corrected, amended, overwritten or silently replaced. Correction or qualification
SHALL create a \textbf{new AuditEvent}, preserve the earlier event unchanged and refer
back to it as appropriate. The earlier assertion remains part of historical record.



Journey is the governed representation of observed Harmonia activity/progression;
AuditEvent is the authoritative audit representation of a significant auditable
occurrence. Journey is an important audit source, \textbf{not the audit boundary}.
A Journey may relate to many AuditEvents; an AuditEvent may exist without Journey.
\textbf{Journey participation SHALL NOT be a prerequisite for Audit Capture.}



\begin{quote}


\textbf{Reference by default; replicate only where necessary to preserve longitudinal
audit evidence.}



\end{quote}


Kleio SHALL NOT duplicate authoritative source information merely because it is
audit-relevant. Where historical state can be reliably resolved longitudinally,
Kleio references it with sufficient \textbf{identity, version, provenance and temporal
detail}. Version identity alone is not assumed to supply all temporal detail.
Only where the authoritative source cannot guarantee retrieval of the information
as it existed at the audited point in time may Kleio clone/replicate the necessary
source information within or in association with the authoritative AuditEvent to
preserve required longitudinal evidence. Detailed FHIR attachment/reference
mechanisms remain unresolved.







\par\medskip\textbf{Auditability and R1.x/R2.x availability}\par


\begin{quote}


\textbf{Harmonia SHALL NOT perform an auditable activity unless the required audit
information can be authoritatively and durably captured.}



\textbf{If Kleio cannot authoritatively capture required audit information, Harmonia
is DOWN.}



\end{quote}


Audit Capture is not best-effort and SHALL NOT be assumed reconstructable later
from logs, Journey information or incidental operational state. Transaction,
retry/failover and health/readiness mechanisms are deliberately unresolved.



Kleio retains audit application ownership; Mneme provides governed Context/Pragma access
and Mnemosyne \textbf{Harmonia Operations \& Assurance} establishes durable audit state.
For R1.x/R2.x, multiple \textbf{independent authoritative audit persistence targets}
through Mneme/Mnemosyne are permitted. At least \textbf{one configured target capable
of durably establishing the required audit information} is sufficient for
continued Harmonia operation. If none can establish it, Harmonia is unavailable
and fails closed. Target count, selection and whether one/some/all receive an
event are unresolved; targets SHALL NOT be called replicas without established
replication semantics.



\begin{quote}


\textbf{For R1.x/R2.x, audit data harmonisation between independent Mnemosyne audit
persistence targets is outside Harmonia scope — Somebody Else's Problem (SEP).}



\end{quote}


R1.x/R2.x guarantees neither identical histories, complete copies of all
AuditEvents, automatic convergence, reconciliation after disconnection nor unified
cross-target longitudinal history. Different targets may expose different histories;
an immutable AuditEvent remains authoritative without its target holding the whole
corpus. No future harmonisation/convergence mechanism is selected.



\begin{itemize}


\item \textbf{Standing:} U07 resolves responsibility, representation, immutability,
  longitudinal evidence, availability and release scope; \href{\detokenize{../../../../../docs/markdown/05-information-systems-architecture/gaps.md#kleio-detailed-contracts}}{CDG-43}
  retains detailed capture/query/correlation/retention/processing and persistence
  mechanisms. Historical FHIR/operational storage structures remain provenance,
  without competing with the established O\&A semantic alignment.

\end{itemize}

{\small Recorded unresolved matters: canonicalName; exact scope remains in the canonical record.}\par
{\footnotesize \href{../../../../continuum/application/application-component/fefda110e2274884ad58906733606efd.md}{Canonical record: fefda110e2274884ad58906733606efd}; full qualifications and provenance.}\par
\section*{IV.3.11 Mneme}
\phantomsection
\label{navigation:IV.3.11}
\addcontentsline{toc}{section}{IV.3.11 Mneme}

\subsection*{IV.3.11.1 Application Overview}
\phantomsection
\label{navigation:IV.3.11.1}
\addcontentsline{toc}{subsection}{IV.3.11.1 Application Overview}

\subsubsection*{IV.3.11.1.1 Mneme}
\phantomsection
\label{navigation:IV.3.11.1.1}
\addcontentsline{toc}{subsubsection}{IV.3.11.1.1 Mneme}

\label{element:1af8a56358f7429fabfb57be66542a99}
\index{Mneme (IV.3.11.1.1)}
{\small ApplicationComponent · Candidate — accepted Step-B baseline; final Solution Architecture unapproved · current}\par

\begin{itemize}


\item \textbf{Role/responsibility:} \textbf{Harmonia's Context and Pragma Manager}. Mneme manages
  the information context of activity through Pragma/Context, providing controlled,
  concurrent and seamless access to existing values and management of new Content.
  \href{\detokenize{../../../../../docs/markdown/05-information-systems-architecture/sources.md#u10}}{U10} refines AX-05/11 active representation/coordination.

\item \textbf{Principal exposed services:} \textbf{Context Load, Context Save, Adhoc Resource
  Access, Adhoc Search} — four services. \hyperref[element:49c9eb2bef0c47b8b5328621d5fc5479]{Lifecycle}
  distinguishes Mneme materialisation/save from Ergo Loader/Unloader boundaries.
  Likely resource sets are non-exhaustive; bounded runtime acquisition is supported.

\item \textbf{Internal responsibilities:} Pragma Management C28 retains PragmaFactory and
  Publishing/Subscription/Pattern semantic eligibility/selection. Context management
  composes authoritative resource values into a versioned temporal snapshot;
  Resource Harmonisation Manager C40 manages exceptional delivery conditions.
  C11 supporting contracts do not add exposed services.

\item \textbf{Information:} Context D26 is the working value environment; Pragma C07 is
  the input/output envelope. Changed Context values and new Content expressed as
  Pragma Output D28 remain distinct. Mneme manages their subsequent consequences.

\item \textbf{Persistence:} active state is reconstructable; Mnemosyne alone establishes
  durable state and authoritatively reports persistence outcomes. Working changes,
  cache retention/replication and output production confer no durable authority.
  Context/resource versions and active generation remain distinct.

\item \textbf{Collaboration/exclusions:} Ergo framework anticipates/loads; Ergon primarily
  acts on Context values and may acquire Adhoc information. Resource harmonisation
  uses new governed Harmonisation Ergon work through common Ponos execution, not
  originating-Ergon conflict resolution. Adhoc Search uses Mnemosyne Search Management.
  Petasos retains handoff; Themis owns enforcement. Mneme is not a generic repository,
  execution framework, durable version authority or raw database/cache access path.

\item \textbf{Standing:} U10 adjudicates Context/Pragma responsibility; detailed representation,
  concurrency/access/Output and harmonisation contracts remain CDG-45/46; older
  cache/write-behind and access variants remain CDG-10/11/16/27/31.

\end{itemize}

{\small Recorded unresolved matters: canonicalName; exact scope remains in the canonical record.}\par
{\footnotesize \href{../../../../continuum/application/application-component/1af8a56358f7429fabfb57be66542a99.md}{Canonical record: 1af8a56358f7429fabfb57be66542a99}; full qualifications and provenance.}\par
\subsection*{IV.3.11.2 Component Architecture}
\phantomsection
\label{navigation:IV.3.11.2}
\addcontentsline{toc}{subsection}{IV.3.11.2 Component Architecture}

\subsubsection*{IV.3.11.2.1 PragmaFactory}
\phantomsection
\label{navigation:IV.3.11.2.1}
\addcontentsline{toc}{subsubsection}{IV.3.11.2.1 PragmaFactory}

\label{element:75e77d8dfd1d4870aa6151485d28d1dd}
\index{PragmaFactory (IV.3.11.2.1)}
{\small ApplicationComponent · Candidate — accepted Step-B baseline; final Solution Architecture unapproved · current}\par

\begin{itemize}


\item \textbf{Role:} candidate logical application embodiment associated with the
  ActionableTaskArchetype meta-concept, within Mneme Pragma Management C28.
  \href{\detokenize{../../../../../docs/markdown/05-information-systems-architecture/sources.md#u02}}{U01/U02}. The meta-concept is not itself an instantiated
  Information Asset; the factory is not a new upstream Information Resource.

\item \textbf{Responsibilities/features:} creates/initialises appropriate Pragma
  representations. Detailed instantiation/selection/validation, parameters and
  policy contracts remain \textbf{unknown}; the basic responsibility does not finalise
  a factory feature contract.

\item \textbf{Contained constructs/information/persistence:} subcomponents, realised
  definition Data Objects, inputs/outputs and persistence responsibility unknown.

\item \textbf{Collaborations/exclusions:} management allocation is established by U02;
  detailed callers and contracts remain unresolved. It must not be equated to a
  TaskSequence seeder or PraxisDefinition by inference.

\item \textbf{Standing:} explicitly instructed candidate solution meaning, CDG-04; not
  invented from a matching class and not rejected for missing implementation.

\end{itemize}

{\small Recorded unresolved matters: canonicalName; exact scope remains in the canonical record.}\par
{\footnotesize \href{../../../../continuum/application/application-component/75e77d8dfd1d4870aa6151485d28d1dd.md}{Canonical record: 75e77d8dfd1d4870aa6151485d28d1dd}; full qualifications and provenance.}\par
\subsection*{IV.3.11.3 Information Architecture}
\phantomsection
\label{navigation:IV.3.11.3}
\addcontentsline{toc}{subsection}{IV.3.11.3 Information Architecture}

\subsubsection*{IV.3.11.3.1 Pragma Topic}
\phantomsection
\label{navigation:IV.3.11.3.1}
\addcontentsline{toc}{subsubsection}{IV.3.11.3.1 Pragma Topic}

\label{element:27ad143e7c784e378556af8a57c438f7}
\index{Pragma Topic (IV.3.11.3.1)}
{\small DataObject · Bounded adjudicated application information meaning; Domain05 remains the accepted candidate Step-B baseline, not final Solution Architecture. Detailed representation is deferred; only this defined logical object is projected. · current}\par

\textbf{Mneme manages Pragma.} Context is the value-based working environment;
Pragma is the task/activity information envelope through which inputs and outputs
are expressed. They are related and distinct.



\textbf{Pragma realises \hyperref[element:c34b4929f71946708d37dbe6c1da920a]{Metabolē}}
for governed Harmonia information/activity progression. This is architectural
Realisation, not synonymy or inheritance from Diadromí. Mneme retains Pragma
management; the conceptual construct introduces no common manager, persistence
requirement or shared lifecycle with Pylai.Diadromí.



The conceptual
\hyperref[element:c34b4929f71946708d37dbe6c1da920a]{ActivityContext}
classifies the broad activity and is distinct from request action, this loaded
Context and Themis PEP-Context. Their detailed physical mapping remains deferred.



An Ergon may read existing Context values, act upon/change working Context
values, and generate \textbf{new Content}. New Content is expressed as \textbf{Pragma
Output}, rather than merely an incidental Context mutation:



\begin{quote}\small\ttfamily
Existing Context values → Ergon → changed working Context values\\
Context + bounded business logic → Ergon → Pragma Output → new Content\\
                  Context Unloader returns both, distinctly\\
                                    ↓\\
                 Mneme manages subsequent information consequences\\

\end{quote}


Changed Context values and newly generated Pragma Outputs SHALL remain
semantically distinguishable. A working change or produced output does not
itself establish authoritative durable state. Task/undertaking/outcome ownership
and output rules remain unchanged.



Ergon execution may produce one or more outputs, whose processing may result in
multiple Pragma. \textbf{Only a Pragma playing the Pathway-continuation role identifies
information for continuing the current Praxis along its Pathway}; other resulting
Pragma are independently available for governed Pub/Sub, another Praxis, other
processing or no further processing. Praxis defines continuation semantics;
the execution framework performs applicable output processing and routing,
without making Praxis an orchestrator/choreographer or transferring Mneme's
information management. See the \hyperref[element:27ad143e7c784e378556af8a57c438f7]{shared semantic view}
and \href{\detokenize{../../../../../docs/markdown/05-information-systems-architecture/design-notes/DesignNote-Pragma-Outputs-and-Pathway-Continuation.md}}{DesignNote}.
The continuation role creates no subtype. Whether a new Pragma is constructed,
an output Pragma is designated, or another governed representation is used remains
\href{\detokenize{../../../../../docs/markdown/05-information-systems-architecture/gaps.md#continuation-representation}}{explicitly deferred}. Exact Output
representation, routing mechanisms and Pragma/Information Unit mappings remain open.



Mneme may manage richer resulting Pragma relationships and provenance than the
originating Praxis's Journey represents. Continuation lineage supports meaningful
Journey progression without making Journey a complete provenance graph, execution
log, audit trail or queue history; \hyperref[element:359b901df27f4b43b5a2ff14a8484e34]{Journey's boundary}
and Energeia ownership remain intact.



Source standing: Bounded adjudicated application information meaning; Domain05 remains the accepted candidate Step-B baseline, not final Solution Architecture. Detailed representation is deferred; only this defined logical object is projected.


{\footnotesize \href{../../../../continuum/application/data-object/27ad143e7c784e378556af8a57c438f7.md}{Canonical record: 27ad143e7c784e378556af8a57c438f7}; full qualifications and provenance.}\par
\subsubsection*{IV.3.11.3.2 Context Topic}
\phantomsection
\label{navigation:IV.3.11.3.2}
\addcontentsline{toc}{subsubsection}{IV.3.11.3.2 Context Topic}

\label{element:76ded4866bf443dabc0849d68b215690}
\index{Context Topic (IV.3.11.3.2)}
{\small DataObject · Bounded adjudicated application information meaning; Domain05 remains the accepted candidate Step-B baseline, not final Solution Architecture. Detailed representation is deferred; only this defined logical object is projected. · current}\par

\begin{quote}


\textbf{Context is a versioned, temporal load-snapshot composed from authoritative
resource values materialised for a Harmonia activity's anticipated information
requirements.}



\end{quote}


Context contains \textbf{values, not merely references}. Each value represents the
defined authoritative resource state/version actually materialised when it is
incorporated. Loaded values are composed into the Context bucket:



\begin{quote}\small\ttfamily
Context — identity / Context version / temporal load basis\\
  └── Resource Bucket                         [Composition]\\
        ├── Resource Value A @ authoritative state/version\\
        ├── Resource Value B @ authoritative state/version\\
        └── Resource Value C @ authoritative state/version\\

\end{quote}


\textbf{Composition} concerns loaded values belonging to the snapshot. The Context
does not existentially own the underlying authoritative resources. This is
distinct from the separately established Aggregation relationships for Praxis,
Journey and PragmaFactory.



Context version and constituent resource authoritative state/version SHALL
remain distinct. Neither is an alias for a valid active-state generation or
security PEP-Context. Materialisation does \textbf{not} imply a globally
transactionally consistent distributed snapshot or a universal frozen instant.
Context is the Ergon's managed working information environment, not a closed
universe of all information it may access or produce.



Authoritative constituent resource state may progress while an Ergon operates.
Mneme needs sufficient state/version information to recognise delivery that cannot
safely proceed normally; generic conflict handling remains outside the originating
Ergon. No concurrency or conflict-detection algorithm is established here.



This activity Context does not redefine Domain04's Information Unit
\href{\detokenize{../../../../../docs/markdown/04-information-architecture/metamodel/information-architecture-metamodel.md#65-context}}{Management Context}.
Their detailed mapping, physical bucket representation, Context identity/version
mechanisms and concurrency algorithms remain unresolved under AX-17.



Source standing: Bounded adjudicated application information meaning; Domain05 remains the accepted candidate Step-B baseline, not final Solution Architecture. Detailed representation is deferred; only this defined logical object is projected.


{\footnotesize \href{../../../../continuum/application/data-object/76ded4866bf443dabc0849d68b215690.md}{Canonical record: 76ded4866bf443dabc0849d68b215690}; full qualifications and provenance.}\par
\subsection*{IV.3.11.4 Services}
\phantomsection
\label{navigation:IV.3.11.4}
\addcontentsline{toc}{subsection}{IV.3.11.4 Services}

\subsubsection*{IV.3.11.4.1 Context Load}
\phantomsection
\label{navigation:IV.3.11.4.1}
\addcontentsline{toc}{subsubsection}{IV.3.11.4.1 Context Load}

\label{element:1ccfe420541945bb9c0756a949799d95}
\index{Context Load (IV.3.11.4.1)}
{\small ApplicationService · Candidate — adjudicated Domain05 Step-B baseline; final Solution Architecture unapproved · current}\par

Materialise a managed Context from authoritative resource values for anticipated activity requirements.



Mneme exposes this principal service. Resource Harmonisation Manager remains internal, and supporting contracts are not additional exposed services. The Context/Pragma and authoritative-state boundaries apply.


{\small Recorded unresolved matters: canonicalName; exact scope remains in the canonical record.}\par
{\footnotesize \href{../../../../continuum/application/application-service/1ccfe420541945bb9c0756a949799d95.md}{Canonical record: 1ccfe420541945bb9c0756a949799d95}; full qualifications and provenance.}\par
\subsubsection*{IV.3.11.4.2 Context Save}
\phantomsection
\label{navigation:IV.3.11.4.2}
\addcontentsline{toc}{subsubsection}{IV.3.11.4.2 Context Save}

\label{element:24d4fc42d9524d53a747f0550833f1a4}
\index{Context Save (IV.3.11.4.2)}
{\small ApplicationService · Candidate — adjudicated Domain05 Step-B baseline; final Solution Architecture unapproved · current}\par

Return resulting Context state to Mneme management for appropriate subsequent information handling, alongside distinguishable Pragma Outputs.



Mneme exposes this principal service. Resource Harmonisation Manager remains internal, and supporting contracts are not additional exposed services. The Context/Pragma and authoritative-state boundaries apply.


{\small Recorded unresolved matters: canonicalName; exact scope remains in the canonical record.}\par
{\footnotesize \href{../../../../continuum/application/application-service/24d4fc42d9524d53a747f0550833f1a4.md}{Canonical record: 24d4fc42d9524d53a747f0550833f1a4}; full qualifications and provenance.}\par
\subsubsection*{IV.3.11.4.3 Adhoc Resource Access}
\phantomsection
\label{navigation:IV.3.11.4.3}
\addcontentsline{toc}{subsubsection}{IV.3.11.4.3 Adhoc Resource Access}

\label{element:410f2f55bbda481b973c6c67b924f0be}
\index{Adhoc Resource Access (IV.3.11.4.3)}
{\small ApplicationService · Candidate — adjudicated Domain05 Step-B baseline; final Solution Architecture unapproved · current}\par

Governed runtime acquisition of resource information outside the anticipated loaded Context where the bounded activity requires it.



Mneme exposes this principal service. Resource Harmonisation Manager remains internal, and supporting contracts are not additional exposed services. The Context/Pragma and authoritative-state boundaries apply.


{\small Recorded unresolved matters: canonicalName; exact scope remains in the canonical record.}\par
{\footnotesize \href{../../../../continuum/application/application-service/410f2f55bbda481b973c6c67b924f0be.md}{Canonical record: 410f2f55bbda481b973c6c67b924f0be}; full qualifications and provenance.}\par
\subsubsection*{IV.3.11.4.4 Adhoc Search}
\phantomsection
\label{navigation:IV.3.11.4.4}
\addcontentsline{toc}{subsubsection}{IV.3.11.4.4 Adhoc Search}

\label{element:02bf54d699b44ed0b26fa5b8125f0f01}
\index{Adhoc Search (IV.3.11.4.4)}
{\small ApplicationService · Candidate — adjudicated Domain05 Step-B baseline; final Solution Architecture unapproved · current}\par

Activity/context-facing governed search, using Mnemosyne's underlying authoritative persistence Search Management.



Mneme exposes this principal service. Resource Harmonisation Manager remains internal, and supporting contracts are not additional exposed services. The Context/Pragma and authoritative-state boundaries apply.


{\small Recorded unresolved matters: canonicalName; exact scope remains in the canonical record.}\par
{\footnotesize \href{../../../../continuum/application/application-service/02bf54d699b44ed0b26fa5b8125f0f01.md}{Canonical record: 02bf54d699b44ed0b26fa5b8125f0f01}; full qualifications and provenance.}\par
\subsection*{IV.3.11.5 Internal Functions}
\phantomsection
\label{navigation:IV.3.11.5}
\addcontentsline{toc}{subsection}{IV.3.11.5 Internal Functions}

\subsubsection*{IV.3.11.5.1 Pragma Management}
\phantomsection
\label{navigation:IV.3.11.5.1}
\addcontentsline{toc}{subsubsection}{IV.3.11.5.1 Pragma Management}

\label{element:ab0cdaa7ed4149ee856216efa5a265c4}
\index{Pragma Management (IV.3.11.5.1)}
{\small ApplicationFunction · Candidate — accepted Step-B baseline; final Solution Architecture unapproved · current}\par

\begin{itemize}


\item \textbf{Role/responsibility:} Mneme-associated information-management responsibility
  area with logical Aggregation of PragmaFactory C08 and Publishing C36,
  Subscription C37, Pattern C38 responsibilities, managing Pragma C07.
  \href{\detokenize{../../../../../docs/markdown/05-information-systems-architecture/sources.md#u05}}{U02/U03/U05}. Aggregation/management identifies responsibility,
  not Composition or a required process/service.

\item \textbf{Features/information:} creation/initialisation through PragmaFactory and
  governed management of Pragma representations and common process-management,
  assurance, reporting and platform integration associated with that information.
  Publishing/Subscription definitions, Pattern semantics/evaluation, matching,
  selection, semantic eligibility and eligible recipient/consumer relationships
  belong here; concrete persistent destination/handoff mechanics remain Petasos.
  Pragma expresses activity inputs/outputs; new Content is \textbf{Pragma Output}, distinct
  from changed Context working values. Mneme manages its subsequent consequences
  under \hyperref[element:27ad143e7c784e378556af8a57c438f7]{U10}.
  U18 distinguishes independently resulting Pragma from the Pathway-continuation
  role without transferring Praxis semantics or execution/routing to Mneme.
  Stimulus/Effector and candidate Dokimasia Pragma semantic roles
  reuse these facilities instead of requiring specialist equivalents.

\item \textbf{Persistence/collaboration:} active use/coordination through Mneme; authoritative
  durable establishment through Mnemosyne. Ponos executes, Petasos provides
  messaging/recoverable handoff and Ergon supplies bounded activity behaviour.

\item \textbf{Standing/exclusions:} candidate responsibility allocations established by
  U02/U03/U05;
  detailed management contracts and persistence structures remain unresolved.
  Execution of, behaviour upon, or handoff of managed information does not
  confer information-management responsibility for that information.

\end{itemize}

{\small Recorded unresolved matters: canonicalName; exact scope remains in the canonical record.}\par
{\footnotesize \href{../../../../continuum/application/application-function/ab0cdaa7ed4149ee856216efa5a265c4.md}{Canonical record: ab0cdaa7ed4149ee856216efa5a265c4}; full qualifications and provenance.}\par
\subsubsection*{IV.3.11.5.2 Resource Harmonisation Manager}
\phantomsection
\label{navigation:IV.3.11.5.2}
\addcontentsline{toc}{subsubsection}{IV.3.11.5.2 Resource Harmonisation Manager}

\label{element:a6580e0b9f0f40c9a33d731d20110057}
\index{Resource Harmonisation Manager (IV.3.11.5.2)}
{\small ApplicationFunction · Candidate — accepted Step-B baseline; final Solution Architecture unapproved · current}\par

\textbf{Internal Mneme responsibility/capability, not a fifth exposed service or an
independent execution framework.} \href{\detokenize{../../../../../docs/markdown/05-information-systems-architecture/sources.md#u10}}{U10} establishes the
\hyperref[element:a6580e0b9f0f40c9a33d731d20110057]{exceptional delivery path}.
It identifies/manages conditions that cannot proceed normally and causes appropriate
bounded Harmonisation Ergon work; it does not perform domain harmonisation itself.



\begin{itemize}


\item \textbf{Expected use:} approximately 5\% or less at worst of resource-delivery activity,
  a design expectation, not a hard SLA. Normal handling remains efficient.

\item \textbf{Execution:} new governed work through the established Harmonia work/execution
  contract reuses Ponos; no resumed originating Ergon, duplicate/embedded execution
  framework or originating-Ergon generic conflict resolution.

\item \textbf{Information/reuse:} relevant condition/state/Context is preserved for governed
  handling; a reusable Ergo Library remains a candidate reuse mechanism only.

\item \textbf{Unresolved:} categories, detection/concurrency/harmonisation algorithms,
  Harmonisation Ergon catalogue, detailed work contracts and repeated-harmonisation
  terminal/escalation handling (CDG-46). No UI/forms architecture is selected.
  Kleio cross-target audit harmonisation remains R1.x/R2.x SEP.

\end{itemize}

{\small Recorded unresolved matters: canonicalName; exact scope remains in the canonical record.}\par
{\footnotesize \href{../../../../continuum/application/application-function/a6580e0b9f0f40c9a33d731d20110057.md}{Canonical record: a6580e0b9f0f40c9a33d731d20110057}; full qualifications and provenance.}\par
\section*{IV.3.12 Mnemosyne}
\phantomsection
\label{navigation:IV.3.12}
\addcontentsline{toc}{section}{IV.3.12 Mnemosyne}

\subsection*{IV.3.12.1 Application Overview}
\phantomsection
\label{navigation:IV.3.12.1}
\addcontentsline{toc}{subsection}{IV.3.12.1 Application Overview}

\subsubsection*{IV.3.12.1.1 Mnemosyne}
\phantomsection
\label{navigation:IV.3.12.1.1}
\addcontentsline{toc}{subsubsection}{IV.3.12.1.1 Mnemosyne}

\label{element:7985e647bf9e4ca3ba70faf9e3e4ae02}
\index{Mnemosyne (IV.3.12.1.1)}
{\small ApplicationComponent · Candidate — accepted Step-B baseline; final Solution Architecture unapproved · current}\par

\begin{itemize}


\item \textbf{Role/responsibility:} \textbf{Harmonia's Persistence Manager} under
  \href{\detokenize{../../../../../docs/markdown/05-information-systems-architecture/sources.md#u11}}{U11}, establishing authoritative durable state/version and the
  information/metadata required to sustain activity Context. Persistence conveys
  no originating information/application ownership.

\item \textbf{First-class responsibilities:} transaction oversight; persistent-service
  oversight/reporting; shard management; authoritative persistence reporting;
  persistence logging; compartment validation/reapplication; hosting Themis-owned
  PEP enforcement; \textbf{Search Management}; \textbf{Research Metadata Management}.
  \hyperref[element:7985e647bf9e4ca3ba70faf9e3e4ae02]{The management view}
  establishes each boundary. Underlying services may retain their own transactions;
  Mnemosyne does not become infrastructure management.

\item \textbf{Information/outcomes:} durable representations/version/metadata, persistence
  observations and outcomes, search capability metadata and distinct Research Metadata.
  API/engine success is not automatically the authoritative Harmonia outcome;
  uncertain effects remain explicit. Persistence logging is distinct from Kleio audit.

\item \textbf{Search:} Mneme exposes Adhoc Search; Mnemosyne performs underlying authoritative
  managed-persistence search, without treating a transient cache as complete.
  Persistence-domain/shard/service/search-index capability metadata may be managed.

\item \textbf{Research Metadata:} manages capture, persistence, management and search utilisation
  of AI-derived meaning/relationships across content sets based on proximity. Source
  Content remains unchanged and derived metadata distinguishable. Deriving AI
  processing executes \textbf{outside Mnemosyne} through appropriate Harmonia behaviour/
  execution. Detailed provenance/confidence/schema remain unresolved.

\item \textbf{Persistence domains:} Clinical Information, C\&G and O\&A retain U06 guardrails:
  coherent C\&G authority; O\&A horizontal partitionability preserving identity,
  provenance, relationships and meaning. Domains select no database/service topology
  or shard key. Mnemosyne manages shards within semantic constraints, without
  independently deciding sharding safety.

\item \textbf{Governance:} compartments are validated/reapplied without loss or weakening;
  semantic applicability is not independently determined here. Themis owns PEP
  semantics and Criteria; Mnemosyne controls/prevents the persistence operation.
  Kleio mandatory AuditEvent persistence uses its
  \href{\detokenize{../../../../../docs/markdown/05-information-systems-architecture/component-register.md#kleio-capture-authority}}{governed intrinsic platform authority}, distinct
  from the audited Identity; Mnemosyne gains no equivalent intrinsic authority.

\item \textbf{Existing registry support:} server-side Provider Registry durable integrity,
  preservation/history and recovery responsibilities remain as described in C13
  and the data view; this correction does not redesign their semantic validation.

\item \textbf{Collaboration/exclusions:} Mneme manages Context/Pragma information consequences;
  Mnemosyne establishes/reports persistence. Energeia observes/reports persistence
  operational state, without acquiring management. No application-facing raw storage,
  AI derivation inside Mnemosyne, technology selection or workflow engine is established.

\item \textbf{Standing:} U11 and explicit task decisions settle responsibilities; detailed
  oversight/transaction/compartment/outcome/search/Research Metadata contracts remain
  CDG-47/48, alongside historical structures CDG-12/13/40.

\end{itemize}

{\small Recorded unresolved matters: canonicalName; exact scope remains in the canonical record.}\par
{\footnotesize \href{../../../../continuum/application/application-component/7985e647bf9e4ca3ba70faf9e3e4ae02.md}{Canonical record: 7985e647bf9e4ca3ba70faf9e3e4ae02}; full qualifications and provenance.}\par
\section*{IV.3.13 Ponos}
\phantomsection
\label{navigation:IV.3.13}
\addcontentsline{toc}{section}{IV.3.13 Ponos}

\subsection*{IV.3.13.1 Application Overview}
\phantomsection
\label{navigation:IV.3.13.1}
\addcontentsline{toc}{subsection}{IV.3.13.1 Application Overview}

\subsubsection*{IV.3.13.1.1 Ponos}
\phantomsection
\label{navigation:IV.3.13.1.1}
\addcontentsline{toc}{subsubsection}{IV.3.13.1.1 Ponos}

\label{element:b521762c1ea444b48cd5f66be3752d30}
\index{Ponos (IV.3.13.1.1)}
{\small ApplicationComponent · Candidate — accepted Step-B baseline; final Solution Architecture unapproved · current}\par

\begin{itemize}


\item \textbf{Role/responsibility:} reusable execution framework/component pattern, peer
  to Energeia C39, for governed execution and observable activity progression.
  \href{\detokenize{../../../../../docs/markdown/05-information-systems-architecture/sources.md#u02}}{U02} adjudicates this role; \href{\detokenize{../../../../../docs/markdown/05-information-systems-architecture/sources.md#a02}}{A02} and
  \href{\detokenize{../../../../../docs/markdown/05-information-systems-architecture/sources.md#l02}}{L02/L01/L23} retain responsibility context and historical detail.

\item \textbf{Features/contained constructs:} work ingress and dispatcher, workflow/sequence
  resolution where governed control is established, activity invocation, execution
  concurrency/timeouts, execution retry/re-attempts, checkpoint coordination and
  administrative inspection. Applies Praxis-defined output-processing/continuation
  semantics and routing; not every output-derived Pragma continues that Praxis.
  \href{\detokenize{../../../../../docs/markdown/05-information-systems-architecture/sources.md#u18}}{U18} establishes the distinction, leaving representation open.
  Subscription matching belongs to Mneme; handoff
  redelivery/recovery and dead-letter semantics belong to Petasos.
  Consumer loops/thread pools are documented mechanisms, not logical children.
  Ordinary semantic interest fan-out to different Ergon types is orthogonal to
  ActionableTask.ExecutionConcurrency permitting concurrent undertakings by
  instances of the same type, each through its own FulfillmentTask. DigitalTwin
  specialises execution with serial Behaviour processing within a Twin; common
  concurrency facilities do not permit intra-Twin Behaviour concurrency (U19).

\item \textbf{Information:} executes consumer undertakings using governed information/context
  resolved through Mneme; consults established control definitions and produces
  execution/progression observations contributing to Journey. Petasos handoff
  does not make an opaque payload into a Ponos-owned Pragma.

\item \textbf{Persistence:} coordinates recoverable handoff through Petasos and governed
  information change through Mneme/Mnemosyne. Does not establish durable state
  by completing an activity or writing a checkpoint to a cache.

\item \textbf{Collaboration/exclusions:} tightly coupled to peer Ergo/Ergon responsibilities,
  executing/managing bounded Ergon activity conforming to the Ergo contract;
  does not structurally contain Ergo or Ergon. Uses Praxis, Erga, Themis, Mneme
  Pragma Management and Petasos. DigitalTwin C19 and Dokimasia C20 specialise/
  reuse this pattern; they are not children inside a Ponos runtime. Conceptual
  supertype language does not mandate literal implementation inheritance.
  Does not own Pragma information management, Publishing/Subscription/Pattern
  semantic selection, workflow definitions, concrete domain logic, broker lifecycle,
  external standards/transport or universal clinical-work determination.

\item \textbf{Standing:} specialisation relationships adjudicated for the candidate by U02;
  detailed model remains candidate. CDG-03/08/16–18 retain incomplete control
  and continuation semantics. Numerical worker/retry defaults are not baselined.

\end{itemize}

{\small Recorded unresolved matters: canonicalName; exact scope remains in the canonical record.}\par
{\footnotesize \href{../../../../continuum/application/application-component/b521762c1ea444b48cd5f66be3752d30.md}{Canonical record: b521762c1ea444b48cd5f66be3752d30}; full qualifications and provenance.}\par
\subsection*{IV.3.13.2 Component Architecture}
\phantomsection
\label{navigation:IV.3.13.2}
\addcontentsline{toc}{subsection}{IV.3.13.2 Component Architecture}

\subsubsection*{IV.3.13.2.1 Ergon}
\phantomsection
\label{navigation:IV.3.13.2.1}
\addcontentsline{toc}{subsubsection}{IV.3.13.2.1 Ergon}

\label{element:49c9eb2bef0c47b8b5328621d5fc5479}
\index{Ergon (IV.3.13.2.1)}
{\small ApplicationComponent · Candidate — accepted Step-B baseline; final Solution Architecture unapproved · current}\par

\begin{itemize}


\item \textbf{Role/responsibility:} a single well-defined software component encapsulating
  a bounded fine-grained task-completion activity set within the Ergo contract,
  executed/managed through Ponos. \href{\detokenize{../../../../../docs/markdown/05-information-systems-architecture/sources.md#u02}}{U01/U02}, \href{\detokenize{../../../../../docs/markdown/05-information-systems-architecture/sources.md#l03}}{L03/L19}.

\item \textbf{Features:} specialised transformation/mapping, validation, identifier/reference
  checking, enrichment, routing, result processing, fan-out or registry-change
  behaviour. Reuses C05 common facilities; \texttt{\seqsplit{ErgonBase}} is historical contract
  support, not an architectural inheritance requirement.

\item \textbf{Information:} acts upon or generates Pragma, payload/resource information,
  outputs, checkpoints and observations through governed contracts. It primarily
  acts on loaded Context values, may change working values and may generate new
  Content as \textbf{Pragma Output}, distinct from Context mutation. Bounded \textbf{Adhoc
  Resource Access / Adhoc Search} remains permitted for emerging needs. \textbf{No discrete
  Information Architecture representation or Ergon Information Resource exists
  merely because the component is an Ergon.} Activity is not information about it.

\item \textbf{Output progression:} one or more execution outputs may yield multiple Pragma.
  The \hyperref[element:27ad143e7c784e378556af8a57c438f7]{Pathway-continuation role}
  identifies information continuing the current Praxis; independently resulting
  Pragma remain available for governed use, without automatically advancing it.

\item \textbf{Persistence/collaboration:} requests governed information use/change through
  Mneme; does not own Pragma management or durable state. Erga collects components;
  Ponos executes them; Petasos effects persistent handoff. Every Ergon exposes
  IngressEndpoint and EgressEndpoint C29; traversal does not expose internal logic.
  Praxis relationship detail remains open. The originating Ergon owns no persistence
  mechanics, shard/database transactions or generic conflict resolution; exceptional
  harmonisation is new governed work under U10.

\item \textbf{Specialisations:} Behaviour C25, Dokimasia Ergon C26 and, through the latter,
  AI Assurance Ergon C27. They retain normal Ergo/Ergon facilities, constrained
  by the specialised role and Dokimasia's terminal rules where applicable.

\item \textbf{Standing/exclusions:} bounded behaviour established for the candidate;
  competing atomic/pure/idempotent/branching claims remain CDG-05. An ordinary
  Ergon may intentionally invoke an AI service internally as an implementation
  choice without becoming AI Assurance Ergon or transferring authority.

\end{itemize}

{\small Recorded unresolved matters: canonicalName; exact scope remains in the canonical record.}\par
{\footnotesize \href{../../../../continuum/application/application-component/49c9eb2bef0c47b8b5328621d5fc5479.md}{Canonical record: 49c9eb2bef0c47b8b5328621d5fc5479}; full qualifications and provenance.}\par
\subsection*{IV.3.13.3 Interfaces}
\phantomsection
\label{navigation:IV.3.13.3}
\addcontentsline{toc}{subsection}{IV.3.13.3 Interfaces}

\subsubsection*{IV.3.13.3.1 PathwayEndpoint}
\phantomsection
\label{navigation:IV.3.13.3.1}
\addcontentsline{toc}{subsubsection}{IV.3.13.3.1 PathwayEndpoint}

\label{element:800789f54ba24f53a354600d5cd43f6f}
\index{PathwayEndpoint (IV.3.13.3.1)}
{\small ApplicationInterface · Candidate — accepted Step-B baseline; final Solution Architecture unapproved · current}\par

Application Interface conceptually equivalent to ArchiMate ApplicationInterface.
Every Ergon exposes IngressEndpoint and EgressEndpoint. These surfaces do not
expose internal processing; endpoint identity and technical interfaces remain open.


{\small Recorded unresolved matters: canonicalName; exact scope remains in the canonical record.}\par
{\footnotesize \href{../../../../continuum/application/application-interface/800789f54ba24f53a354600d5cd43f6f.md}{Canonical record: 800789f54ba24f53a354600d5cd43f6f}; full qualifications and provenance.}\par
\section*{IV.3.14 Pylai}
\phantomsection
\label{navigation:IV.3.14}
\addcontentsline{toc}{section}{IV.3.14 Pylai}

\subsection*{IV.3.14.1 Application Overview}
\phantomsection
\label{navigation:IV.3.14.1}
\addcontentsline{toc}{subsection}{IV.3.14.1 Application Overview}

\subsubsection*{IV.3.14.1.1 Pylai}
\phantomsection
\label{navigation:IV.3.14.1.1}
\addcontentsline{toc}{subsubsection}{IV.3.14.1.1 Pylai}

\label{element:e3e4147fea3640198dc251ed41ee854d}
\index{Pylai (IV.3.14.1.1)}
{\small ApplicationComponent · Candidate — accepted Step-B baseline; final Solution Architecture unapproved · current}\par

\begin{itemize}


\item \textbf{Role/responsibility:} \textbf{Harmonia's External Interoperability Service},
  enabling/managing system-to-system communication irrespective of style, protocol,
  standard or technology. Messaging, synchronous/asynchronous RMI, APIs, files and
  polling are supported architectural possibilities; RMI does not mean Java RMI.
  \href{\detokenize{../../../../../docs/markdown/05-information-systems-architecture/sources.md#u14}}{U14}, \href{\detokenize{../../../../../docs/markdown/05-information-systems-architecture/pylai-external-interoperability.md}}{canonical Pylai/Horos view}.

\item \textbf{Constructs/patterns:} Terma C43 acts at the logical external interface;
  Diadromí C44 carries/represents transient interaction. PylaiPathway connects
  governed Praxis/Ergo work; PylaiBroker serves executing BusinessLogic;
  PylaiServer changes use Pathway and READ/SEARCH uses Diadromí-Controller C45.
  These patterns do not mandate classes, services or deployment boundaries.

\item \textbf{Information/persistence:} Communication Information Unit preserves raw payload,
  communication metadata and applicable external identity; exact schema unresolved.
  Ingress Loader establishes Task.input/Context, with semantic mapping in paired
  BusinessLogic. Durable acceptance requires authoritative Pragma/Task persistence
  through Mneme/Mnemosyne. Protocol ACK and queue/cache success are distinct.
  Diadromí is transient and non-recoverable; Petasos owns durable Parcel handoff.

\item \textbf{Security/evidence:} carries external identity; Themis establishes Harmonia
  authority/enforcement. Kleio constructs authoritative AuditEvent from reported
  significant transactions; Dokimasia assures; Journey records applicable observed
  business progression or interoperability failure. None of these are synonyms.

\item \textbf{Exclusions/standing:} no execution framework, persistence manager, authority
  inference, fabricated READ/SEARCH work or mandatory Controller for all interactions.
  AX-02/13 non-destructive fail-closed publication remains. Technology/handover
  details remain CDG-32/34/50; CDG-51 records the resolved Strategy boundary.
  Kleio mandatory capture uses its own governed platform authority, not the
  external/acting Identity's grants; Pylai gains no such authority.

\end{itemize}

{\small Recorded unresolved matters: canonicalName; exact scope remains in the canonical record.}\par
{\footnotesize \href{../../../../continuum/application/application-component/e3e4147fea3640198dc251ed41ee854d.md}{Canonical record: e3e4147fea3640198dc251ed41ee854d}; full qualifications and provenance.}\par
\subsection*{IV.3.14.2 Information Architecture}
\phantomsection
\label{navigation:IV.3.14.2}
\addcontentsline{toc}{subsection}{IV.3.14.2 Information Architecture}

\subsubsection*{IV.3.14.2.1 Pylai.Diadromí Topic}
\phantomsection
\label{navigation:IV.3.14.2.1}
\addcontentsline{toc}{subsubsection}{IV.3.14.2.1 Pylai.Diadromí Topic}

\label{element:978f5c2bd2d343ffaf3a13dff700c258}
\index{Pylai.Diadromí Topic (IV.3.14.2.1)}
{\small DataObject · Bounded adjudicated application information meaning; Domain05 remains the accepted candidate Step-B baseline, not final Solution Architecture. Detailed representation is deferred; only this defined logical object is projected. · current}\par

\textbf{Pylai.Terma} is the logical external interoperability interface/boundary.
\texttt{\seqsplit{Pylai.IngressTerma}} and \texttt{\seqsplit{Pylai.EgressTerma}} specialise it where useful. Terma owns
external communication mechanics: connection/session handling, listening, polling,
transport/protocol interaction, framing, segmentation/reassembly, protocol
acknowledgements, API/RMI/file interaction, interaction-local communication retry
and external identity carriage. One Terma may use one or more technical endpoints/connectors;
it is not reducible to a socket, URL, Camel endpoint or deployed adapter.



\begin{quote}


\textbf{Pylai.Terma manages the communication stack; the applicable Harmonia
execution framework manages the application meaning.}



\end{quote}


Terma owns the mechanics necessary to establish, conduct and complete an external
communication interaction. It does not own the Harmonia application semantics
of the information communicated. Application interpretation and processing use
the applicable Harmonia execution framework, principally Praxis/Ergo/Ergon,
except where an explicitly defined interaction pattern provides otherwise.
The \href{\detokenize{../../../../../docs/markdown/05-information-systems-architecture/pylai-external-interoperability.md#pylai-server-controller}}{Server READ/SEARCH pattern} is such a defined
pattern; it grants Terma no application/business interpretation responsibility.



\begin{quote}\small\ttfamily
External System → Pylai.Terma → Pylai.Diadromí → Praxis / Ergo / Ergon\\
                  communication mechanics       application interpretation,\\
                                                transformation / processing\\

\end{quote}


Ingress/Egress describe \textbf{information direction}, irrespective of client/server
role, who opened the connection or push/pull behaviour. An IngressTerma may poll;
an EgressTerma may expose an external file interaction.



\begin{quote}


\textbf{Pylai.Diadromí is the transient carriage and interaction context through which
a Pylai.Terma exchanges communication with its Harmonia-side participant for the
duration of an external interoperability interaction.}



\textbf{Terma acts; Diadromí carries/represents transient interaction.}



\end{quote}


A Diadromí may span multiple protocol exchanges, acknowledgements, retries within
the active interaction and segmented communication. An intermediate protocol
acknowledgement need not establish logical handover or authoritative Harmonia
acceptance. Diadromí is neither an Ergon, Task/Pragma alias nor a durable queue.
Its exact representation is unresolved.



\textbf{Pylai.Diadromí realises
\hyperref[element:c34b4929f71946708d37dbe6c1da920a]{Metabolē}}
for transient information progression through Pylai interoperability activity.
This is genuine architectural Realisation, not analogy, synonymy with Metabolē
or inheritance from Pragma. Pylai retains responsibility; Diadromí's transient,
non-recoverable lifecycle does not acquire Pragma persistence or Mneme management
requirements. Conceptual ActivityContext does not replace a specific interaction
action or grant authorisation.



Source standing: Bounded adjudicated application information meaning; Domain05 remains the accepted candidate Step-B baseline, not final Solution Architecture. Detailed representation is deferred; only this defined logical object is projected.


{\footnotesize \href{../../../../continuum/application/data-object/978f5c2bd2d343ffaf3a13dff700c258.md}{Canonical record: 978f5c2bd2d343ffaf3a13dff700c258}; full qualifications and provenance.}\par
\subsection*{IV.3.14.3 Services}
\phantomsection
\label{navigation:IV.3.14.3}
\addcontentsline{toc}{subsection}{IV.3.14.3 Services}

\subsubsection*{IV.3.14.3.1 External Interoperability Service}
\phantomsection
\label{navigation:IV.3.14.3.1}
\addcontentsline{toc}{subsubsection}{IV.3.14.3.1 External Interoperability Service}

\label{element:43b79574792443faa57daa76b9a47a19}
\index{External Interoperability Service (IV.3.14.3.1)}
{\small ApplicationService · Candidate — adjudicated Domain05 Step-B baseline; final Solution Architecture unapproved · current}\par

\begin{quote}


\textbf{Pylai is Harmonia's External Interoperability Service, responsible for enabling
and managing system-to-system communication between Harmonia and external systems
and clients, irrespective of communication style, protocol, standard or technology.}



\end{quote}


Pylai may support asynchronous messaging, synchronous or asynchronous remote
method invocation (RMI), APIs, files, polling and other external interoperability
mechanisms. RMI is an architectural interaction concept; it does not mean Java RMI.
HL7/FHIR gateway slices are examples, not the complete Pylai architecture.
Apache Camel or another integration framework may realise external connectors/
technical endpoints; it is a candidate implementation mechanism, not Pylai itself.



External representations conform to their applicable interoperability contracts.
Publication SHALL be non-destructive and fail-closed: only permitted information/
extensions may leave. Harmonia-private authority, security-context, cache,
persistence, governance and concurrency metadata are not automatically exportable.
Egress terminates management of the emitted representation (AX-02/13); retained
provenance does not extend control over the external copy.



The \href{\detokenize{../../../../../docs/markdown/05-information-systems-architecture/gaps.md#pylai-strategy-conflict}}{Strategy/Terma reconciliation} distinguishes
external communication mechanics from Harmonia application meaning and generic
network/infrastructure management.


{\small Recorded unresolved matters: canonicalName; exact scope remains in the canonical record.}\par
{\footnotesize \href{../../../../continuum/application/application-service/43b79574792443faa57daa76b9a47a19.md}{Canonical record: 43b79574792443faa57daa76b9a47a19}; full qualifications and provenance.}\par
\section*{IV.3.15 Themis}
\phantomsection
\label{navigation:IV.3.15}
\addcontentsline{toc}{section}{IV.3.15 Themis}

\subsection*{IV.3.15.1 Application Overview}
\phantomsection
\label{navigation:IV.3.15.1}
\addcontentsline{toc}{subsection}{IV.3.15.1 Application Overview}

\subsubsection*{IV.3.15.1.1 Themis}
\phantomsection
\label{navigation:IV.3.15.1.1}
\addcontentsline{toc}{subsubsection}{IV.3.15.1.1 Themis}

\label{element:899d34ec491948b68aad09e507dea6d1}
\index{Themis (IV.3.15.1.1)}
{\small ApplicationComponent · Candidate — accepted Step-B baseline; final Solution Architecture unapproved · current}\par

\begin{itemize}


\item \textbf{Role/responsibility:} \textbf{Themis owns and manages Harmonia's Authorisation
  Model}, including PAP, PRP, PDP, Criteria/collections, AuthorisationRoles,
  Identity assignments, role/group resolution, Criteria compilation and the
  distributed PEP framework. Themis consumes trusted authenticated Identity;
  it does not authenticate that Identity. \href{\detokenize{../../../../../docs/markdown/05-information-systems-architecture/sources.md#u12}}{U12} refines
  \href{\detokenize{../../../../../docs/markdown/05-information-systems-architecture/sources.md#u08}}{U08} within \hyperref[element:5de2ae2c3b794c2d9f3a14efd51d564f]{AX-07}.

\item \textbf{PAP — Policy Administration Point:} administration and lifecycle management
  of the Themis-owned model. Management/review/approval activity uses
  \href{\detokenize{../../../../../docs/markdown/05-information-systems-architecture/component-register.md#authorisation-administration}}{Ergo/Praxis and Iris}; detailed Praxes/Erga,
  UI and versioning mechanisms remain deferred. No general-purpose policy
  language is implied by policy administration.

\item \textbf{PRP — Policy Repository Point:} authoritative repository semantics for the
  Authorisation Model. Durable persistence may be delegated through \textbf{Themis PRP
  → Mneme → Mnemosyne Harmonia Configuration \& Governance}. This requires no
  dedicated Themis physical database and transfers no model ownership to persistence.

\item \textbf{PDP — Policy Decision Point:} resolves authoritative policy, authenticated
  Identity and applicable context into precomputed Identity-specific
  \textbf{PEP-Criteria}, supplied through \textbf{PEP-Context} for local enforcement.
  An immediate decision remains possible; it uses the same governed semantics.

\item \textbf{Governed actions:} \textbf{Can I ACCESS XYZ? Can I CHANGE ABC? Can I FORWARD DEF?}
  ACCESS, CHANGE and FORWARD are distinct; they SHALL NOT collapse into CRUD.
  FORWARD may depend on destination/context independently of ACCESS.
  Unauthenticated or unauthorised operations default to \textbf{DENY}; unavailable
  authoritative policy determination remains distinct from that decision.

\end{itemize}






\par\medskip\textbf{Identity.Authentication — separately provided trusted Identity}\par


\textbf{Identity.Authentication establishes WHO an Individual or System is. It is a
Separately Provided Capability (SEP).} The primary pattern is standards-based
OIDC/OAuth integration with an externally managed Identity Provider (IdP).
Harmonia does not own or provide an identity-management, credential-management,
password-management or MFA framework.



\begin{quote}


\textbf{Keycloak + PostgreSQL SHALL provide the default FOSS Identity Provider for
development, test and POC environments.}



\end{quote}


Keycloak remains an external, replaceable implementation component. Harmonia's
Identity.Authentication architecture SHALL NOT acquire Keycloak-specific
dependencies. This target does not prescribe a production IdP or an internal
Harmonia credential store. Authentication establishes trusted Identity only;
it SHALL NOT determine Harmonia Function permissions, Function.Scope or PEP-Criteria.



\textbf{Identity.Authorisation} establishes what that authenticated Identity may do
within Harmonia and is Themis-owned. The same authorisation model applies to
human and system Identities; differing allowable assignments do not establish
separate architectures.







\par\medskip\textbf{Identity and role terminology}\par


\begin{longtable}{>{\raggedright\arraybackslash}p{\dimexpr\linewidth/2-2\tabcolsep\relax}>{\raggedright\arraybackslash}p{\dimexpr\linewidth/2-2\tabcolsep\relax}}
Term & Candidate meaning and boundary \\
\textbf{Identity} & The authenticated Individual or System consumed by Themis. This use does not redefine Information Unit Identity or the Domain04 professional identifier assertions. \\
\textbf{Practitioner} & The individual practitioner, retaining the \href{\detokenize{../../../../../docs/markdown/04-information-architecture/information-families/practitioner.md#31-practitioner}}{Domain04 workforce-entity meaning}; registration status is a separate assertion. \\
\textbf{FHIR PractitionerRole} & The formal practitioner role/service relationship represented through FHIR. \\
\textbf{PractitionerCapacity} & The selected operational capacity in which a Practitioner is currently acting within Harmonia; it does not itself confer authority. \\
\textbf{AuthorisationRole} & A named, governed, assignable collection of additive Themis Criteria and/or Criteria collections. It is distinct from PractitionerCapacity, FHIR PractitionerRole and a Domain03 Business Role. \\
\textbf{Criterion / Criteria} & One governed permitted combination / multiple such combinations in the Themis Authorisation Model. \\
\textbf{PEP-Criteria} & The precomputed, Identity-specific set of permitted combinations derived by Themis for constrained local enforcement. \\
\textbf{PEP-Context} & The Themis-derived operational enforcement context for an authenticated Identity, through which applicable PEP-Criteria is supplied. It is distinct from the general Mneme working Context. \\
\end{longtable}


The \href{\detokenize{../../../../../docs/markdown/04-information-architecture/information-families/practitioner.md#36-practitioner-role-binding}}{Domain04 Practitioner Role Binding}
and \href{\detokenize{../../../../../docs/markdown/03-business-architecture/actors-roles/roles.md}}{Domain03 Business Roles}
retain their meanings. Neither lexical similarity nor FHIR convention establishes
their mapping to PractitionerCapacity or AuthorisationRole.







\par\medskip\textbf{Criterion conjunction and additive grants}\par


A Criterion represents a \textbf{governed permitted combination} across applicable
dimensions, using \href{\detokenize{../../../../../docs/markdown/05-information-systems-architecture/component-register.md#authorisation-nomenclature}}{Calliope-controlled nomenclature}:



\begin{quote}\small\ttfamily
Criterion =\\
    Information.Compartment\\
    AND Information.Classification\\
    AND Function.Name\\
    AND Function.Scope\\
    AND applicable Temporal constraints\\

\end{quote}


The four principal dimensions describe \textbf{what} is authorised; applicable
temporal constraints describe \textbf{when}. Within an individual Criterion,
applicable dimensions are \textbf{conjunctive}. Between Criteria, grants are
\textbf{additive / UNION}: each Criterion independently contributes its permitted
combination. Only explicitly governed combinations have authorisation meaning.
Absence of a matching effective Criterion results in \textbf{DENY}.



\textbf{Explicit negative/DENY Criteria are not part of this Candidate Architecture.}
Neither arbitrary Boolean policy expressions, deny-overrides semantics nor a
general-purpose policy language is introduced by this model.







\par\medskip\textbf{Criteria containment, AuthorisationRoles and Identity assignments}\par


Criteria support containment, collection and nesting for management and reuse:



\begin{quote}\small\ttfamily
CriteriaCollection                 (conceptual label; final name unresolved)\\
  ├── Criterion\\
  ├── Criterion\\
  └── CriteriaCollection\\
       ├── Criterion\\
       └── Criterion\\

\end{quote}


Flattening nested collections produces the same additive set of effective
Criteria. Nesting creates no arbitrary Boolean semantics and requires no Role
hierarchy. The exact implementation name/class hierarchy and representation
remain deferred; this diagram establishes no lifecycle/ownership semantics
beyond the stated management/reuse requirement.



An \textbf{AuthorisationRole} is named, governed and assignable and contains Criteria
and/or Criteria collections. Its contents are additive:



\begin{quote}\small\ttfamily
EffectiveCriteria(AuthorisationRole)\\
    = UNION(all contained Criteria)\\
\\
IdentityEffectiveCriteria =\\
    DirectCriteria\\
    UNION AuthorisationRole1.Criteria\\
    UNION AuthorisationRole2.Criteria\\
    UNION ...\\

\end{quote}


An Identity may have directly assigned Criteria, one or more AuthorisationRoles,
or both. This applies to human and system Identities. The UNION is of applicable
grants, subject to governed assignments, temporal validity, revocation and
applicable operational context; it does not remove a Criterion's conjunctions
or require every contained Criterion to match simultaneously.







\par\medskip\textbf{Predetermined temporal validity and explicit revocation}\par


\textbf{Temporal validity} establishes predetermined applicability/validity bounds.
Temporal constraints may apply to a Criterion, a direct Criterion assignment,
an AuthorisationRole assignment and applicable operational context such as
PractitionerCapacity. The exact temporal expression model remains deferred.



\textbf{Revocation} withdraws an otherwise valid authorisation. It SHALL NOT be
modelled merely as temporal expiry. Both concepts are required. Propagation of
revocation and invalidation of already-issued PEP-Criteria/PEP-Context remain
unresolved; no mechanism or continued validity after revocation is inferred.







\par\medskip\textbf{PractitionerCapacity and operational responsibility}\par


For a Practitioner, the available PractitionerCapacity set is constrained by
the practitioner's governed profile and applicable operational context. The
Practitioner selects \textbf{one active PractitionerCapacity} for the applicable
session/operational or roster context. It supplies operational context used by
Themis when deriving effective authorisation; \textbf{it does not itself confer authority}.
Selection UI and the exact PractitionerCapacity ↔ FHIR PractitionerRole mapping
remain deferred.



\begin{quote}


\textbf{Authentication establishes who the Individual is. PractitionerCapacity
establishes the capacity in which a Practitioner is currently acting. Themis
establishes what that Identity, acting in that capacity, may do.}



\end{quote}


\begin{quote}\small\ttfamily
graph TD\\
    Auth["Authenticated Practitioner"] --> Eligible["Eligible PractitionerCapacity set"]\\
    Profile["Governed practitioner profile"] --> Eligible\\
    Context["Applicable operational / roster context"] --> Eligible\\
    Eligible -->|"Practitioner selects ONE"| Capacity["Active PractitionerCapacity: context, no authority by itself"]\\
    Capacity -->|"authorisation context"| PDP["Themis: resolve governed Criteria / assignments / policy state"]\\
    PDP --> Criteria["Identity-specific precomputed PEP-Criteria through PEP-Context"]\\
    Criteria --> PEP["Constrained local Themis PEP"]\\
    Capacity -.->|"operational responsibility context"| Responsibility["Agora: membership intent; detailed context use deferred"]\\

\end{quote}


Themis answers \textbf{“May this Identity perform this action on this information?”}
Future Agora collaboration detail concerns \textbf{“Is this Practitioner currently
responsible for receiving/responding to this operational activity?”} These are
separate responsibilities even where they consume the same PractitionerCapacity.
Receipt of an escalation or approval notification grants no information access.
Following it into Iris or another application surface requires normal Themis PEP
enforcement using current valid PEP-Criteria. U16 establishes
\hyperref[element:d1c4053e68384080b8eb79d2008632ba]{Agora control and Omilia operations};
detailed membership, escalation and approval-routing mechanics remain deferred.







\par\medskip\textbf{PEP-Context and distributed enforcement}\par


\begin{quote}


\textbf{PEP-Context is the Themis-derived policy enforcement context established for
an authenticated Identity. It supplies the precomputed PEP-Criteria applicable
to that Identity that a Themis Policy Enforcement Point requires for local
enforcement decisions.}



\end{quote}


\textbf{PEP-Criteria is a precomputed, Identity-specific set of permitted combinations},
not a general-purpose executable policy language distributed to PEPs. Themis
resolves direct Criteria, AuthorisationRole-derived Criteria, nested collections,
applicable temporal constraints, PractitionerCapacity context where applicable
and authoritative policy state. PEPs SHALL NOT receive the complete model merely
because they need enforcement information.



\begin{quote}\small\ttfamily
Policy                Authorisation Model managed by Themis\\
Criteria              governed permitted combinations in the Authorisation Model\\
PEP-Criteria          precomputed Identity-specific permitted combinations\\
PEP-Context           operational context supplying applicable PEP-Criteria\\
Enforcement Decision  local PEP determination for a particular proposed action\\

\end{quote}


Enforcement occurs at the component capable of preventing the action. A PEP may
be a Themis-provided library/framework element embedded within a host module.
\textbf{The distributed PEP capability, implementation and enforcement semantics are
Themis-owned/derived/built.} The host controls the actual operation; Themis
supplies enforcement mechanism and semantics. Hosts SHALL NOT independently
reinterpret the Authorisation Model in place of the Themis PEP. No separate
central Policy Enforcement Software Component is created.



A local PEP SHALL perform deliberately \textbf{constrained, deterministic enforcement}
against authenticated Identity context, effective Information.Compartment and
Information.Classification, Function.Name, Function.Scope, applicable temporal
state and the Identity's valid PEP-Criteria. It SHALL NOT become a general-purpose
policy interpreter.



Themis SHALL retain sufficient provenance to explain effective PEP-Criteria
derivation: authoritative Criteria, collections, direct/AuthorisationRole
assignments and applicable policy state. Runtime PEP-Criteria need not carry
the full derivation history; stable references may support it. No final wire/data
structure, cache, refresh or distribution mechanism is prescribed.



PEP-Context and existing security-context contracts are operational context;
they do not automatically become persisted resource content. Their detailed
relationship, schema and distribution mechanisms remain unresolved;
PEP-Context is not silently equated with a historical class or transport carrier.







\par\medskip\textbf{Untagged information and effective security attributes}\par


\begin{quote}


\textbf{Absence of explicit information security metadata does not imply absence
of security control. Untagged information is not unrestricted information.}



\end{quote}


Where required Information.Compartment and/or Information.Classification
attributes are absent, effective values SHALL be established through governed,
deterministic rules. If required effective attributes cannot be established,
enforcement SHALL fail closed.



Deriving effective attributes for an authorisation decision is distinct from
persistently assigning/changing security metadata on an Information Unit.
Themis enforces the resulting attributes; it is not assumed to own all semantic
derivation rules. Exact derivation ownership and precedence, including possible
context/containment relationships, remain \href{\detokenize{../../../../../docs/markdown/05-information-systems-architecture/gaps.md#identity-deferred-design}}{deferred}.
Mnemosyne's established compartment validation/reapplication responsibility does
not settle those semantic rules.







\par\medskip\textbf{Availability and delegation}\par


An in-flight activity with a \textbf{valid PEP-Context sufficient for its required
enforcement decisions} may continue during temporary PDP unavailability.
PDP availability is required for a \textbf{new, renewed or replacement PEP-Context}.
If the PDP cannot establish that required context, the activity SHALL NOT proceed
beyond the point requiring it and must wait. \textbf{PDP unavailability is not policy
denial}: policy permits, policy denies, and authoritative policy determination
cannot presently be established remain distinct. This grants no unauthenticated
or unauthorised action permission and weakens no required audit-capture boundary.
Detailed validity, timeout/retry, revocation and renewal mechanisms are unresolved.



At this stage delegation is \textbf{Identity / PEP-Criteria pass-off (forwarding)},
subject to applicable \textbf{FORWARD} authority. Provenance SHALL preserve both the
originating Identity and receiving/acting Identity; the receiver does not become
the originator. No separate general-purpose delegation subsystem is introduced.
Detailed forwarding lifecycle/representation, PDP interaction and governance
mechanisms remain deferred; forwarding establishes no unrestricted impersonation
or authority-broadening rule.







\par\medskip\textbf{Governed authorisation administration and target allocation}\par


Authorisation administration is governed Harmonia work. Criteria creation/change,
collection management, AuthorisationRole management, direct Identity Criteria
assignment, AuthorisationRole assignment, temporal changes and revocation SHALL
be subject to appropriate governed management/review/approval processes through
\textbf{Ergo/Praxis}, with \textbf{Iris} providing human interaction. Detailed PAP management
Praxes/Erga and approval rules remain deferred; this does not make every change
require an identical approval sequence.



\begin{longtable}{>{\raggedright\arraybackslash}p{\dimexpr\linewidth/2-2\tabcolsep\relax}>{\raggedright\arraybackslash}p{\dimexpr\linewidth/2-2\tabcolsep\relax}}
Responsibility & Candidate target \\
Identity authentication & Keycloak / external OIDC IdP; SEP, Keycloak + PostgreSQL default for development/test/POC \\
Controlled authorisation nomenclature & Calliope \\
Authorisation model, Criteria/collections, AuthorisationRoles and Identity assignments & Themis-owned \\
Role/group resolution & Themis-owned; no separate human/system authorisation architecture \\
Criteria compilation & Themis-owned \\
Policy repository & Themis PRP, persisted through Harmonia via Mneme / Mnemosyne C\&G \\
Policy administration & Themis PAP \\
Governed management / review / approval workflows & Ergo/Praxis \\
Administration UI & Iris; no authorisation-model ownership or detailed reusable-form design \\
Runtime enforcement & Themis PEP library at action-controlling hosts \\
Untagged information handling & Governed effective-attribute resolution; exact ownership/precedence deferred \\
\end{longtable}


Harmonia owns its authorisation semantics. \textbf{jCasbin, OPA, Cedar or another
general-purpose policy framework is not an architectural dependency.} FOSS
implementations may later supply implementation references where appropriate;
their conceptual models SHALL NOT redefine Harmonia semantics.



\begin{itemize}


\item \textbf{Collaboration/exclusions:} governed access, activity, publication and
  collaboration hosts use Themis-owned PEP capability. Petasos hosts governed
  FORWARD enforcement because it controls the queue operation. Praxis may supply
  relevant context but is not the Authorisation Model; authorisation does not
  establish Praxis conformance. Themis is neither Audit nor Assurance. Significant
  policy/enforcement occurrences are reported through Kleio Audit Capture under
  applicable policy; clients need not construct FHIR AuditEvent.
  \textbf{Calliope defines security/assurance nomenclature; Themis decides/enforces;
  Dokimasia assures; Kleio audits.} Governance, Authorisation, Assurance and Audit
  remain distinct responsibilities; Praxis conformance and assurance findings
  are not authorisation grants. Dokimasia terminal constraints remain unchanged.

\item \textbf{Standing:} U12 consolidates the adjudicated Identity Candidate Architecture
  and refines U08's Criteria, policy-language deferral and delegation wording.
  It does not accept all Domain05 candidate material or begin Step B.
  \href{\detokenize{../../../../../docs/markdown/05-information-systems-architecture/gaps.md#themis-detailed-contracts}}{CDG-42} retains detailed design.
  U15 resolves the former U08 §15 conflict through Kleio-specific
  \href{\detokenize{../../../../../docs/markdown/05-information-systems-architecture/component-register.md#kleio-capture-authority}}{mandatory capture authority}, governed under AX-07.
  Historical fixed gate counts and automatic context persistence remain
  CDG-15/30 provenance, without reopening established security boundaries.

\end{itemize}

{\small Recorded unresolved matters: canonicalName; exact scope remains in the canonical record.}\par
{\footnotesize \href{../../../../continuum/application/application-component/899d34ec491948b68aad09e507dea6d1.md}{Canonical record: 899d34ec491948b68aad09e507dea6d1}; full qualifications and provenance.}\par
\subsection*{IV.3.15.2 Information Architecture}
\phantomsection
\label{navigation:IV.3.15.2}
\addcontentsline{toc}{subsection}{IV.3.15.2 Information Architecture}

\subsubsection*{IV.3.15.2.1 PEP-Context}
\phantomsection
\label{navigation:IV.3.15.2.1}
\addcontentsline{toc}{subsubsection}{IV.3.15.2.1 PEP-Context}

\label{element:90c3f7e03448452baec5c425c34fd77a}
\index{PEP-Context (IV.3.15.2.1)}
{\small DataObject · Candidate — adjudicated Domain05 Step-B baseline; final Solution Architecture unapproved · current}\par

\begin{quote}


\textbf{PEP-Context is the Themis-derived policy enforcement context established for
an authenticated Identity. It supplies the precomputed PEP-Criteria applicable
to that Identity that a Themis Policy Enforcement Point requires for local
enforcement decisions.}



\end{quote}


\textbf{PEP-Criteria is a precomputed, Identity-specific set of permitted combinations},
not a general-purpose executable policy language distributed to PEPs. Themis
resolves direct Criteria, AuthorisationRole-derived Criteria, nested collections,
applicable temporal constraints, PractitionerCapacity context where applicable
and authoritative policy state. PEPs SHALL NOT receive the complete model merely
because they need enforcement information.



\begin{quote}\small\ttfamily
Policy                Authorisation Model managed by Themis\\
Criteria              governed permitted combinations in the Authorisation Model\\
PEP-Criteria          precomputed Identity-specific permitted combinations\\
PEP-Context           operational context supplying applicable PEP-Criteria\\
Enforcement Decision  local PEP determination for a particular proposed action\\

\end{quote}


Enforcement occurs at the component capable of preventing the action. A PEP may
be a Themis-provided library/framework element embedded within a host module.
\textbf{The distributed PEP capability, implementation and enforcement semantics are
Themis-owned/derived/built.} The host controls the actual operation; Themis
supplies enforcement mechanism and semantics. Hosts SHALL NOT independently
reinterpret the Authorisation Model in place of the Themis PEP. No separate
central Policy Enforcement Software Component is created.



A local PEP SHALL perform deliberately \textbf{constrained, deterministic enforcement}
against authenticated Identity context, effective Information.Compartment and
Information.Classification, Function.Name, Function.Scope, applicable temporal
state and the Identity's valid PEP-Criteria. It SHALL NOT become a general-purpose
policy interpreter.



Themis SHALL retain sufficient provenance to explain effective PEP-Criteria
derivation: authoritative Criteria, collections, direct/AuthorisationRole
assignments and applicable policy state. Runtime PEP-Criteria need not carry
the full derivation history; stable references may support it. No final wire/data
structure, cache, refresh or distribution mechanism is prescribed.



PEP-Context and existing security-context contracts are operational context;
they do not automatically become persisted resource content. Their detailed
relationship, schema and distribution mechanisms remain unresolved;
PEP-Context is not silently equated with a historical class or transport carrier.


{\small Recorded unresolved matters: canonicalName; exact scope remains in the canonical record.}\par
{\footnotesize \href{../../../../continuum/application/data-object/90c3f7e03448452baec5c425c34fd77a.md}{Canonical record: 90c3f7e03448452baec5c425c34fd77a}; full qualifications and provenance.}\par
\subsubsection*{IV.3.15.2.2 PEP-Criteria}
\phantomsection
\label{navigation:IV.3.15.2.2}
\addcontentsline{toc}{subsubsection}{IV.3.15.2.2 PEP-Criteria}

\label{element:7e10287b7e634bce8aa2481eb481bb76}
\index{PEP-Criteria (IV.3.15.2.2)}
{\small DataObject · Candidate — adjudicated Domain05 Step-B baseline; final Solution Architecture unapproved · current}\par

\textbf{PEP-Criteria is a precomputed, Identity-specific set of permitted combinations},
not a general-purpose executable policy language distributed to PEPs. Themis
resolves direct Criteria, AuthorisationRole-derived Criteria, nested collections,
applicable temporal constraints, PractitionerCapacity context where applicable
and authoritative policy state. PEPs SHALL NOT receive the complete model merely
because they need enforcement information.



Themis derives this information; PEP-Context supplies it to constrained local enforcement. It is not a general-purpose policy language or a persisted resource by default.


{\small Recorded unresolved matters: canonicalName; exact scope remains in the canonical record.}\par
{\footnotesize \href{../../../../continuum/application/data-object/7e10287b7e634bce8aa2481eb481bb76.md}{Canonical record: 7e10287b7e634bce8aa2481eb481bb76}; full qualifications and provenance.}\par
\subsubsection*{IV.3.15.2.3 Criterion}
\phantomsection
\label{navigation:IV.3.15.2.3}
\addcontentsline{toc}{subsubsection}{IV.3.15.2.3 Criterion}

\label{element:b5be4886ec77466e860209443a03771b}
\index{Criterion (IV.3.15.2.3)}
{\small DataObject · Candidate — adjudicated Domain05 Step-B baseline; final Solution Architecture unapproved · current}\par

A Criterion represents a \textbf{governed permitted combination} across applicable
dimensions, using \href{\detokenize{../../../../../docs/markdown/05-information-systems-architecture/component-register.md#authorisation-nomenclature}}{Calliope-controlled nomenclature}:



\begin{quote}\small\ttfamily
Criterion =\\
    Information.Compartment\\
    AND Information.Classification\\
    AND Function.Name\\
    AND Function.Scope\\
    AND applicable Temporal constraints\\

\end{quote}


The four principal dimensions describe \textbf{what} is authorised; applicable
temporal constraints describe \textbf{when}. Within an individual Criterion,
applicable dimensions are \textbf{conjunctive}. Between Criteria, grants are
\textbf{additive / UNION}: each Criterion independently contributes its permitted
combination. Only explicitly governed combinations have authorisation meaning.
Absence of a matching effective Criterion results in \textbf{DENY}.



\textbf{Explicit negative/DENY Criteria are not part of this Candidate Architecture.}
Neither arbitrary Boolean policy expressions, deny-overrides semantics nor a
general-purpose policy language is introduced by this model.


{\small Recorded unresolved matters: canonicalName; exact scope remains in the canonical record.}\par
{\footnotesize \href{../../../../continuum/application/data-object/b5be4886ec77466e860209443a03771b.md}{Canonical record: b5be4886ec77466e860209443a03771b}; full qualifications and provenance.}\par
\subsubsection*{IV.3.15.2.4 AuthorisationRole}
\phantomsection
\label{navigation:IV.3.15.2.4}
\addcontentsline{toc}{subsubsection}{IV.3.15.2.4 AuthorisationRole}

\label{element:bad6431f87bd4bf59645c0962cc219a7}
\index{AuthorisationRole (IV.3.15.2.4)}
{\small DataObject · Candidate — adjudicated Domain05 Step-B baseline; final Solution Architecture unapproved · current}\par

Criteria support containment, collection and nesting for management and reuse:



\begin{quote}\small\ttfamily
CriteriaCollection                 (conceptual label; final name unresolved)\\
  ├── Criterion\\
  ├── Criterion\\
  └── CriteriaCollection\\
       ├── Criterion\\
       └── Criterion\\

\end{quote}


Flattening nested collections produces the same additive set of effective
Criteria. Nesting creates no arbitrary Boolean semantics and requires no Role
hierarchy. The exact implementation name/class hierarchy and representation
remain deferred; this diagram establishes no lifecycle/ownership semantics
beyond the stated management/reuse requirement.



An \textbf{AuthorisationRole} is named, governed and assignable and contains Criteria
and/or Criteria collections. Its contents are additive:



\begin{quote}\small\ttfamily
EffectiveCriteria(AuthorisationRole)\\
    = UNION(all contained Criteria)\\
\\
IdentityEffectiveCriteria =\\
    DirectCriteria\\
    UNION AuthorisationRole1.Criteria\\
    UNION AuthorisationRole2.Criteria\\
    UNION ...\\

\end{quote}


An Identity may have directly assigned Criteria, one or more AuthorisationRoles,
or both. This applies to human and system Identities. The UNION is of applicable
grants, subject to governed assignments, temporal validity, revocation and
applicable operational context; it does not remove a Criterion's conjunctions
or require every contained Criterion to match simultaneously.


{\small Recorded unresolved matters: canonicalName; exact scope remains in the canonical record.}\par
{\footnotesize \href{../../../../continuum/application/data-object/bad6431f87bd4bf59645c0962cc219a7.md}{Canonical record: bad6431f87bd4bf59645c0962cc219a7}; full qualifications and provenance.}\par
